\documentclass{article}
\usepackage{graphicx}
\usepackage{capt-of}
\usepackage{amsmath}
\usepackage{amssymb}
\usepackage[colorlinks=true, linkcolor=blue, citecolor=blue, urlcolor=blue]{hyperref}
\usepackage[a4paper, left=2.5cm, right=3cm, top=2.5cm, bottom=2.5cm]{geometry}

\usepackage{tablefootnote}

\usepackage[backend=biber, style=ieee]{biblatex}
\addbibresource{refs_howlong.bib}
\addbibresource{howlong_extra.bib}
\begin{filecontents*}[overwrite]{howlong_restart.bib}
@article{pal2017firstpassage,
  author  = {Pal, Arnab and Reuveni, Shlomi},
  title   = {First Passage under Restart},
  journal = {Physical Review Letters},
  volume  = {118},
  pages   = {030603},
  year    = {2017},
  doi     = {10.1103/PhysRevLett.118.030603}
}
\end{filecontents*}
\addbibresource{howlong_restart.bib}

\newtheorem{definition}{Definition}
\newtheorem{proposition}{Proposition}
\newtheorem{principle}{Principle}

\title{Sequential Achievement and the Sharp Crossover in Outlier Productivity}
\author{Ajax Benander}
\date{First drafted March 2026\\ {\small This version: September 2026 --- draft}}

\begin{filecontents}[overwrite]{howlong_extra.bib}
% Entries added at referee fix (Aug 2026); ALL VERIFY-flagged --- details from
% the referee model's memory; confirm each against the source at merge, then
% migrate to refs.bib and delete this block.
@article{merton1968matthew,
  author = {Merton, Robert K.},
  title = {The {M}atthew Effect in Science},
  journal = {Science}, volume = {159}, number = {3810},
  pages = {56--63}, year = {1968}}
@article{price1976cumulative,
  author = {Price, Derek de Solla},
  title = {A general theory of bibliometric and other cumulative advantage processes},
  journal = {Journal of the American Society for Information Science},
  volume = {27}, pages = {292--306}, year = {1976}}
@article{lotka1926,
  author = {Lotka, Alfred J.},
  title = {The frequency distribution of scientific productivity},
  journal = {Journal of the Washington Academy of Sciences},
  volume = {16}, pages = {317--323}, year = {1926}}
@article{oboyle2012,
  author = {O'Boyle, Jr., Ernest and Aguinis, Herman},
  title = {The best and the rest: revisiting the norm of normality of individual performance},
  journal = {Personnel Psychology}, volume = {65}, pages = {79--119}, year = {2012}}
@article{sinatra2016quantifying,
  author = {Sinatra, Roberta and Wang, Dashun and Deville, Pierre and Song, Chaoming and Barab{\'a}si, Albert-L{\'a}szl{\'o}},
  title = {Quantifying the evolution of individual scientific impact},
  journal = {Science}, volume = {354}, number = {6312}, pages = {aaf5239}, year = {2016}}
@article{luchins1942,
  author = {Luchins, Abraham S.},
  title = {Mechanization in problem solving: The effect of {E}instellung},
  journal = {Psychological Monographs}, volume = {54}, number = {6}, year = {1942}}
@article{bilalic2008,
  author = {Bilali{\'c}, Merim and McLeod, Peter and Gobet, Fernand},
  title = {Why good thoughts block better ones: The mechanism of the pernicious {E}instellung (set) effect},
  journal = {Cognition}, volume = {108}, pages = {652--661}, year = {2008}}
@article{clauset2009power,
  author = {Clauset, Aaron and Shalizi, Cosma Rohilla and Newman, M. E. J.},
  title = {Power-law distributions in empirical data},
  journal = {SIAM Review}, volume = {51}, pages = {661--703}, year = {2009}}
@article{ericsson1995ltwm,
  author = {Ericsson, K. Anders and Kintsch, Walter},
  title = {Long-term working memory},
  journal = {Psychological Review},
  volume = {102}, number = {2}, pages = {211--245}, year = {1995}
}
@article{chase1973perception,
  author = {Chase, William G. and Simon, Herbert A.},
  title = {Perception in chess},
  journal = {Cognitive Psychology},
  volume = {4}, number = {1}, pages = {55--81}, year = {1973}
}
@article{gobet1996templates,
  author = {Gobet, Fernand and Simon, Herbert A.},
  title = {Templates in chess memory: A mechanism for recalling several boards},
  journal = {Cognitive Psychology},
  volume = {31}, number = {1}, pages = {1--40}, year = {1996}
}
@misc{NexicFoundation, author={Benander, Ajax}, title={The Copernican Method Against the Copernican Principle: Anthropic Typicality over Partial Information States}, year={2026}, howpublished={Preprint}, url={https://philpapers.org/rec/BENTCM-7}}

@book{westfall1980never, author={Westfall, Richard S.}, title={Never at Rest: A Biography of {Isaac Newton}}, publisher={Cambridge University Press}, address={Cambridge}, year={1980}}

@misc{conduitt1727memoranda, author={Conduitt, John}, title={Memoranda Relating to {Sir Isaac Newton}}, howpublished={Keynes Ms.~130.02, King's College, Cambridge; transcribed by the Newton Project, THEM00165}, year={1727--1728}}

@article{keynes1995personality, author={Keynes, Milo}, title={The Personality of {Isaac Newton}}, journal={Notes and Records of the Royal Society of London}, volume={49}, number={1}, pages={1--56}, year={1995}, doi={10.1098/rsnr.1995.0001}}

@article{ohuma2023national, author={Ohuma, Eric O. and others}, title={National, Regional, and Global Estimates of Preterm Birth in 2020, with Trends from 2010: A Systematic Analysis}, journal={The Lancet}, volume={402}, number={10409}, pages={1261--1271}, year={2023}}

@article{bhutta2002cognitive, author={Bhutta, Adnan T. and Cleves, Mario A. and Casey, Patrick H. and Cradock, Mary M. and Anand, K. J. S.}, title={Cognitive and Behavioral Outcomes of School-Aged Children Who Were Born Preterm: A Meta-analysis}, journal={JAMA}, volume={288}, number={6}, pages={728--737}, year={2002}, doi={10.1001/jama.288.6.728}}

@article{twilhaar2018cognitive, author={Twilhaar, E. Sabrina and others}, title={Cognitive Outcomes of Children Born Extremely or Very Preterm Since the 1990s and Associated Risk Factors: A Meta-analysis and Meta-regression}, journal={JAMA Pediatrics}, volume={172}, number={4}, pages={361--367}, year={2018}}

@article{eves2021association, author={Eves, Robert and others}, title={Association of Very Preterm Birth or Very Low Birth Weight With Intelligence in Adulthood: An Individual Participant Data Meta-analysis}, journal={JAMA Pediatrics}, volume={175}, number={8}, pages={e211058}, year={2021}, doi={10.1001/jamapediatrics.2021.1058}}

@article{marlow2005neurologic, author={Marlow, Neil and Wolke, Dieter and Bracewell, Melanie A. and Samara, Muthanna}, title={Neurologic and Developmental Disability at Six Years of Age after Extremely Preterm Birth}, journal={New England Journal of Medicine}, volume={352}, number={1}, pages={9--19}, year={2005}}

@article{hallin2010followup, author={Hallin, Anna-Lena and Hellstr{\"o}m-Westas, Lena and Stjernqvist, Karin}, title={Follow-up of Adolescents Born Extremely Preterm: Cognitive Function and Health at 18 Years of Age}, journal={Acta Paediatrica}, volume={99}, number={9}, pages={1401--1406}, year={2010}}

@book{isaacson2007einstein, author={Isaacson, Walter}, title={Einstein: His Life and Universe}, publisher={Simon \& Schuster}, address={New York}, year={2007}}

@article{eisenstadt1978parental, author={Eisenstadt, J. Marvin}, title={Parental Loss and Genius}, journal={American Psychologist}, volume={33}, number={3}, pages={211--223}, year={1978}, doi={10.1037/0003-066X.33.3.211}}

@incollection{damian2014diversifying, author={Damian, Rodica Ioana and Simonton, Dean Keith}, title={Diversifying Experiences in the Development of Genius and Their Impact on Creative Cognition}, booktitle={The Wiley Handbook of Genius}, editor={Simonton, Dean Keith}, publisher={Wiley-Blackwell}, year={2014}}

@article{gobet1996recall, author={Gobet, Fernand and Simon, Herbert A.}, title={Recall of Rapidly Presented Random Chess Positions Is a Function of Skill}, journal={Psychonomic Bulletin \& Review}, volume={3}, number={2}, pages={159--163}, year={1996}}

\end{filecontents}

\begin{document}

\maketitle

\begin{abstract}
Bessis \cite{bessis2024mathematica,bessis2025genes} observes that mathematical talent, as perceived by practicing mathematicians, follows a Pareto-like distribution with orders-of-magnitude gaps between individuals, and argues that this tells against genetic determinism as a primary explanation, an inference that needs care, since a nonlinear map from inputs to achievement can change the output distribution whatever caused the inputs. The question of how a Gaussian input could yield a Pareto-like output is left there as an informal conjecture. Taking Ramanujan as a motivating example, this paper proposes a quantitative mechanism. In Part~I, we develop a method for inferring the parameters of a population's ability distribution from the observable characteristics of an elite subpopulation and a single extreme observation, and show that the scaling character of the rank-$\#1$ residual serves as a two-stage diagnostic that makes the log-normal a testable candidate selected against two comparators rather than a convention: proportionality of the overshoot separates it from the linear-Gaussian alternative, and the overshoot's attenuation with population size separates it from power-law alternatives, the diagnostic being a model test against those comparators and not an identification of the family against every alternative. In Part~II, we model complex achievement as a sequential stochastic process requiring consecutive successes, where any single failure resets progress. We derive a sharp finite-size crossover in expected completion time: below a depth-dependent crossover score $s_c(n) \sim \ln n/(-\ln(1-p_0))$, completion time grows exponentially in the task depth; above it, completion time is approximately linear; and at the crossover the overhead over the ideal is an explicit scaling function of the depth-scaled score. This consecutive-success mechanism is the bridge: in the low-success regime the rate is approximately log-normal, its log-standard deviation amplified by the task depth and its log-variance by the square of the depth, so that modest, continuous differences become differences of many orders of magnitude in realized output, a log-normal broad enough to closely mimic power-law behavior over a substantial finite empirical range, the Pareto-like appearance Bessis identifies, whether a given window distinguishes the two being an empirical question of its width, sample size, and the amplified variance; at high scores the exact rate saturates, bounded above by $\lambda/n$, so the transformed distribution is not globally log-normal. We propose a realist interpretation of the underlying score as \emph{divergence resistance}---the absence of a common tendency to extrapolate inferential structures incorrectly---inverting the standard framing by treating the outlier not as possessing a rare gift but as lacking a common impediment. An illustrative calibration on stipulated anchors, not an empirical validation, shows that the framework simultaneously demystifies the outlier, who is a draw from the same distribution as everyone else, and, under its full-restart assumption, makes their irreplaceability quantitative and practical: modest reliability differences can separate completion probabilities so widely that, within a specified population and time budget, additional lower-reliability workers provide negligible compensation for losing a higher-reliability one.
\end{abstract}

\section{Introduction}\label{sec:intro}

In December 1919, Srinivasa Ramanujan lay dying in a nursing home in
Madras. He was thirty-two. In the preceding five years he had produced a
body of work in pure mathematics (partition identities, modular forms,
$q$-series, mock theta functions, and much else) whose depth and
originality remain, a century later, undiminished and actively
generative: entire subfields of modern number theory trace their lineage
to results that Ramanujan recorded without proof in notebooks that
mathematicians are still unpacking today
\cite{hardy1940ramanujan,kanigel1991man}.

How does one account for this?

The standard framing offers two unsatisfying extremes. On the one hand, a
na\"{\i}ve reading of the ``great man'' model \cite{hook2004great} would treat Ramanujan as
categorically distinct from the rest of humanity, a framing whose
historical entanglement with the systematic exclusion of marginalized
groups from education and scientific institutions makes it both
inadequate and harmful. On the other hand, the strong
inevitability thesis \cite{kuhn1970structure}, which holds that any sufficiently developed mathematical
community would have produced the same results eventually, sits uneasily
with the fact that much of Ramanujan's output has taken decades of
collective effort merely to verify, let alone to have independently
discovered.

A recent contribution to this debate comes from Bessis \cite{bessis2024mathematica,bessis2025genes,bessis2026jaimungal}, who observes that the perceived distribution of mathematical talent resembles a Pareto distribution---with orders-of-magnitude gaps between individuals---rather than the Gaussian distribution expected of highly heritable polygenic traits like height. Bessis argues that this mismatch tells against genetic determinism as the primary explanation for mathematical genius---an inference that needs care, since a nonlinear transformation from latent inputs to achievement can change the output distribution whatever caused the inputs, and the amplification mechanism proposed here is an instance of one, so a Gaussian-versus-Pareto-like contrast alone does not rule out a genetic contribution--- and proposes instead that the gap arises from ``cognitive transformation and neuroplasticity''---a capitalization process in which abstractions are stacked upon abstractions over years of practice. However, Bessis stops short of providing a formal mechanism: the question of {how} a Gaussian input (ability) becomes a Pareto-like output (achievement) is left as an informal conjecture about the brain's connectome. The present paper proposes to fill this gap.
Specifically, this paper proposes a quantitative account. We argue that extreme outlier productivity of the kind exemplified by Ramanujan can be modeled as a {quantitative consequence} of two interacting features: (i)~a log-normal
distribution of latent ability across the population, whose form we propose to diagnose from the empirical scaling behavior of record-holder residuals, through an illustrative inference scheme rather than an analyzed dataset with a selection model, instead of assuming it by convention; and (ii)~a sequential-success
mechanism governing the production of complex output, in which the
individual must succeed at a chain of steps consecutively and in order,
with any failure resetting the process from the beginning.

The central result is that these two features compose to produce a sharp
{crossover} in expected completion time. Below a critical
score, the expected time to complete a task of depth~$n$ grows
exponentially in~$n$; above the crossover score, the expected time is
approximately linear in~$n$, essentially routine. The transition between these
regimes is governed by the individual's position in the log-normal
ability distribution, and the amplification is polynomial in the
individual's score with exponent equal to the task depth. This means
that modest differences in latent ability, fully
consistent with a smooth, continuous human distribution, translate into
differences of many orders of magnitude in realized output.
The crucial philosophical reframe is this: the mechanism does not
distinguish between individuals. Everyone, from the population median
to Ramanujan, is running the same consecutive-success process with the
same structural rules. Nobody is categorically unable to complete the
task. What differs is the {time complexity} of completion: a task
that takes three weeks for an individual above the crossover score takes
longer than the age of the universe for an individual below it. The
apparent categorical difference between outlier producers and everyone
else is therefore not a difference in {kind} but a quantitative
consequence of a universal mechanism applied uniformly to all, whose
mathematical structure converts continuous differences in input into
discrete-looking differences in output.

The explanatory payoff is twofold. First, the model {demystifies}
the outlier: Ramanujan is not a categorical alien but a draw from the
same distribution as everyone else, located at an extreme but finite
quantile.
Second, and simultaneously, under its full-restart assumption the model {makes the irreplaceability} of the outlier quantitative: modest reliability differences can produce completion probabilities so widely separated that, within a specified population and time budget, additional individuals below the crossover score provide negligible compensation, the finite-horizon completion probability of \eqref{eq:horizon} saying how negligible---so that the loss of such an
individual represents a genuine, quantifiable deficit to the collective
record. These two conclusions---continuity of the individual with the
population, and irreproducibility of the individual's output by the
collective---are usually treated as contradictory. The present framework
shows them to be simultaneous and inseparable consequences of the same
mathematics.

The paper is organized as follows.
Section~\ref{sec:inference} develops a method for inferring the
parameters of a log-normal ability distribution from observable
characteristics of an elite subpopulation and a single extreme
observation, and selects the log-normal specification, as a testable candidate against the linear-Gaussian and power-law comparators, from the scaling character of rank-$\#1$ residuals.
Section~\ref{sec:completion} introduces the consecutive-success process,
derives the crossover in expected completion time, and connects
the score from Section~\ref{sec:inference} to the amplification ratio
that governs the concentration of output among top performers.
Section~\ref{sec:results} applies the framework to specific empirical
domains, instantiating the general methodology with concrete data.


\section{Inference of Parent-Population Parameters from Elite Subpopulation
         Characteristics}\label{sec:inference}

\subsection{Motivation: Diagnosing the Distributional Family from
            Record-Holder Residuals}\label{ssec:motivation}

Before developing the formal framework, we present an empirical argument
that constrains the distributional family of the trait under study. The
argument proceeds in five steps.

\medskip
\noindent\textbf{Step 1: Fitting the elite order statistics.}\;
Consider the top $m$ performers in a population~$A$, ranked by a
quantitative trait~$X$. One fits a smooth curve (either a parametric
order-statistics model or a nonparametric interpolant) to the values of
ranked individuals $\#2$ through $\#m$, and extrapolates to predict the
expected value at rank~$\#1$.

\medskip
\noindent\textbf{Step 2: Observing the residual at rank $\#1$.}\;
In many competitive domains one finds that, during epochs in which a
dominant individual is active (e.g.\ a period containing a figure such as
Bolt in sprinting or Carlsen in chess), the actual rank-$\#1$ value
systematically overshoots the extrapolated curve. Denote this overshoot
by $\Delta$.

\medskip
\noindent\textbf{Step 3: Attribution to generational selection bias.}\;
The overshoot is naturally explained as a selection effect: the analyst
tends to perform the analysis precisely during those epochs in which an
outlier is present, and such an outlier is effectively the maximum drawn not
from the contemporaneous population of size~$N$ alone but from a larger
cumulative pool across multiple generations. Let $K = N_{\mathrm{eff}} / N$
denote the ratio of this effective population to the contemporaneous one.
Since the expected number of individuals exceeding a level~$x$ is
$N \cdot \Pr(X > x)$, observing one individual above a higher level
requires proportionally more draws, and the overshoot~$\Delta$ is
consistent with the rank-$\#1$ slot being filled from a pool of
size~$K N$.

\medskip
\noindent\textbf{Step 4: Diagnosing the scaling character of
$\boldsymbol{\Delta}$.}\;
Under extreme-value theory, enlarging the population from $N$ to $KN$
shifts the typical maximum by an amount that is \emph{additive in the
natural scale of the distribution}
\cite{leadbetter1983,embrechts1997}:
\[
  \mathbb{E}[Y_{(KN)}] - \mathbb{E}[Y_{(N)}]
  \;\approx\;
  \sigma\bigl(z_{\max}(KN) - z_{\max}(N)\bigr),
\]
where $Y$ is the trait expressed in the scale in which it is Gaussian,
$\sigma$ is the standard deviation in that scale, and $z_{\max}(\cdot)$ is
the EVT asymptotic (\ref{eq:zmax} below). If $Y = X$ (Gaussian in
original scale), the overshoot~$\Delta$ is a fixed additive quantity
depending on $\sigma$ and $K$ but not on the current performance level. If
instead $Y = \ln X$ (Gaussian in log-scale), the overshoot in log-space is
$\delta = \sigma(z_{\max}(KN) - z_{\max}(N))$, which translates to an
overshoot in original scale of
\[
  \Delta
  \;=\;
  X_{(N)}\bigl(e^{\delta} - 1\bigr)
  \;\propto\;
  X_{(N)}\,,
\]
i.e.\ a \emph{proportional} overshoot that scales with the current level.

\medskip
\noindent\textbf{Step 5: Conclusion.}\;
If the empirically observed $\Delta$ is proportional to the performance
level rather than constant across different levels or domains, the pattern is consistent with Gaussian structure in log-space, which the paper adopts as its working specification:
\begin{equation}\label{eq:deduction}
  \Delta \propto X_{(N)} \quad\text{is consistent with}\quad \ln X \sim \mathcal{N}(\mu,\sigma^2),
\end{equation}
The log-normal specification is therefore not an assumption of convenience
but a \emph{selection} forced by the scaling character of rank-$\#1$
residuals, at the grade fixed by the scope remark below: proportionality
places the Gaussian structure in log-space or heavier, and the
$N$-attenuation of the log-overshoot then separates the log-normal from
its power-law rivals\footnote{
I believe it is also worth mentioning a note about the intuition for this idea, as I see it: effectively it is an instance of selection bias. Because a ``greatest of all time" person with the height of Robert Wadlow or the speed of Usain Bolt, etcetera, only comes around once every several decades or even centuries in some cases, when we look at the score rankings of the top individuals at said person's peak, we are actually looking at a sample with the bias of ``having to exist in time during which this person is operating". Tangibly speaking, if we look at the 100-meter dash after the retirement of Usain Bolt, suddenly the new winner of the race has hardly any residual at all from the extrapolated curve of the other participants. In this sense, it is as if Usain Bolt is drawn from a ``best of the century" widened distribution of people, whereas the second-place and after holders are drawn from a ``best of the decade" distribution. In most races, the first-place residual is, as expected, roughly zero, but for it to be higher than this is itself a capturing of the selection bias which necessarily comes from comparing any ``greatest of all time" individual with their contemporaries. \par
Given this setup itself, my instinct with actually carrying out the statistical analysis is to use the standard high-ranking Gumbel procedure, because my direct intuition would force through an arbitrary threshold of ``draw the top person from the wider-timescale distribution and others from the contemporary-timescale distribution", which is messy in the sense of nonlocality due to forcing about more global distributions. Hence, the reason why the Gumbel process version is superior, (which, note, I had not heard of while first devising this, I just had the intuition to look online for a ``top members of a lognormal distribution for a finite population" distribution due to the nonlocality concern), in the sense that we can re-encode the same information of globally-widening the population from which the top person is drawn by virtue of simply tweaking that top person's quantile in the Gumbel distribution such that, roughly speaking, the distance to one is shrunk in proportion to the timescale over which they operate. In effect, this allows us to capture the nonlocal intuition with a perfectly local methodological deduction, and in doing so allows the distribution itself to pick the ``threshold" aspect rather than having to force it through as an arbitrarily-chosen free parameter of sorts. 
The interesting thing is that in addition to getting rid of the arbitrary threshold concept, this also essentially encodes a naturally-nonzero residual in the top person anyway, in the sense that there is a continuous interpolation in ``their natural timescale population from which they are drawn" between perfectly contemporary and being measured in centuries. For instance, even if the top person say fits to a ``decade-longer natural population", then we would see that reflected in their residual amongst contemporaries perturbing up just a little bit. \par
This is also relevant to the Ramanujan analysis later on in the paper, where we consider his ``intrinsic score" to be extra-boosted from someone like Bolt due to Bolt being the best over the timescale of the Olympics (which of course is still quite good), while Ramanujan is the best over the timescale of, effectively, the entire written history of pure mathematics. Hence, even pertaining to his intrinsic score Ramanujan gets an extreme residual boost, which then compounds exponentially when discussing his total sequential achievement-pertaining score, (which all contributes to his so vastly outpacing anyone else we know of in the field).
}.

\medskip

The individual ingredients of this argument---log-normal models
of positive traits \cite{gibrat1931}, Gumbel limits for maxima
\cite{fisher1928,gnedenko1943}, and the use of order-statistics residuals to
diagnose distributional families \cite{embrechts1997}---are of course all standard practice.
The intention of the present argument lies in the specific inferential
chain: running from the empirical scaling of the rank-$\#1$ overshoot
through the population-ratio mechanism to the identification of the
log-normal family. Where the standard EVT pipeline classifies domains of
attraction given a known distribution, the present argument proposes to reverse
the arrow, using the empirical signature of the extreme to identify the
distribution.

\medskip
\noindent\textbf{Remark (scope of the diagnostic).}\;
Proportionality of the overshoot separates the linear-Gaussian family
(additive overshoot) from the \emph{class} of multiplicative-tail
families, of which the log-normal is one member: a regularly varying
(Pareto-type) tail also yields an overshoot proportional to the current
level. The two members of the multiplicative class are then separated by
the $N$-dependence of the \emph{log-space} overshoot~$\delta$: for a
log-normal parent, $\delta \approx \sigma \ln K / \sqrt{2\ln N}$,
attenuating as the population grows, whereas a Pareto tail of
index~$\alpha$ gives $\delta = \ln K / \alpha$, independent of~$N$. The
diagnostic is accordingly two-stage---proportionality first, attenuation
second---and the second stage is not decorative: the thesis of Section~\ref{ssec:pareto} is that a log-normal \emph{input} can produce an \emph{output} broad enough to look Pareto-like over a finite window, so the input-side diagnostic must be able to tell the two families apart, and it can.
(For the general methodology of discriminating distributional families
from tail data, see \cite{clauset2009power}.)

\subsection{Setup}\label{ssec:setup}

Let $A$ denote a parent population of size $N$ in which a quantitative
trait~$X$ takes strictly positive values and is log-normally distributed:
\begin{equation}\label{eq:parent}
  X \sim \operatorname{LogNormal}(\mu,\,\sigma^{2}),
\end{equation}
so that the log-transformed trait $Y = \ln X$ is Gaussian:
\begin{equation}\label{eq:logtrait}
  Y \sim \mathcal{N}(\mu,\,\sigma^{2}).
\end{equation}
All distributional identities in the sequel are standard normal results
applied to~$Y$.

Let $B \subset A$ denote an elite subpopulation consisting approximately of
the individuals with the largest values of~$X$. Under a sharp selection
model in which $B$ comprises exactly those individuals exceeding some
threshold~$\tau$ in original scale, the log-transformed trait conditional on
$Y > \ln\tau$ follows a left-truncated normal distribution. Writing the
selection fraction as $p = \Pr(X > \tau)$, the log-threshold satisfies
\begin{equation}\label{eq:threshold}
  \ln\tau \;=\; \mu + \sigma\,\Phi^{-1}(1-p)\,,
\end{equation}
where $\Phi^{-1}$ is the standard-normal quantile function. The truncated
distribution of~$Y$ has mean
\begin{equation}\label{eq:trunc-mean}
  \mu_B
  \;=\;
  \mu + \sigma\,\frac{\varphi\!\bigl(\Phi^{-1}(1-p)\bigr)}{p}\,,
\end{equation}
where $\varphi$ denotes the standard-normal density, and variance
\begin{equation}\label{eq:trunc-var}
  \sigma_B^{2}
  \;=\;
  \sigma^{2}\!\left[
    1
    \;+\;
    \frac{\Phi^{-1}(1-p)\;\varphi\!\bigl(\Phi^{-1}(1-p)\bigr)}{p}
    \;-\;
    \left(\frac{\varphi\!\bigl(\Phi^{-1}(1-p)\bigr)}{p}\right)^{\!2}
  \right].
\end{equation}

The truncated normal is not itself Gaussian; it is right-skewed, truncated on the left with its tail extending right, and has its mode at the truncation point~$\ln\tau$ rather than at~$\mu_B$. In practice, however, the empirical distribution of an elite
subpopulation often appears approximately normal in log-space, a phenomenon
attributable to the fact that selection in realistic settings is rarely a
sharp univariate threshold but instead involves multiple correlated criteria,
measurement noise, and stochastic developmental factors. These
considerations soften the truncation boundary and yield an effective
distribution closer to Gaussian. Nothing below needs $B$'s central tendency: the quantile condition~\eqref{eq:eq1} uses the logarithm of $B$'s mode, $\ln m_B$, directly, and that is not in general the mode of $\ln X$ on~$B$, the two differing by the Jacobian of the transformation.

\subsection{Observed Data and Unknowns}\label{ssec:data}

We assume the following quantities are observed or estimable:
\begin{enumerate}
  \item[\textbf{D1.}]
    The mode of~$B$ in original scale, denoted $m_B$.
  \item[\textbf{D2.}]
    The quantile of $m_B$ within the parent population~$A$, denoted~$q$,
    so that $\Pr(X \le m_B) = q$.
  \item[\textbf{D3.}]
    A single extreme observation $v_{\max}$, representing the largest
    recorded value across all of~$A$. We parametrize this as
    $v_{\max} = k\,m_B$ for a known constant~$k > 1$.
  \item[\textbf{D4.}]
    The size $N$ of the parent population~$A$.
\end{enumerate}
The unknowns to be recovered are the log-space parameters $\mu$
and~$\sigma$. The size~$n$ of the elite subpopulation is not
required for the primary inference; it enters only if one additionally
wishes to characterize the spread of~$B$ via $\sigma_B$ (see
Section~\ref{ssec:sigmaB}).

\subsection{Derivation of the Log-Space Parameters}\label{ssec:derivation}

All identities in this subsection are stated in log-space, where the
Gaussian structure of~$Y = \ln X$ is directly available.

\subsubsection{First Equation: Quantile Condition}

Since $\Pr(X \le m_B) = q$ is equivalent to
$\Pr(Y \le \ln m_B) = q$, we have
\begin{equation}\label{eq:eq1}
  \ln m_B \;=\; \mu + \sigma\,z_q\,,
  \qquad
  z_q \;\equiv\; \Phi^{-1}(q).
\end{equation}

\subsubsection{Second Equation: The Typical Maximum}

The maximum order statistic $Y_{(N)} = \ln X_{(N)}$ of $N$ independent
draws from $\mathcal{N}(\mu,\sigma^2)$ has, for large~$N$, the Gumbel limit $Y_{(N)} \approx \mu + \sigma\big(z_{\max} + G/\sqrt{2\ln N}\big)$ with $G$ a standard Gumbel variable \cite{leadbetter1983,david2003}, so that its most likely value, the centering, is
\begin{equation}\label{eq:emax}
  \operatorname{mode}\big(Y_{(N)}\big)
  \;\approx\;
  \mu + \sigma\,z_{\max}\,,
\end{equation}
while its mean exceeds this by $\sigma\gamma_E/\sqrt{2\ln N}$, with $\gamma_E \approx 0.5772$ the Euler--Mascheroni constant. The estimates below match an observed maximum to this typical value, a \emph{typical-maximum matching} estimate. It is not in general the maximum-likelihood estimate here: the mode is the maximum-likelihood location only when the scale is known, and with $\sigma$ inferred jointly with $\mu$ under the quantile condition the likelihood carries a factor $1/\sigma$, so matching the standardized maximum to its mode does not maximize it; where the calibration is to be read as inference, the order-statistic likelihood of Section~\ref{ssec:uncertainty} is to be optimized directly. Matching the maximum to the mean instead lowers the calibrated dispersion of Section~\ref{ssec:physics} from $0.5911$ to $0.5682$, a sensitivity every downstream quantity carries. Here
\begin{equation}\label{eq:zmax}
  z_{\max}
  \;=\;
  \sqrt{2\ln N}
  \;-\;
  \frac{\ln(\ln N) + \ln(4\pi)}{2\sqrt{2\ln N}}\,.
\end{equation}
Since $v_{\max} = k\,m_B$ in original scale, the log-maximum satisfies
$\ln v_{\max} = \ln k + \ln m_B$. Equating with its expectation gives
\begin{equation}\label{eq:eq2}
  \ln k + \ln m_B
  \;\approx\;
  \mu + \sigma\,z_{\max}\,.
\end{equation}
We discuss the uncertainty inherent in this single-observation
identification in Section~\ref{ssec:uncertainty}.

\subsubsection{Solution}

Subtracting \eqref{eq:eq1} from \eqref{eq:eq2} eliminates $\mu$ and yields
\begin{equation}\label{eq:sigma-sol}
  \boxed{\;
    \sigma
    \;=\;
    \frac{\ln k}{\,z_{\max} - z_q\,}
  \;}
\end{equation}
provided $z_{\max} > z_q$, which holds whenever the population~$A$ is large
enough that its typical maximum exceeds the elite mode in the upper
tail, a condition satisfied for any practically relevant $N$ and~$q$.

Back-substituting into \eqref{eq:eq1} gives
\begin{equation}\label{eq:mu-sol}
  \boxed{\;
    \mu
    \;=\;
    \ln m_B \;-\; \frac{z_q\,\ln k}{\,z_{\max} - z_q\,}
  \;}
\end{equation}

Note that, in contrast to a linear-scale Gaussian model, no constraint on
the sign of~$\mu$ or on the relationship between $k$ and the ratio
$z_{\max}/z_q$ is required. Since $X = e^Y > 0$ by construction, the
log-normal framework is consistent for all real values of~$\mu$ and all
$k > 1$. This is one of the principal practical advantages of the
log-normal formulation over the linear-scale alternative.

\subsubsection{Original-Scale Summary Statistics}

The log-normal parameters $(\mu,\sigma)$ determine all moments in
original scale. For reference:
\begin{equation}\label{eq:orig-stats}
  \mathbb{E}[X] = e^{\mu + \sigma^2/2},
  \qquad
  \operatorname{Var}(X) = e^{2\mu+\sigma^2}\!\bigl(e^{\sigma^2}-1\bigr),
  \qquad
  \operatorname{mode}(X) = e^{\mu - \sigma^2}.
\end{equation}

\subsection{Spread of the Elite Subpopulation}\label{ssec:sigmaB}

If the size~$n$ of the elite subpopulation is known, one may define the
selection fraction $p = n/N$ and obtain the log-space standard deviation
of~$B$ by substituting $\sigma$ from \eqref{eq:sigma-sol} into
\eqref{eq:trunc-var}:
\begin{equation}\label{eq:sigmaB-sol}
  \sigma_B^{2}
  \;=\;
  \sigma^{2}\!\left[
    1
    + \frac{\alpha\,\varphi(\alpha)}{p}
    - \left(\frac{\varphi(\alpha)}{p}\right)^{\!2}
  \right],
  \qquad
  \alpha \equiv \Phi^{-1}(1-p).
\end{equation}

\subsection{Consistency Conditions}\label{ssec:consistency}

The log-normal framework imposes only mild requirements on the data:
\begin{enumerate}
  \item[\textbf{C1.}]
    $k > 1$, i.e.\ the record-holder exceeds the elite mode. This
    ensures $\sigma > 0$ in \eqref{eq:sigma-sol}.
  \item[\textbf{C2.}]
    $z_{\max} > z_q$, i.e.\ the typical maximum of a sample of
    size~$N$ exceeds the $q$-th percentile in standardized units. For
    any $q < 1$ and $N$ of practical size this is automatically
    satisfied.
  \item[\textbf{C3.}]
    \emph{Positivity of the truncated variance} (relevant only when $n$
    is supplied).\;
    The bracketed expression in \eqref{eq:sigmaB-sol} is automatically
    positive for all $p \in (0,1)$: it is the ratio of truncated to
    parent variance, which lies strictly between $0$ and $1$. No
    additional condition is imposed.
\end{enumerate}

\subsection{Quantification of Uncertainty and the Record-Holder
            Amplification Effect}\label{ssec:uncertainty}

\subsubsection{Log-Space Uncertainty}

The principal source of uncertainty in the preceding estimates is the identification of $\ln v_{\max}$ with its typical value, and one consequence of that identification must be respected downstream. Once the observed maximum has been set equal to the typical maximum to solve for $\sigma$, the record-holder is, under the inferred parameters, the \emph{typical} record-holder; a second, independent Gumbel fluctuation may not then be added to move the record-holder above the typical maximum, since that counts the same observation's extremeness twice. Where a realization above the typical maximum is wanted for illustration, the calibration below flags it as an illustrative choice and not as an inference. The clean statistical route, used in place of the two-step procedure wherever the calibration is to be read as inference, is the order-statistic likelihood: the observed maximum $y$ of a sample of size $N$ enters once, through $f_{Y_{(N)}}(y \mid \mu, \sigma) = N f_Y(y \mid \mu, \sigma)\, F_Y(y \mid \mu, \sigma)^{N-1}$, jointly with whatever independent information constrains $(\mu, \sigma)$, and the resulting profile likelihood carries the maximum's uncertainty without a second fluctuation being introduced by hand. The log-maximum
$Y_{(N)}$ converges in distribution to a Gumbel variate
\cite{fisher1928,gnedenko1943}, with scale parameter
\begin{equation}\label{eq:gumbel-scale}
  \beta_N
  \;=\;
  \frac{\sigma}{\sqrt{2\ln N}}\,.
\end{equation}
The standard deviation of the Gumbel distribution is
$(\pi/\sqrt{6})\,\beta_N$, giving
\begin{equation}\label{eq:sd-logmax}
  \operatorname{sd}(Y_{(N)})
  \;\approx\;
  \frac{\pi\,\sigma}{\sqrt{12\,\ln N}}\,.
\end{equation}
Propagating through \eqref{eq:sigma-sol} by a first-order expansion, with the log-maximum the uncertain quantity, $\delta\sigma \approx \delta Y_{(N)}/(z_{\max} - z_q)$, so the relative uncertainty in $\sigma$ is
\begin{equation}\label{eq:rel-uncert-sigma}
  \frac{\delta\sigma}{\sigma}
  \;\approx\;
  \frac{\pi}{\sqrt{12\,\ln N}\,(z_{\max} - z_q)}\,,
\end{equation}
the denominator $z_{\max} - z_q$ carried from \eqref{eq:sigma-sol} and not to be dropped---at the calibration below it is $\approx 2.12$, which halves the figure the bare Gumbel scale would suggest---and the absolute uncertainty in $\mu$ is
\begin{equation}\label{eq:abs-uncert-mu}
  \delta\mu
  \;=\;
  |z_q|\;\delta\sigma
  \;=\;
  \frac{\pi\,|z_q|\,\sigma}{\sqrt{12\,\ln N}\,(z_{\max} - z_q)}\,.
\end{equation}

\begin{table}[ht]
  \centering
  \caption{Relative uncertainty in $\sigma$ as a function of parent population size~$N$, at the calibrated $z_{\max} - z_q \approx 2.12$; the bare Gumbel scale $\pi/\sqrt{12\ln N}$ is shown for comparison.}
\label{tab:uncert}
\medskip
\begin{tabular}{r c c c}
\hline
$N$ & $\sqrt{12\,\ln N}$ & $\pi/\sqrt{12\ln N}$ & $\delta\sigma/\sigma$ \\
\hline
$10^{3}$ & $9.10$ & $34.5\%$ & $16.3\%$ \\
$10^{6}$ & $12.87$ & $24.4\%$ & $11.5\%$ \\
$10^{9}$ & $15.76$ & $19.9\%$ & $9.4\%$ \\
$10^{12}$ & $18.19$ & $17.3\%$ & $8.1\%$ \\
\hline
\end{tabular}
\end{table}

\subsubsection{Multiplicative Amplification in Original Scale}
\label{ssec:bolt}

The record-holder amplification effect identified in
Section~\ref{ssec:motivation} now admits a precise quantitative treatment.
The observed maximum in original scale is
\begin{equation}\label{eq:max-orig}
  X_{(N)}
  \;=\;
  \exp\!\bigl(Y_{(N)}\bigr)
  \;\approx\;
  e^{\mu + \sigma z_{\max}} \cdot e^{\varepsilon}\,,
\end{equation}
where $\varepsilon$ is a Gumbel-distributed random variable with standard
deviation given by \eqref{eq:sd-logmax}. The multiplicative factor
$e^{\varepsilon}$ has, to first order, a multiplicative uncertainty scale---not the exact coefficient of variation of the exponentiated Gumbel, which would require its moment expression, but the first-order scale that suffices here---
\begin{equation}\label{eq:bolt-spread}
  \operatorname{CV}(X_{(N)})
  \;\approx\;
  \exp\!\!\left(\frac{\pi\,\sigma}{\sqrt{12\,\ln N}}\right) - 1\,.
\end{equation}

This is precisely the mechanism underlying Steps~3--4 of
Section~\ref{ssec:motivation}. The Gumbel fluctuation~$\varepsilon$ is
additive in log-space and has magnitude independent of the current
performance level; exponentiation converts it into a proportional overshoot
in original scale. The rank-$\#1$ residual~$\Delta$ identified empirically
in Step~2 is thus
\begin{equation}\label{eq:residual}
  \Delta
  \;=\;
  X_{(N)}\bigl(e^{\varepsilon} - 1\bigr)
  \;\propto\; X_{(N)}\,,
\end{equation}
confirming the proportional scaling that, per Step~5, diagnoses the
log-normal family.

Crucially, the entire framework uses a single population~$A$ of
contemporaneous size~$N$. The generational selection bias described in
Step~3 is not required as a separate mechanism: it is \emph{equivalent} to
the Gumbel fluctuation. Enlarging the population from $N$ to $KN$ shifts
$z_{\max}$ upward by a modest amount; a single Gumbel draw landing
$\varepsilon$ above the mean of the $N$-population maximum produces an
identical effect in original scale. The two descriptions (``the outlier
comes from a larger pool'' and ``the outlier is a routine upward fluctuation
amplified through exponentiation'') are dual, but the latter is more
parsimonious and avoids the complications of defining a stationary trait
distribution across historical cohorts.

\subsubsection{Inter-Rank Spacing and the Gap Distribution}
\label{ssec:gap}

The preceding analysis characterizes the fluctuation of the
maximum around its expected value. Equally important for the
consecutive-success model of Section~\ref{sec:completion} is the
\emph{spacing} between the record-holder and their nearest
competitor, since it is this gap that the crossover amplifies
into the completion-time separation documented in
Section~\ref{sec:results}.

The Poisson-process representation of extreme order statistics
\cite{resnick1987} gives the following result. Let
$X_{(1)} > X_{(2)} > \cdots$ denote the top order statistics of $N$
independent draws from the log-normal population. Their log-space
spacings are
\begin{equation}\label{eq:logspace-spacing}
  Y_{(j)} - Y_{(j+1)}
  \;=\;
  \ln X_{(j)} - \ln X_{(j+1)}
  \;\sim\;
  \frac{\sigma\, b_N}{j}\, E_j\,,
\end{equation}
where $E_1, E_2, \ldots$ are independent standard exponential random
variables, $b_N = 1/\sqrt{2\ln N}$ is the Gumbel scale of the standard Gaussian, so that $\sigma b_N$ is that of $Y$, and the approximation becomes exact in the $N \to
\infty$ limit. Converting to the score ratio between the
record-holder and the second-ranked individual:
\begin{equation}\label{eq:score-gap}
  \frac{s_{(1)}}{s_{(2)}}
  \;=\;
  \exp\!\bigl(\sigma\,b_N\, E_1\bigr)
  \;=\;
  \exp\!\bigl(\beta_N\, \Delta\bigr)\,,
\end{equation}
where $\Delta \sim \operatorname{Exp}(1)$ and
$\beta_N = \sigma\,b_N = \sigma/\sqrt{2\ln N}$ is the same Gumbel
scale parameter~\eqref{eq:gumbel-scale} that governs the uncertainty
of the maximum.

Equation~\eqref{eq:score-gap} yields three immediate consequences.
First, the \emph{median} gap between the record-holder and the
second-ranked individual is
\begin{equation}\label{eq:median-gap}
  \operatorname{med}\!\left(\frac{s_{(1)}}{s_{(2)}}\right)
  \;=\;
  e^{\,\beta_N \ln 2}\,,
\end{equation}
which evaluates to approximately $1.063$, or a $6.3\%$ advantage,
at $\sigma = 0.591$ and $N = 8 \times 10^9$. Second, the probability
that the gap exceeds a fraction~$g$ is
\begin{equation}\label{eq:gap-tail}
  P\!\left(\frac{s_{(1)}}{s_{(2)}} > 1 + g\right)
  \;=\;
  (1+g)^{-1/\beta_N}\,,
\end{equation}
since $P(\Delta > x) = e^{-x}$ and $(1+g)^{-1/\beta_N} =
e^{-\ln(1+g)/\beta_N}$. At the present calibration
$1/\beta_N = \sqrt{2\ln N}/\sigma \approx 11.4$, giving
$P(\text{gap} > 10\%) \approx 34\%$ and
$P(\text{gap} > 20\%) \approx 12\%$. Gaps of this magnitude are not
rare flukes; they are routine features of the Gumbel process.

Third, the spacings between \emph{lower}-ranked individuals compress:
the gap between ranks $j$ and $j+1$ in~\eqref{eq:logspace-spacing}
has mean $b_N / j$, immediate from~\eqref{eq:logspace-spacing}. The record-holder's gap is
therefore the widest, and the remaining competitors cluster
progressively more tightly below. This top-heavy spacing structure is
what permits the crossover of Section~\ref{sec:completion} to
separate the record-holder from the field: a modest percentage
advantage in score, arising from the natural statistics of the Gumbel
process, becomes an exponential advantage in completion time once it
is processed through the consecutive-success mechanism.

\begin{figure}[ht]
  \centering
  \includegraphics[width=\textwidth]{gap_distribution.pdf}
  \caption{Gap distribution between the record-holder and the
    second-ranked individual for log-normal scores.
    \textbf{(a)}~Probability density of the percentage gap
    $g = s_{(1)}/s_{(2)} - 1$. The median gap is $6.3\%$ (dashed);
    the Wadlow height gap of $3.0\%$ (dotted) is shown for comparison.
    \textbf{(b)}~Tail probability
    $P(\text{gap} > g) = (1+g)^{-1/\beta_N}$, with key
    thresholds marked. A gap exceeding $10\%$ occurs in $34\%$ of
    realizations.
    \textbf{(c)}~Comparison of the log-normal score distribution
    ($\sigma = 0.59$, $\beta_N = 0.0875$; solid) with normal
    heights ($\sigma_h/\mu = 0.04$, $\beta_h = 0.0058$; dashed).
    The log-normal produces gaps $15\times$ wider because
    $\beta_N \propto \sigma$. Wadlow's $3\%$ gap is anomalous
    for normal heights ($99.4$th percentile) but routine for
    log-normal scores ($29$th percentile).}
  \label{fig:gap-dist}
\end{figure}

\medskip
\noindent\textbf{Empirical illustration: extreme heights.}\;
As an independent check on the gap formula~\eqref{eq:gap-tail}, consider
the distribution of extreme human heights. The tallest verified person in
recorded history is Robert Wadlow (272~cm), followed by John Rogan
(264~cm), a gap of approximately $3.0\%$
\cite{guinness2024wadlow}. The next several individuals cluster tightly
below Rogan: 264, 257, 254, 254, 251, 251, 249, 249~cm. This
pattern of a single outlier standing apart from a tight cluster of
competitors is precisely the top-heavy spacing structure predicted by
the Poisson process.

To test this quantitatively, model heights as approximately
normal with $\mu_h = 175$~cm and $\sigma_h = 7$~cm. For a normal
distribution, the gap between the first and second order
statistics of $N$ draws is exponentially distributed with mean
$\sigma_h / \sqrt{2\ln N}$. At $N \approx 3 \times 10^{10}$
(the approximate population of the recorded era), this gives a
median gap of $0.7$~cm. Wadlow's $8$~cm gap exceeds the $99$th
percentile of this distribution, suggesting that the extreme tail
of human heights is heavier than Gaussian, consistent with the
known role of pituitary gigantism in producing the tallest
individuals. The important point for the present framework is that
when the underlying distribution \emph{is} log-normal, the
log-space dispersion~$\sigma$ is large enough (here $\sigma \approx
0.59$) that the Gumbel-process gaps are substantial: a median gap
of $6.3\%$ and a $34\%$ probability of exceeding $10\%$. The gap that
drives the crossover of Section~\ref{sec:completion} is not a
rare fluctuation but a routine feature of log-normal extremes.

\medskip
\noindent\textbf{Derivation note.}\;
Equations~\eqref{eq:logspace-spacing}--\eqref{eq:gap-tail} follow
from the Poisson-process representation of extreme order statistics
in the Gumbel domain of attraction
\cite{resnick1987,leadbetter1983}. The representation states that,
after normalizing by the Gumbel location and scale parameters, the
top order statistics of an i.i.d.\ sample converge in distribution
to the points of a unit-rate Poisson process on $(0,\infty)$, whose
inter-arrival times are i.i.d.\ $\operatorname{Exp}(1)$. The
exponential distribution of the gap $\Delta$ in~\eqref{eq:score-gap}
is a direct consequence. Monte Carlo simulations (generating the top
50 order statistics via the Poisson-process construction, assigning
each individual a random birth-time to define ``historical'' and
``contemporaneous'' subpopulations, and computing the score ratios)
confirm that the analytic formula matches the simulated gap
distribution to within sampling error at all tested percentiles.

\subsection{Remark on the Generational Interpretation}\label{ssec:generations}

Although the Gumbel-amplification mechanism renders generational population
inflation unnecessary, the generational perspective remains available as an
equivalent reformulation. If one prefers to attribute the rank-$\#1$
overshoot to the record-holder being drawn from $N_{\mathrm{eff}} = K N$
individuals across $K$~generational cohorts, the only modification is to
evaluate $z_{\max}$ at $N_{\mathrm{eff}}$ rather than~$N$
in~\eqref{eq:eq2}:
\begin{equation}\label{eq:sigma-gen}
  \sigma
  \;=\;
  \frac{\ln k}
       {\,z_{\max}(N_{\mathrm{eff}}) - z_q\,}\,,
  \qquad
  \mu
  \;=\;
  \ln m_B - \sigma\, z_q\,,
\end{equation}
with the quantile equation~\eqref{eq:eq1} unchanged (since the elite
subpopulation~$B$ is drawn from a single cohort). Because $z_{\max}$ grows
only as $\sqrt{2\ln N_{\mathrm{eff}}}$, even order-of-magnitude uncertainty
in~$K$ produces only modest changes in the inferred parameters.

\subsection{Remark on the Normality Approximation for $B$}
\label{ssec:normality}

As noted in Section~\ref{ssec:setup}, the left-truncated normal is not
Gaussian. Its density on the support $(\ln\tau,\infty)$ is
\[
  f_{B}(y)
  \;=\;
  \frac{1}{\sigma}\,
  \frac{\varphi\!\bigl((y-\mu)/\sigma\bigr)}
       {1 - \Phi\!\bigl((\ln\tau - \mu)/\sigma\bigr)}\,,
  \qquad y > \ln\tau\,,
\]
which has its mode at~$\ln\tau$ and is right-skewed. The degree of skewness
diminishes as~$p$ increases (less extreme selection), and the distribution
converges to the full Gaussian in the limit $p \to 1$. In empirical
settings the effective distribution of~$B$ may be closer to normal in
log-space than the strict truncation model predicts, for the reasons
discussed above (multidimensional selection, developmental noise, etc.).
The framework developed here therefore provides first-order estimates whose
accuracy improves as these softening mechanisms become more pronounced.



\section{Sequential Achievement and the Exponential-to-Linear Crossover}\label{sec:completion}

\subsection{Overview}\label{ssec:overview-s2}

Section~\ref{sec:inference} developed a framework for locating an
individual within a log-normally distributed population via a scalar
score~$s$. We now ask what this score \emph{implies} for an individual's
capacity to produce complex output. The central thesis of this section is
that many forms of significant achievement are well modelled as
\emph{sequential stochastic processes requiring consecutive successes}, and
that the score~$s$ enters these processes through a modified success
probability in a manner that induces a sharp \emph{crossover} in
expected completion time: from exponential in the sequence length for
low-scoring individuals to approximately linear for high-scoring ones.
This crossover provides a mechanistic explanation for the extreme
concentration of output among top performers: the observation that
individuals only modestly separated in latent ability can be separated by
orders of magnitude in realized achievement.

\subsection{Formal Definitions}\label{ssec:defs}

We introduce a small family of formal objects. The notation is informed by
type theory and is chosen to emphasize compositionality: complex
achievements are assembled from typed components with explicit parameters.

\begin{definition}[Achievement type]\label{def:achievement}
An \emph{achievement type} is a tuple
\[
  \mathfrak{a} \;=\; (\mathcal{K},\; p_0,\; \lambda)
\]
where:
\begin{itemize}
  \item $\mathcal{K}$ is a \emph{category} (a label identifying the domain
        of competence, e.g.\ a subdiscipline, a skill, a class of problem);
  \item $p_0 \in (0,1)$ is the \emph{objective success probability}---the
        baseline probability that a single attempt at this achievement
        succeeds for an individual of median ability in category~$\mathcal{K}$;
  \item $\lambda > 0$ is the \emph{occurrence rate}---the rate parameter of
        the exponential distribution governing the inter-arrival time of
        opportunities to attempt the achievement.
\end{itemize}
\end{definition}

\begin{definition}[Completion type]\label{def:completion}
A \emph{completion type} is a tuple
\[
  \mathfrak{c}
  \;=\;
  \bigl(\,\mathfrak{a}_1,\;\mathfrak{a}_2,\;\ldots,\;\mathfrak{a}_n\,\bigr)
\]
consisting of an ordered sequence of $n$ achievement types.
Completion requires that all $n$ achievements be realized
\emph{consecutively and in order}: the individual must succeed at
$\mathfrak{a}_1$, then at $\mathfrak{a}_2$, and so on through
$\mathfrak{a}_n$, with any single failure resetting the process to
$\mathfrak{a}_1$. The integer $n$ is called the \emph{depth} of the
completion.
\end{definition}

\begin{definition}[Individual score]\label{def:score}
Let $X$ be the trait value of an individual drawn from the log-normal
population of Section~\ref{sec:inference}, and let
$X_{\mathrm{ref}} = e^{\mu}$ denote the population median. The
\emph{score} of the individual relative to achievement category~$\mathcal{K}$
is
\begin{equation}\label{eq:score}
  s \;=\; \frac{X}{X_{\mathrm{ref}}} \;=\; e^{Y - \mu}\,,
\end{equation}
where $Y = \ln X \sim \mathcal{N}(\mu,\sigma^2)$. Thus $\ln s \sim
\mathcal{N}(0,\sigma^2)$, the median individual has $s = 1$, and the
score inherits the log-normal distribution of~$X$.
\end{definition}

In the homogeneous case, all achievement types in a completion share a
common category, success probability, and occurrence rate:
$\mathfrak{a}_i = (\mathcal{K}, p_0, \lambda)$ for all~$i$. We develop
the theory first for this case and generalize in
Section~\ref{ssec:hetero}.

\subsection{Score-Modified Success Probability}\label{ssec:modified-p}

The score $s$ from Section~\ref{sec:inference} modifies an individual's
effective success probability. A high-scoring individual's rich
processing across nearby steps in the achievement sequence produces, at
the level of each individual step, an emergent enhancement that is
well described by the following model: the step behaves \emph{as if}
the individual had $s$ independent attempts at it, succeeding whenever
at least one attempt succeeds. The probability that all $s$ putative
attempts fail is $(1-p_0)^s$, giving an effective probability of
\begin{equation}\label{eq:ptilde}
  \widetilde{p}(s)
  \;=\;
  1 - (1 - p_0)^{s}\,.
\end{equation}
The individual is not literally running $s$ parallel trials; the
enhancement arises from the interaction between consecutive steps, with
each step's resolution informed by the individual's processing of its
neighbors (see the remark on the independence framework below).
The formula~\eqref{eq:ptilde} captures the net probabilistic effect of
this interaction within a tractable independent-events framework.
The expression $(1-p_0)^s$ is well-defined for all real $s > 0$, not
only for integers, so~\eqref{eq:ptilde} extends smoothly from the
combinatorial interpretation at integer scores to the full
log-normally distributed population. For $s = 1$ (median individual),
$\widetilde{p}(1) = p_0$, recovering the baseline. The formulation
also permits the achievement types in the sequence to differ
(Section~\ref{ssec:hetero}), since each step carries its own
$p_{0,j}$ and the modified probability
$\widetilde{p}_j(s_j) = 1 - (1-p_{0,j})^{s_j}$ is computed per-step.

A few additional properties:
\begin{enumerate}
  \item[\textbf{(i)}]
    $\widetilde{p}(s) \to 1$ as $s \to \infty$, and
    $\widetilde{p}(s) \to 0$ as $s \to 0^+$.
  \item[\textbf{(ii)}]
    For small $p_0$ and moderate~$s$ (such that $s\,p_0 \ll 1$), a
    first-order expansion gives
    \begin{equation}\label{eq:ptilde-approx}
      \widetilde{p}(s) \;\approx\; s\,p_0\,,
    \end{equation}
    so the effective probability scales linearly with the score in this
    regime.
\end{enumerate}

\noindent\textbf{Phenomenological interpretation.}\;
The emergent enhancement captured by~\eqref{eq:ptilde} arises, at the
cognitive level, from the high-scoring individual's ability to process
nearby steps in the chain in relation to one another. Because the
individual's internal representation of the problem is richly connected,
each step is not encountered in isolation but is \emph{triangulated}
against the structure of its neighbors, cross-checked from multiple
directions before the individual commits. This inter-step reinforcement
is what the independent model represents, in compressed form, as $s$
parallel approaches to a single step.

The cognitive role of~$s$ is best understood not as a compression ratio
but as a measure of how completely it closes the gap between the
baseline probability~$p_0$ and certainty. At $s = 1$,
$\widetilde{p} = p_0$; at $s = 42$, $\widetilde{p} \approx 0.97$. When
$\widetilde{p}$ is near unity, consecutive steps become so reliable that
they merge phenomenologically: the individual perceives several steps as
a single confident move, producing the felt sensation of \emph{inferential
compaction}. The degree of compression varies with the difficulty of the
step. For a step with small~$p_0$ (where even the enhanced
$\widetilde{p}(s)$ is moderate), the individual experiences the step as
effortful and distinct; for a step with larger~$p_0$ (where
$\widetilde{p}(s)$ is near unity), the step dissolves into its neighbors.
The observed compaction ratio is therefore a downstream consequence of
high~$s$, not a direct reading of~$s$ itself.

To make this connection precise, define the \emph{felt compaction
factor}~$c$ as the number of consecutive steps that merge into a single
perceived move. In the idealized case where $c$ is an integer that
divides~$n$, the individual would need $n/c$ consecutive successes each
at probability~$p_0$, giving a conjunction probability of $p_0^{n/c}$.
The per-step probability required to reproduce this conjunction in
$n$~individual trials is $\widetilde{p} = p_0^{1/c}$. Equating this
with~\eqref{eq:ptilde} yields the relationship
\begin{equation}\label{eq:c-from-s}
  1 - (1-p_0)^s \;=\; p_0^{1/c}\,,
\end{equation}
which implicitly defines $c$ as a function of $s$ and $p_0$. We work in
terms of~$s$ rather than~$c$ for three reasons. First, $s$~is directly
determined by the log-normal distribution of
Section~\ref{sec:inference}, whereas $c$ is not. Second, the
$s$-based formulation passes the empirical test cases of
Section~\ref{sec:results} (the crossover, the linearity at high
scores, and the unremarkability of Einstein at $n = 1$), whereas a
formulation in which $c$ is treated as the primary parameter does not
(see Section~\ref{ssec:phase}). Third, and most instructively, for
\emph{fixed}~$s$ the implied $c$ varies with~$p_0$: harder steps
(smaller~$p_0$) yield a smaller~$c$, and easier steps (larger~$p_0$)
yield a larger one. This means the felt compression is not a fixed
property of the individual but an interaction between the individual's
score and the difficulty of each step, exactly as one would expect
phenomenologically.

The critical observation for linking the two pictures is that
$\widetilde{p}(s)$ is monotonically increasing in both $s$ and $c$:
higher score always implies higher per-step probability, which always
implies greater felt compression. The score therefore functions as a
\emph{sufficient proxy} for the phenomenological compaction, without
requiring~$c$ to be specified directly.

As a concrete illustration, consider Einstein's 1905 derivation of
mass-energy equivalence \cite{einstein1905inertia}. The published paper
contains on the order of five distinct inferential steps. The model
predicts that Einstein's internal experience of the derivation compressed
these into approximately two moves, each carrying the others as implicit
connotations: (1)~emit the two light beams, with the connotation of
antisymmetric (odd-function) structure\footnote{The difference in effect here is substantial. While a classical step-by-step version essentially has to lead with the light beams and (with luck) reach the antisymmetric cancellation by surprise, in the case of Einstein he is able to lead directly with the realization of that antisymmetric structure straight from the equation, essentially cutting his probability of missing the idea in half by means of leading with the punch and simply establishing the light beams as a “manifestation” instance of the structure. This allows him to systematically reach deeper into the derivation with each next step than would a normal theorist, because initiation for him looks like conclusions for them.}; (2)~review the resulting changes,
with the connotation of the invariance-variance divide that yields the
mass-energy relation. What the reader parses as five steps, the author
traversed as two richly connoted ones, a felt compaction of roughly
$5/2 = 2.5$ for this particular derivation. The compaction ratio need
not be constant: it depends on the per-step $p_0$ values in the
achievement sequence, with harder steps (lower~$p_0$) resisting
compression and easier steps (higher~$p_0$) dissolving more readily
into their neighbors. This variability is naturally accommodated by the
heterogeneous model of Section~\ref{ssec:hetero}.

\medskip
\noindent\textbf{Remark on the independence framework.}\;
The actual cognitive process underlying high-scoring performance is
almost certainly not one of independent trials. Approaches to nearby
steps \emph{reinforce} one another: a promising line of attack at
step~$j$ feeds structural information into the approaches available at
step~$j+1$, and partial progress along one thread constrains and informs
the others. In a graph-theoretic picture, the steps can talk to each
other, and the resulting dependencies alter the joint probabilities
throughout the chain. Modeling these correlations directly would require
specifying the full dependency structure, which we deliberately avoid
(\emph{hypotheses non fingo}).
Instead, the formula~\eqref{eq:ptilde} should be understood as the
\emph{correction} that reproduces the net effect of the correlated
process from within the independent-events framework. The score~$s$ is
the weight that accomplishes this translation: it inflates the per-step
success probability by exactly the amount needed to make the independent
model output the same completion-time distribution as the true correlated
one. In this reading, $\widetilde{p}(s) = 1 - (1-p_0)^s$ is not a
description of the cognitive mechanism but a \emph{sufficient encoding}
of its probabilistic consequences, analogous to a sufficient statistic
that compresses a complex likelihood into a lower-dimensional summary
without loss of inferential content.

\medskip
\noindent\textbf{Remark on rate-score trade-offs and achievement-class variation.}\;
The inferential compaction picture suggests that a high-scoring
individual's cognitive style may differ qualitatively from that of a
lower-scoring one: deeper, more deliberate rumination over
conceptual structure, at the expense of raw throughput on routine
steps. Within the model, this would manifest as a \emph{lower}
occurrence rate~$\lambda$ for certain achievement classes, partially
offsetting the higher~$\widetilde{p}(s)$. The high scorer still
dominates at sufficient depth, since the exponential dependence on
$\widetilde{p}$ in~\eqref{eq:expected-trials} eventually overwhelms
any linear penalty from reduced~$\lambda$, but at low~$n$ the
advantage may be modest or even imperceptible. This is consistent
with the empirical observation that extreme outliers do not necessarily
distinguish themselves on shallow, routine tasks.

Separately, not all achievement types need share the same baseline
probability. A routine technical derivation step may have a relatively
large~$p_0$, placing its single-step reliability point $1/a \approx 1/p_0$, and with it the crossover $s_c(n) = \ln n/a$ at any given depth, within reach
of many competent practitioners. A profound conceptual reorientation,
however, such as Einstein's recognition of the equivalence principle
\cite{pais1982subtle}, may correspond to an achievement type with very
small~$p_0$, reflecting the difficulty of the ontological move and
pushing its crossover far into the tail. For an individual well above the
crossover score of even the hardest step, all classes in the chain
are tractable; for those near or below the crossover of the low-$p_0$ classes,
these steps become the binding constraint, and the completion fails
there regardless of how easily the routine steps are handled.
The heterogeneous model of Section~\ref{ssec:hetero}
accommodates this variation naturally through the per-step
parameters $(\mathcal{K}_j, p_{0,j}, \lambda_j)$.

\subsection{A Realist Interpretation: Divergence Resistance}\label{ssec:divergence}

The score-modified probability~\eqref{eq:ptilde} provides a sufficient statistical description of the consecutive-success process, but it leaves open the question of what the score \emph{is} at the cognitive level. This subsection proposes a realist interpretation that inverts the usual framing of the genius phenomenon.

\subsubsection{The Standard Question and Its Inversion}

The standard question about extreme outliers is: \emph{what special capacity does the genius possess that others lack?} This framing leads naturally to mystification, since any proposed answer (extraordinary memory, superhuman pattern recognition, divine inspiration) attributes to the genius a positive quality that must then be explained in its own right. We propose to invert the question: \emph{what does the typical individual possess that the genius lacks?}

The answer we suggest is \emph{divergence}---a spontaneous tendency, at each step of an inferential chain, to extrapolate the structure of the problem incorrectly. In this framing, the genius is not doing something extraordinary; rather, everyone else is doing something detrimental. The baseline cognitive process, uncorrupted by divergence, would produce the correct next step of any sufficiently constrained inferential chain. What prevents most individuals from sustaining this process over many consecutive steps is the accumulation of small deflections---incorrect extrapolations, paradigmatic biases, premature commitments to suboptimal framings---that compound across the chain and eventually produce a fatal error, resetting the process.

\subsubsection{Empirical Motivation}

The divergence interpretation is motivated by a specific empirical observation: when individuals engage with a novel theoretical framework for the first time, they tend to extrapolate its next structural implications incorrectly, and they do so with striking regularity. The errors are not random in a uniform sense; they are \emph{systematically biased} toward familiar structures. A probabilist will extrapolate toward a probability framework. A graph theorist will extrapolate toward a familiar graph topology. A physicist trained in Copenhagen quantum mechanics will extrapolate toward Copenhagen-compatible structures. In each case, the individual's prior cognitive commitments act as an attractor basin, pulling their inference away from the actual structure and toward a more familiar one. This attractor dynamic has a named experimental literature: the Einstellung effect \cite{luchins1942}, in which a practiced solution schema blocks discovery of a better available one, with eye-tracking evidence in chess experts showing the familiar idea continuing to direct attention even while the subject reports searching for alternatives \cite{bilalic2008}. Divergence, as the term is used here, is the chained form of that effect: not one blocked solution but a per-step deflection probability compounding across an inferential sequence.

The key observation is that this divergence is \emph{the norm, not the exception}. What requires explanation is not the rare individual who avoids it, but the ubiquitous tendency to succumb to it. From this perspective, the per-step failure probability $1 - \widetilde{p}(s)$ in the consecutive-success process is not a measure of ``lacking genius'' but a measure of \emph{carrying divergence friction}: the probability that, at a given step, the individual's cognitive apparatus deflects them away from the correct extrapolation and toward a familiar but incorrect one.

\subsubsection{The Score as Divergence Resistance}

Under this interpretation, the score~$s$ acquires a concrete cognitive meaning: it is a measure of \emph{divergence resistance}, the degree to which an individual's inferential process is free from the systematic biases that deflect typical individuals away from the correct structural extrapolation at each step. A high score does not indicate the presence of a special faculty; it indicates the \emph{absence} of a common impediment. The genius is not superhuman; the genius is \emph{unimpeded}.

This inversion has several consequences for the interpretation of the model.

First, it explains the phenomenology of high-level creative work as reported by practitioners. When mathematicians and theoretical physicists describe their deepest insights as ``obvious in retrospect,'' or as arriving through ``inner knowing'' rather than sequential construction, this is precisely what the absence of divergence would feel like from the inside. The correct next step simply \emph{appears}, because nothing is deflecting the thinker away from it. The experience of revelation is the experience of frictionless inference. The experience of struggle, conversely, is the experience of friction---of one's cognitive apparatus pulling toward incorrect extrapolations that must be effortfully corrected.

Second, it explains why the score~$s$ can be treated as a stable, 
individual-level parameter across instances of the same achievement 
class. If divergence resistance within a given domain reflects a 
general property of an individual's engagement with that domain's 
inferential structures---the degree to which their apparatus is free 
from paradigmatic rigidity, premature commitment, or systematic bias 
\emph{specific to that class of problems}---then an individual who 
resists divergence at one step of a chain is likely to resist it at 
subsequent steps as well, and across different chains of the same 
categorical type. This is what licenses the assignment of a single 
scalar~$s$ per achievement class: it encodes the empirical 
extrapolation that an individual's divergence rate on the early legs 
of a journey is predictive of their divergence rate on later legs, 
\emph{within that class}. Note, however, that this stability does 
not extend to comparisons \emph{across} achievement classes. An 
individual's divergence resistance in theoretical physics need not 
translate to divergence resistance in, say, musical composition or 
literary criticism, since the cognitive commitments and structural 
biases that generate divergence are class-specific. The heterogeneous 
model of Section~\ref{ssec:hetero}, which assigns per-step parameters 
$(\mathcal{K}_j, p_{0,j}, \lambda_j)$ to each achievement type, 
accommodates this naturally: the score~$s_j$ is category-specific, 
and cross-class comparisons of scores are not well-defined within the 
framework (see the caveat in Section~\ref{ssec:physics}).

Third, and most importantly, the inversion redirects the explanatory burden from the rare to the common. Instead of asking ``what makes geniuses special?'' (a question that invites mystification), we ask ``what makes typical individuals diverge?'' (a question that invites tractable cognitive science). The answer plausibly involves some combination of: paradigmatic training that creates strong prior commitments to specific structural forms; cognitive biases toward familiar patterns; noise in the inferential process that accumulates across steps; and the difficulty of maintaining a neutral, uncommitted stance toward the structure of a problem while simultaneously engaging with it deeply. Each of these is a known cognitive phenomenon, and their joint effect---a per-step divergence probability that, while small, compounds across a long chain to produce near-certain failure---is precisely the mechanism the model formalizes.

\subsubsection{Relation to the Statistical Model}

The divergence interpretation does not replace the statistical model of the preceding sections; it provides a \emph{realist substrate} for it. The per-step probability $\widetilde{p}(s) = 1 - (1-p_0)^s$ remains the operative quantity in all calculations. What the divergence interpretation adds is a cognitive-level account of why the score~$s$ should be expected to behave as a stable individual parameter: it reflects a general property of the individual's inferential process (freedom from divergence friction) rather than a task-specific ability. This is what licenses the model's central assumption that a single scalar can govern an individual's performance across a long and heterogeneous chain of inferential steps.

The interpretation also sharpens the model's relationship to Bessis's \cite{bessis2024mathematica, bessis2025genes, bessis2026jaimungal} account of mathematical intuition. Bessis argues that the central product of mathematics is the development of intuition, and that formal proofs are byproducts of intuitive insight rather than the substance of mathematical work. Under the divergence interpretation, ``developing intuition'' is precisely the process of reducing one's divergence rate: training the cognitive apparatus to extrapolate structures correctly, so that the correct next step of an inferential chain is perceived directly rather than arrived at through effortful correction. The score~$s$ is, in this reading, a summary measure of how far this training has progressed.

\subsection{The Consecutive-Success Process}\label{ssec:process}

We now define the stochastic process governing completion. Fix a
completion type $\mathfrak{c} = (\mathfrak{a}_1, \ldots, \mathfrak{a}_n)$
and an individual with score~$s$. The process evolves as follows:
\begin{enumerate}
  \item
    The individual begins at stage $j = 1$ (no progress).
  \item
    At stage~$j$, an opportunity for achievement~$\mathfrak{a}_j$ arrives
    after an exponentially distributed waiting time with rate~$\lambda_j$.
  \item
    Upon arrival, a Bernoulli trial with success probability
    $\widetilde{p}_j = 1 - (1 - p_{0,j})^{s_j}$ is conducted, where
    $s_j$ is the individual's score in category~$\mathcal{K}_j$.
  \item
    If the trial succeeds, the individual advances to stage $j+1$.
  \item
    If the trial fails, the individual returns to stage $j = 1$
    (full restart).
  \item
    Completion occurs upon advancing past stage~$n$.
\end{enumerate}

\subsection{Expected Completion Time: Homogeneous Case}\label{ssec:time}

In the homogeneous case ($\widetilde{p}_j = \widetilde{p}$ and
$\lambda_j = \lambda$ for all~$j$), the expected number of trials to achieve $n$~consecutive successes is a classical result. The process is a first-passage problem under restart, whose general theory is developed in \cite{pal2017firstpassage}; what this paper contributes is the application to sequential achievement, the scaling analysis and its interpretation, not the restart mathematics. Let $M_n$ denote
this random variable. Defining $E_k$ as the expected additional trials
when $k$ consecutive successes have been accumulated, we have the
recurrence
\begin{equation}\label{eq:recurrence}
  E_k = 1 + \widetilde{p}\,E_{k+1} + (1-\widetilde{p})\,E_0\,,
  \qquad k = 0, 1, \ldots, n-1,
  \qquad E_n = 0.
\end{equation}
The solution is \cite{feller1968}
\begin{equation}\label{eq:expected-trials}
  \boxed{\;
    \mathbb{E}[M_n]
    \;=\;
    \sum_{j=1}^{n} \widetilde{p}^{\,-j}
    \;=\;
    \frac{\widetilde{p}^{-n} - 1}{1 - \widetilde{p}}
  \;}
\end{equation}

Since each trial is preceded by an exponential waiting time with
rate~$\lambda$, the expected wall-clock time to completion is
\begin{equation}\label{eq:expected-time}
  \mathbb{E}[T_n]
  \;=\;
  \frac{1}{\lambda}\,\mathbb{E}[M_n]
  \;=\;
  \frac{1}{\lambda}\,
  \frac{\widetilde{p}(s)^{-n} - 1}{1 - \widetilde{p}(s)}\,.
\end{equation}

\subsection{The Sharp Crossover}\label{ssec:phase}

Equation~\eqref{eq:expected-time} exhibits a sharp transition in its
dependence on the completion depth~$n$, governed by the score~$s$.

\subsubsection{Low-Score Regime ($\widetilde{p} \ll 1$)}

When $s$ is small (near or below the median), $\widetilde{p}(s) \approx
s\,p_0 \ll 1$ by \eqref{eq:ptilde-approx}. The dominant term in
\eqref{eq:expected-trials} is $\widetilde{p}^{-n}$, giving
\begin{equation}\label{eq:exp-regime}
  \mathbb{E}[T_n]
  \;\approx\;
  \frac{1}{\lambda}\,(s\,p_0)^{-n}
  \;=\;
  \frac{1}{\lambda}\,s^{-n}\,p_0^{-n}\,.
\end{equation}
The expected time grows \emph{exponentially} in the depth~$n$. For even
moderate~$n$, completion is effectively impossible within a human
lifetime.

\subsubsection{High-Score Regime ($\widetilde{p} \to 1$)}

When $s$ is large, $\widetilde{p}(s) = 1 - (1-p_0)^s \to 1$. Writing
$\widetilde{p} = 1 - \epsilon$ with $\epsilon = (1-p_0)^s \ll 1$, a
Taylor expansion gives
\begin{equation}\label{eq:lin-regime}
  \mathbb{E}[M_n]
  \;=\;
  \frac{(1-\epsilon)^{-n} - 1}{\epsilon}
  \;\approx\;
  \frac{e^{n\epsilon} - 1}{\epsilon}
  \;\approx\;
  n \qquad \text{when } n\epsilon \ll 1\,.
\end{equation}
The expected number of trials is approximately~$n$, \emph{linear} in
the depth, because the individual almost never fails and therefore almost
never restarts---but the condition is $n\epsilon \ll 1$ and not merely $\epsilon \ll 1$, since $(1-\epsilon)^{-n} \approx e^{n\epsilon}$ and the depth multiplies the failure rate; that factor of $n$ makes the crossover depth-dependent, and the next subsubsection carries it. The expected time is simply
\begin{equation}\label{eq:time-linear}
  \mathbb{E}[T_n] \;\approx\; \frac{n}{\lambda}\,.
\end{equation}

\subsubsection{The Crossover Score, and the Finite-Size Scaling Function}
The transition between regimes is not at a fixed score but at a score that grows with the depth. Writing $a = -\ln(1-p_0)$, so that $\epsilon = (1-p_0)^s = e^{-as}$, the exponential regime is $n\epsilon \gg 1$ and the linear regime $n\epsilon \ll 1$; the crossover is therefore at $n e^{-a s} \sim 1$, that is at
\begin{equation}\label{eq:crossover}
s_c(n) \;\sim\; \frac{\ln n}{a} \;=\; \frac{\ln n}{-\ln(1-p_0)} \;\approx\; \frac{\ln n}{p_0} \qquad\text{for small } p_0.
\end{equation}
A crossover defined by $\widetilde{p}(s)$ alone being bounded away from $0$ and $1$---by $(1-p_0)^{s} = 1/e$, say, giving $s = 1/a \approx 1/p_0$---locates the point at which a \emph{single} step becomes reliable, and that is the wrong question for an $n$-step task: with $p_0 = 0.1455$ and $n = 7700$, the values of Section~\ref{sec:results}, the single-step point sits at $s \approx 6.3$ while the crossover \eqref{eq:crossover} sits at $s_c \approx 56.7$, and the $\ln n$ factor is the whole difference. The crossover has an exact scaling form. Put
\begin{equation}\label{eq:scaling-variable}
s \;=\; \frac{\ln n + c}{a}, \qquad\text{so that}\qquad n e^{-a s} \;=\; e^{-c},
\end{equation}
and then, as $n \to \infty$ at fixed $c$,
\begin{equation}\label{eq:scaling-function}
\frac{\mathbb{E}[T_n]}{n/\lambda} \;\longrightarrow\; F(e^{-c}), \qquad F(x) \;:=\; \frac{e^{x} - 1}{x},
\end{equation}
since $\mathbb{E}[M_n]/n = ((1-\epsilon)^{-n} - 1)/(n\epsilon) \to (e^{x} - 1)/x$ with $x = n\epsilon$. Three regimes follow: for $s \ll s_c(n)$, $x \to \infty$ and $\mathbb{E}[T_n]$ is exponential in $n$ at fixed $s$, and stretched-exponential along $s = \alpha\ln n/a$ with $0 < \alpha < 1$, where $\log\big(\mathbb{E}[T_n]/(n/\lambda)\big) \sim n^{1-\alpha}$; for $s = s_c(n) + O(1/a)$, the overhead over the ideal $n/\lambda$ is the finite factor $F(e^{-c})$; for $s \gg s_c(n)$, $x \to 0$, $F \to 1$, and $\mathbb{E}[T_n] \sim n/\lambda$---linear over the finite range of depths $n \ll e^{as}$, since at any fixed score $\widetilde{p}(s) < 1$ and the growth is eventually exponential. The transition is a sharp finite-size crossover, and it is called that throughout. The scaling function $F$ is smooth in $c$, so \eqref{eq:scaling-function} exhibits no nonanalyticity and does not by itself license the phrase ``phase transition''; what does is the coarser scaled variable $\alpha := as/\ln n$, taken with $\alpha > 0$ since the per-step failure probability $e^{-as} = n^{-\alpha}$ must be a probability, in which $n e^{-as} = n^{1-\alpha}$, so that as $n \to \infty$ the overhead $\mathbb{E}[T_n]/(n/\lambda) = \big((1 - n^{-\alpha})^{-n} - 1\big)/n^{1-\alpha}$ tends to $1$ for $\alpha > 1$ and to $+\infty$ for $0 < \alpha < 1$. It is asymptotic to $F(n^{1-\alpha})$ in ratio only for $\alpha > \tfrac12$, since $(1 - n^{-\alpha})^{-n} = \exp\big(n^{1-\alpha} + \tfrac12 n^{1-2\alpha} + \cdots\big)$; at $\alpha = \tfrac12$ the ratio tends to $e^{1/2}$, and for $0 < \alpha < \tfrac12$ the two agree only in their leading logarithm, $\ln$ of the overhead $\sim n^{1-\alpha}$. The step is the same either way---a limiting step at $\alpha = 1$, the finite-size scaling \eqref{eq:scaling-function} resolving the step's interior at the finer scale $c$. The word is used only in that limiting sense, and the title does not use it. The calibrated case of Section~\ref{sec:results} sits inside the crossover and not beyond it: at $s = 47.7$, $x = n e^{-as} \approx 4.26$ and $F(4.26) \approx 16.3$, which is the roughly fifteen-fold overhead---sixteen task-depths' worth of evaluations---that the tables there exhibit, so the tables were computing the correct object and the named threshold was not.

The right quantity for ``effectively unable to complete'' is not the expected time but the probability of completion within a horizon $H$. Everyone with $p_e(s) > 0$ completes eventually with probability one, so no score makes completion impossible; what a low score makes is completion within a specified finite budget improbable, and the finite-horizon statement below is exact. Between full restarts the process is a sequence of independent runs from stage~$1$, each succeeding outright with probability $p_e^{\,n}$, so to leading order in $p_e^{\,n}$ the completion time is geometric in runs and
\begin{equation}\label{eq:horizon}
P(T_n \le H) \;\approx\; 1 - \exp\!\big(-H/\mathbb{E}[T_n]\big) \;\le\; \frac{H}{\mathbb{E}[T_n]},
\end{equation}
the inequality holding without approximation---not by Markov's inequality, which bounds the opposite tail, but because $T_n$ is new better than used in expectation: from any state of partial progress the expected remaining time is at most the expected time from a fresh start, and for such a time with mean $\mu$ the survival function dominates the equilibrium-excess survival, whence $P(T_n > H) \geq 1 - H/\mu$. The exact survival function is also available: with $Q_s$ the transient generator of the chain on progress states $0, \ldots, n-1$ and $\mathbf{e}_0$ the initial state, $P(T_n > H) = \mathbf{e}_0^{\top} e^{Q_s H}\mathbf{1}$, and for independent workers $P(\text{some completion by } H) = 1 - \prod_i P(T^{(i)}_n > H)$. For the median-quantile individual of Section~\ref{sec:results}, with $\mathbb{E}[T_n]$ of order $10^{6440}$ years and any human or cosmological $H$, the right side is zero to every decimal place that could ever be written; that is the barrier's content, and ``impossibility'' is replaced throughout by this quantity.

\subsubsection{A Continuous Generalization: Fractional Depth through the Ideal Duration}\label{sssec:continuous}
% REFS: cites pal2017firstpassage, whose entry is embedded in the preamble (filecontents, howlong_restart.bib).
The depth $n$ enters \eqref{eq:expected-time} only through the power $\widetilde{p}^{\,-n}$, so the formula itself does not need $n$ to be an integer. Write
\begin{equation}\label{eq:ideal-duration}
\tau \;:=\; \frac{n}{\lambda}
\end{equation}
for the expected completion time when no step ever fails, $\widetilde{p} = 1$, where the time to $n$ successes is Erlang with mean $n/\lambda$: the duration the task would take if every opportunity succeeded, its \emph{ideal duration}. Substituting $n = \lambda\tau$ into \eqref{eq:expected-time} gives
\begin{equation}\label{eq:continuous}
\boxed{\;\mathcal{T}(\tau) \;:=\; \frac{\widetilde{p}^{\,-\lambda\tau} - 1}{\lambda\,(1 - \widetilde{p})}\;}
\end{equation}
which is exact for the model whenever $\lambda\tau$ is an integer and is a smooth function of $\tau > 0$ for every real $\tau$. That is the continuous generalization: the exact formula, read at fractional depth. As $\widetilde{p} \to 1$ it reduces to the ideal duration itself, $\mathcal{T}(\tau) \to \tau$, and with $\epsilon = 1 - \widetilde{p}$,
\begin{equation}\label{eq:continuous-expansion}
\mathcal{T}(\tau) \;=\; \tau\left[1 + \frac{\lambda\tau + 1}{2}\,\epsilon + O(\epsilon^{2})\right]
\end{equation}
at fixed $\lambda$ and $\tau$, the correction growing with the depth $\lambda\tau$, as the condition $n\epsilon \ll 1$ of \eqref{eq:lin-regime} requires. The reduction concerns the mean: at $\widetilde{p} = 1$ and integer depth the completion time is still Erlang, with variance $\tau/\lambda$, (gamma at noninteger depth in the realization recorded below, with the same variance), so $\tau$ is what the task takes on average and not on every occasion.

Two clarifications fix what has and has not been generalized. Time was already continuous: opportunities arrive after exponential waits, and what was discrete is progress, the count of steps. The generalization extends that count to fractional values. It changes nothing in the process of Section~\ref{ssec:process} or in any result derived from it; it is an aside on the formula, not a revision of the model. And a formula for the mean does not by itself define a stochastic process at noninteger depth. A realization does exist, and one is recorded below; it is a modeling choice rather than something the integer formula forces, and what remains open is which realization suits a given application, and so what a fractional step \emph{is} there.\footnote{Thin the opportunity process into successes at rate $\lambda\widetilde{p}$ and failures at rate $r = \lambda(1 - \widetilde{p})$. Without failures the time to $n$ successes is Gamma with shape $n$ and rate $\lambda\widetilde{p}$. Extend it to $S_\nu \sim \mathrm{Gamma}(\nu, \lambda\widetilde{p})$ with $\nu = \lambda\tau > 0$, and restart with a fresh draw at each failure. Since $\mathbb{E}\big[e^{-rS_\nu}\big] = \big(\lambda\widetilde{p}/(\lambda\widetilde{p} + r)\big)^{\nu} = \widetilde{p}^{\,\nu}$, the renewal calculation for restart at rate $r$ gives a mean completion time of $(1 - \widetilde{p}^{\,\nu})/(r\,\widetilde{p}^{\,\nu}) = \mathcal{T}(\tau)$ exactly, for every real $\nu > 0$. Its mean uninterrupted duration is $\tau/\widetilde{p}$, which tends to $\tau$ as $\widetilde{p} \to 1$.}

The extension makes the question of a time scale precise. In the linear regime the formula carries none of its own: $\mathcal{T}(\tau) = \tau$, the duration of the task and nothing else. The depth, read back from the ideal duration, is $n = \lambda\tau$, the ideal duration counted in units of the mean opportunity interval $1/\lambda$; that interval is a mean and not a minimum, exponential waits being arbitrarily short, so $1/\lambda$ is not a quantum of time but the unit in which the depth counts. Away from $\widetilde{p} = 1$ the formula introduces two further rates, $r := \lambda(1 - \widetilde{p})$ and $\kappa := -\lambda\ln\widetilde{p}$, and in terms of them
\begin{equation}\label{eq:continuous-rates}
\mathcal{T}(\tau) \;=\; \frac{e^{\kappa\tau} - 1}{r}, \qquad \frac{\mathcal{T}(\tau)}{\tau} \;=\; \frac{\kappa}{r}\,F(\kappa\tau),
\end{equation}
with $F$ the scaling function of \eqref{eq:scaling-function}. Here $1/r$ is the mean interval between failures and $1/\kappa$ the stretch of ideal duration over which the exponential factor grows $e$-fold; the two coincide as $\widetilde{p} \to 1$, where $\kappa/r \to 1$ and the overhead is $F(\kappa\tau)$, as in \eqref{eq:scaling-function}. The crossover between the regimes is then the comparison of the ideal duration $\tau$ with this scale. For $\tau \ll 1/\kappa$, $\mathcal{T}(\tau) = (\kappa/r)\,\tau + O(\tau^{2})$, linear in the duration, and for $\tau \gg 1/\kappa$ the growth is exponential. Linear is not the same as close to the ideal duration: $\mathcal{T}(\tau) \approx \tau$ needs $\kappa/r \approx 1$ as well, which holds near $\widetilde{p} = 1$. At $\widetilde{p} = 0.1$, $\kappa/r = -\ln(0.1)/0.9 \approx 2.56$, and the curve is locally linear with slope about $2.56$. The integer model's near-linear regime, negligible overhead over $n/\lambda$, is the case in which both hold. A time scale is thus introduced not by the depth but by failure. The linear regime is scale-free in $\tau$, and restart supplies the characteristic duration against which the length of a task is measured. What that scale means, whether it belongs to the worker, to the task, or to their encounter, is a question the formula poses rather than answers.

\noindent\textit{An aside: the limit of fine steps.} A different construction, not to be confused with the extension above, keeps the steps integer but makes them small. Take $n = \lambda\tau$ and $\widetilde{p} = 1 - r/\lambda$ with $\tau$ and $r$ fixed, and let $\lambda \to \infty$ along values for which $n$ is an integer: each step becomes a vanishing increment of progress carrying a vanishing failure probability, with the failure rate per unit of ideal duration held at $r$. Then
\begin{equation}\label{eq:fine-step-limit}
\mathbb{E}[T_n] \;\longrightarrow\; \frac{e^{r\tau} - 1}{r} \;=\; \tau\,F(r\tau),
\end{equation}
the scaling function now exact rather than asymptotic. The limit is a process in its own right: work advances at unit speed, completion needs $\tau$ units of uninterrupted work, failures arrive at Poisson rate $r$, and each failure resets progress to zero. Its mean follows from one renewal equation. With $R \sim \mathrm{Exp}(r)$, an attempt lasts $\min(R, \tau)$ and is repeated if $R < \tau$, so the mean $m$ satisfies $m = \mathbb{E}[\min(R, \tau)] + P(R < \tau)\,m$, and $\mathbb{E}[\min(R, \tau)] = (1 - e^{-r\tau})/r$ gives $m = (e^{r\tau} - 1)/r$. This is first passage under restart, an established setting \cite{pal2017firstpassage}: for $r\tau \ll 1$, $m = \tau + r\tau^{2}/2 + O(r^{2}\tau^{3})$, and for $r\tau \gg 1$, $m \sim e^{r\tau}/r$. In this limit the opportunity scale $1/\lambda$ disappears and only the restart scale $1/r$ remains, so the crossover compares the ideal duration of the task with the mean time between resets. The limit is joint: letting $\lambda \to \infty$ with $\widetilde{p} < 1$ held fixed instead sends the failure rate, and with it the mean completion time, to infinity. Like the extension, the limit leaves the model and the results of this section unchanged.

\noindent\textit{Provenance of the functional form.} The function in \eqref{eq:fine-step-limit} has an earlier source in the author's working notes. A journal entry dated September 30, 2024, proposes, by an interpolation argument, the function $f_k(x) = (e^{kx} - 1)/k$, chosen for its limit $f_k(x) \to x$ as $k \to 0$, and offers its reciprocal as a completion rate. The continuum model of the aside recovers the same function with explicit meanings for its parameters: $x \leftrightarrow \tau$, the uninterrupted work required, and $k \leftrightarrow r$, the rate at which resets destroy accumulated progress. The match is to the continuum limit; at finite opportunity rate the formula is $(e^{\kappa\tau} - 1)/r$, with $\kappa \neq r$ in general. The entry anticipated the functional form. It did not contain the restart derivation, which the present construction supplies, and the correspondence concerns the form alone, independently of the entry's numerical calibration.\footnote{The same entry integrates the proposed rate over $k$. For $x > 0$, $\int_0^\infty k\,(e^{kx} - 1)^{-1}\,dk = x^{-2}\int_0^\infty u\,(e^{u} - 1)^{-1}\,du = \pi^{2}/(6x^{2})$ exactly, since $(e^{u} - 1)^{-1} = \sum_{m \ge 1} e^{-mu}$ and $\int_0^\infty u\,e^{-mu}\,du = 1/m^{2}$, so the appearance of $\pi^{2}/6$ is the reciprocal-square series built into the integral, not a numerical coincidence. Three qualifications attach. The constant needs the infinite upper limit; a finite cutoff leaves a tail term. Reading the integral as a population's output requires a population measure, and an unweighted integral over $k$ fixes one implicitly, a constant density in that coordinate, without deriving the distribution of people. And summing reciprocal mean completion times gives a long-run aggregate throughput under repeated-task assumptions; its reciprocal is not in general the mean time until the first completion, for which the survival function behind \eqref{eq:horizon} is the right object: two workers who each finish deterministically after $\tau$ first complete at $\tau$, while the reciprocal of their summed rates is $\tau/2$.}

\subsubsection{The Amplification Ratio}

For fixed depth~$n$, the ratio of expected completion times between a
median individual ($s = 1$) and an individual with score~$s$ is, in the
regime $s\,p_0 \ll 1$,
\begin{equation}\label{eq:amp-ratio}
  \frac{\mathbb{E}[T_n \mid s=1]}{\mathbb{E}[T_n \mid s]}
  \;\approx\;
  s^{\,n}\,.
\end{equation}
This is \emph{polynomial in~$s$ of degree~$n$}. For a completion of
depth $n = 10$ and an individual at score $s = 5$, the speed-up is
$5^{10} \approx 10^7$: a task requiring ten million effective lifetimes for
a median individual is completed by the high-scorer within one. The
polynomial exponent~$n$ is precisely the depth of the achievement
sequence, connecting the \emph{complexity} of the task to the \emph{sensitivity} of completion time to ability. The approximation $\widetilde{p}(s) \approx s\,p_0$ behind \eqref{eq:amp-ratio} is raised to the $n$th power, so its one-step error compounds: since $\widetilde{p}(s) < s\,\widetilde{p}(1)$ for $s > 1$, $s^{n}$ overstates the exact ratio $\big[(\widetilde{p}(1)^{-n} - 1)/(1 - \widetilde{p}(1))\big] \big/ \big[(\widetilde{p}(s)^{-n} - 1)/(1 - \widetilde{p}(s))\big]$, by a factor that grows with $n$. For the example, the exact ratio is about $7.7 \times 10^{6}$ at $p_0 = 0.01$, against $5^{10} \approx 9.8 \times 10^{6}$, and about $2.9 \times 10^{5}$ at the calibrated $p_0 = 0.1455$ of Section~\ref{sec:results}, where $s\,p_0 \approx 0.73$ lies outside the regime $s\,p_0 \ll 1$ altogether. The degree-$n$ sensitivity is the point; the figure $5^{10}$ is an order of magnitude only under small $p_0$.

\subsection{What Carries the Amplification: Two Sensitivities}\label{ssec:sensitivity}
Two modelling choices carry most of the spectacular amplification, and neither is derived from the log-normal distribution or from the sequential process; both are hypotheses, and the model's robustness is a question about them before it is a question about $k$. The first is the link \eqref{eq:ptilde}, $\widetilde{p}(s) = 1 - (1-p_0)^{s}$. Any monotone saturating link with $\widetilde{p}(s) \approx p_0 s$ for small $s$ and $1 - \widetilde{p}(s) \sim e^{-a s}$ for large $s$ reproduces both regimes and the crossover \eqref{eq:crossover} with $a$ its tail rate; a link saturating only polynomially, $1 - \widetilde{p}(s) \sim s^{-\gamma}$, replaces the logarithmic crossover by a power-law one, $s_c(n) \sim n^{1/\gamma}$, and the exponential barrier persists over a far wider range of scores. The second, and the stronger, is full restart: that a single failure returns the process to stage~$1$. Let failure instead return the process $r$ stages, $j \mapsto \max(1, j - r)$, with $r$ bounded. For $j > r$ the chain is translation-invariant with drift $\widetilde{p} \cdot 1 - (1 - \widetilde{p})\, r$ per attempt, so
\begin{equation}\label{eq:rollback}
\mathbb{E}[M_n] \;=\; \Theta(n) \quad\text{iff}\quad \widetilde{p} \;>\; \frac{r}{1+r}, \qquad \mathbb{E}[M_n] \;\text{exponential in } n \quad\text{iff}\quad \widetilde{p} \;<\; \frac{r}{1+r}:
\end{equation}
the exponential barrier survives, but at a threshold in $\widetilde{p}$ that is \emph{independent of the depth} $n$---$\widetilde{p} = \tfrac12$ for $r = 1$, $\tfrac23$ for $r = 2$---and the $\ln n$ dependence of \eqref{eq:crossover} is the signature of unbounded rollback alone. At the threshold itself, $\widetilde{p} = r/(1+r)$, the drift vanishes and $\mathbb{E}[M_n] = \Theta(n^{2})$: quadratic, neither linear nor exponential. Full restart is the extreme case of memory loss, $r \to \infty$, and it is what makes the crossover climb with the depth. The empirical question the model therefore poses about a real task is not whether its steps are hard but how far a failure sets one back: a task whose failures lose $O(1)$ recent stages has a fixed reliability threshold and linear completion above it, while a task whose failures lose everything has the depth-dependent barrier of this section. The sequential-achievement interpretation of the historical cases below assumes the latter, and says so.

\subsection{Heterogeneous Achievement Sequences}\label{ssec:hetero}

In the general case, the achievement types
$\mathfrak{a}_1, \ldots, \mathfrak{a}_n$ may differ in category,
baseline probability, and occurrence rate. An individual may also have
different scores $s_j$ for different categories~$\mathcal{K}_j$ (though
in practice scores across related categories are highly correlated).
Let $\widetilde{p}_j = 1 - (1-p_{0,j})^{s_j}$ denote the modified
success probability at stage~$j$.

\subsubsection{Recursive Expected Trial Count}

Define $E_k$ as the expected number of events (from stage~$k+1$ onward,
having just completed stage~$k$) until completion. The recurrence is
\begin{equation}\label{eq:hetero-recurrence}
  E_k
  \;=\;
  1
  \;+\;
  \widetilde{p}_{k+1}\,E_{k+1}
  \;+\;
  (1 - \widetilde{p}_{k+1})\,E_0\,,
  \qquad
  k = 0, 1, \ldots, n-1,
  \qquad
  E_n = 0.
\end{equation}
This system can be solved by backward substitution. The system has the exact solution $E_0 = \sum_{j=1}^{n} \prod_{k=j}^{n} \widetilde{p}_k^{\,-1}$, which reduces to~\eqref{eq:expected-trials} when the $\widetilde{p}_k$ are equal; and the qualitative structure is preserved: $E_0$ is dominated by
$\prod_{j=1}^{n}\widetilde{p}_j^{-1}$ (the product of inverse
probabilities), which reduces to $\widetilde{p}^{-n}$ in the homogeneous
case and retains the exponential-to-linear crossover in the
general case.

\subsubsection{Expected Time with Heterogeneous Rates}

When the occurrence rates $\lambda_j$ differ across stages, the expected wall-clock time incorporates the waiting time at each stage. Stage~$j$ is attempted, on average, $\big(\prod_{k=j}^{n} \widetilde{p}_k\big)^{-1}$ times before completion, and each attempt there takes mean time $1/\lambda_j$, so
\begin{equation}\label{eq:hetero-time}
  \mathbb{E}[T_n] \;=\; \sum_{j=1}^{n} \frac{1}{\lambda_j \prod_{k=j}^{n} \widetilde{p}_k}\,,
\end{equation}
which reduces to $E_0/\lambda$ when the rates are equal. The dominant effect, the exponential-to-linear transition, is driven by the success probabilities rather than the rates.

\subsection{Connection to Section~\ref{sec:inference}}\label{ssec:connection}

The score $s$ entering the completion model is precisely the trait
value~$X$ from the log-normal population of
Section~\ref{sec:inference}, normalized by the population median:
$s = X / e^{\mu}$. The log-normal distribution of~$X$ therefore implies
that $\ln s \sim \mathcal{N}(0,\sigma^2)$, and the entire distribution
of completion times across the population is determined by the two
parameters $(\mu, \sigma)$ inferred in Section~\ref{sec:inference}
together with the achievement-specific parameters $(p_0, \lambda, n)$.

The record-holder amplification effect of Section~\ref{ssec:bolt} now
acquires a second, sharper interpretation. In Section~\ref{sec:inference},
the effect manifested as multiplicative volatility of the top performer's
score. Here, that multiplicative shift in~$s$ is raised to the $n$-th
power in the completion-time ratio~\eqref{eq:amp-ratio}, transforming a
modest advantage in latent ability into an overwhelming advantage in
realized output. The composition of log-normal scores with the
consecutive-success mechanism thus provides a quantitative explanation for
the empirical observation that output distributions are far more
concentrated than ability distributions.

\subsection{Illustrative Application: Extreme Mathematical
            Output}\label{ssec:ramanujan}

We conclude with an informal application to the phenomenon of extreme
outlier productivity in pure mathematics, exemplified by figures such as
Ramanujan.

Model the production of a significant mathematical result as a completion
type of depth~$n$: a sequence of~$n$ conceptual steps (insight, formal
development, verification, etc.) that must each succeed consecutively
without a fundamental error that would invalidate the chain. Let $p_0$
denote the per-step success probability for a median-ability mathematician,
and let $s_R$ denote Ramanujan's score in the relevant category.

If $p_0 = 0.05$ and $n = 8$ (a moderately deep result requiring eight
consecutive correct steps), then:
\begin{itemize}
  \item A median individual ($s = 1$) has $\widetilde{p} = 0.05$ and
        $\mathbb{E}[M_8] = \bigl((0.05)^{-8}-1\bigr)/(1-0.05) \approx 2.7 \times
        10^{10}$ trials.
  \item An individual at $s \approx 20$ (the 99.99th percentile when
        $\sigma \approx 0.8$) has
        $\widetilde{p} = 1 - (0.95)^{20} \approx 0.64$ and
        $\mathbb{E}[M_8] \approx 96$ trials.
  \item At $s \approx 60$ (with $a = -\ln 0.95 \approx 0.051$, the crossover $s_c(n) = \ln n / a$ is $\approx 45$ at $n = 10$ and $\approx 90$ at $n = 100$, so $s \approx 60$ is beyond the crossover at shallow depth and inside it at the deeper one),
        $\widetilde{p} \approx 0.954$ and $\mathbb{E}[M_8] \approx 10$
        trials, barely more than the depth itself.
\end{itemize}
The ratio of expected trials between the median individual and the
$s = 60$ scorer is approximately $2.7 \times 10^{9}$: roughly three
billion to one. Yet on the log-normal scale, the difference between $s = 1$ and
$s = 60$ is $\ln 60 \approx 4.1$ standard deviations if $\sigma = 1$;
extreme, but entirely within the support of a continuous distribution
over a population of billions.

This is the core explanatory claim: the individual is not operating
outside the human distribution. They are drawn from the same log-normal
population as everyone else, occupying an extreme but finite quantile.
What converts this quantile position into a qualitatively different level
of output is the consecutive-success mechanism, a structural feature of
the task itself, not a special property of the individual. The depth~$n$
of the required achievement sequence acts as an \emph{amplifier},
exponentiating modest differences in~$s$ into vast differences in expected
completion time. The individual's exceptional output is therefore
\emph{predictable from their position in the ability distribution}, given
knowledge of the task structure, without invoking any categorical
distinction between them and the rest of the population.

Under the divergence interpretation of Section~\ref{ssec:divergence}, Ramanujan's extraordinary output admits a particularly clean description. On that interpretation, which is a hypothesis about cognition rather than a result of the model, Ramanujan would not have possessed a special faculty for seeing mathematical truths that others lacked the organ to perceive; he would have lacked the divergence friction that prevents typical mathematicians from sustaining long chains of correct structural extrapolation. His celebrated ability to ``see'' formulas without proofs---the experience he attributed to the goddess Namagiri---would then be what frictionless inference feels like from the inside: the correct next step of the structure simply appears, because nothing is deflecting the thinker away from it. The formal proof, which others must laboriously construct step by step, is a post hoc decomposition of something Ramanujan grasped as a unit---not because he was operating outside the human distribution, but because his position within it placed him in the regime where consecutive steps merge into single perceived moves (Section~\ref{ssec:modified-p}).

% VERIFY before circulation: the Viennot course anecdote's details (venue, student count) against Bessis's own telling.
This account addresses a specific puzzle raised by Bessis \cite{bessis2024mathematica,bessis2025genes,bessis2026jaimungal}: in a semester-long course taught by Xavier Viennot at the \'Ecole Normale Sup\'erieure, thirty to forty gifted students proved one of Ramanujan's formulas using purely visual combinatorial objects, finding the result ``obvious at the end.'' Bessis takes this as evidence that Ramanujan's brain was not categorically different---that a sufficiently developed visual-intuitive pathway can reach the same results\footnote{I also want to affirm my agreement with Bessis \cite{bessis2024mathematica,bessis2025genes,bessis2026jaimungal} on his fundamental point that the conceptual ground which Ramanujan is traversing is, categorically speaking, not beyond the bounds of reasonable reality, and so could hypothetically be traversed by anyone. The problem in practice, then, is not that Ramanujan's trodden ground is incomprehensible or otherwise mystically unavailable, but very simply that normal thinkers lack the sufficient pathfinding gifts to \emph{find} the conceptual road, and then to \emph{stay on} the road without almost surely spontaneously fraying off it. This then indeed reveals a peculiarity bordering on existentially jarring: \emph{the very path which Ramanujan walked exists for anyone to walk, it just remains combinatorially implausible that anyone else will find it}. This reminds me of a common thought experiment I come back to, being the question of what modern mathematics would look like if Galois had never existed. In effect, we would be forever burdened by the having to use the elements of his ideas in isolation, without ever knowing that direct paths through this space existed but had not been found. Accordingly, then, the thing that keeps me up at night is that I would frankly bet on our having missed one or two (or maybe more) Galois theory-caliber insights. Ramanujan's path itself serves perhaps as a canonical example of a fully-comprehensible path which remains obscured in the overgrowth of artificial complexity.}. The present model vindicates this assessment at one level and refines it at another. In the divergence interpretation of Section~\ref{ssec:divergence}, the structure of a sufficiently constrained inferential problem admits a single overwhelmingly dominant path---a route whose probability of success at each step dwarfs all alternatives by orders of magnitude. The fragments of this path that the students find are genuinely of the same kind as the concrete inferential moves Ramanujan makes: the same path, the same structural insights, the same results. Bessis is right that the path is humanly accessible and that the cognition involved is not categorically special. However, there is a subtle but important distinction between finding oneself on the path and placing oneself on it. The students' alignment with the dominant route, while genuinely impressive and far beyond what an untrained individual could achieve, retains an element of fortunate convergence---their process arrives at the right structural moves in part because the pedagogical scaffolding of the course channels them there. Between these guided segments, the underlying frizzing tendency reasserts itself. For Ramanujan, by contrast, alignment is self-directed: his mental process actively orients onto the dominant path and holds it there, so that the good segments are the default rather than the intermittent achievement. The greater depth he reaches is emergent from this sustained self-directed orientation, not the other way around. And yet---returning to the central thesis of this paper---this distinction is one of degree, not of kind. The capacity to orient onto the dominant path varies continuously across the population. The students in Viennot's course are not on the far side of a categorical divide; they are located at an intermediate position on a smooth spectrum, possessing enough sustained orientation to hold the path for stretches that are already far beyond the population baseline. The consecutive-success mechanism is what converts smooth variation along this spectrum into the discrete-looking gap between those who can produce results at Ramanujan's depth and those who cannot.

\subsection{Remark on Dispersion Amplification, and the Pareto Appearance}\label{ssec:pareto}
Bessis \cite{bessis2024mathematica, bessis2025genes, bessis2026jaimungal} has observed that mathematical talent, as subjectively perceived by practicing mathematicians, follows a distribution closer to Pareto than Gaussian: the gaps between individuals feel like orders of magnitude, not percentages. He takes this to be incompatible with direct genetic determination, since highly heritable polygenic traits produce Gaussian outcomes; the inference needs care, because a nonlinear transformation from latent inputs to output can change the output distribution whatever caused the inputs, and the mechanism below is such a transformation. The consecutive-success mechanism supplies a bridge, and what it supplies should be stated exactly, because it is not a literal conversion of a Gaussian into a Pareto law. Let $\ln s \sim N(0, \sigma^2)$ across the population, and consider the completion \emph{rate} $R(s) = 1/\mathbb{E}[T_n(s)]$. In the low-score regime $s p_0 \ll 1$, $\widetilde{p}(s) \approx p_0 s$ and $R(s) \approx \lambda (p_0 s)^{n}$---to leading order in the logarithm only, since the link's relative error, of order $p_0 s$, is raised to the $n$th power and shifts $\ln R$ by about $n p_0 s/2$, so that the statement below holds where the population mass lies in this regime and that shift is small or nearly constant across it---so
\begin{equation}\label{eq:dispersion}
\ln R \;\approx\; \ln\lambda + n \ln p_0 + n \ln s, \qquad \operatorname{Var}(\ln R) \;\approx\; n^{2} \sigma^{2}:
\end{equation}
the log-rate is Gaussian with standard deviation $n\sigma$. The mechanism therefore maps a log-normal of modest dispersion ($\sigma \approx 0.59$) to a log-normal of enormous dispersion, the task depth multiplying log-dispersion linearly and log-variance quadratically, in that approximation---an \emph{amplification} theorem, not a change of family. And the family cannot become Pareto in the model as stated: in the high-score regime $R(s) \to \lambda/n$, so the rate is bounded above and has no power-law upper tail at all. What the mechanism does produce, where the relevant population mass lies in the low-score regime, is an approximately log-normal rate with log-dispersion of order $n\sigma$, broad enough that over a limited empirical window its rank plot can be hard to tell from a power law, whether a given window distinguishes them depending on its range, the sample size, the curvature, and the saturation at $\lambda/n$; the Pareto structure Bessis reports is, on this account, the \emph{appearance} of a very broad log-normal over the range a population actually exhibits, and whether a given dataset's finite window is fit better by the amplified log-normal or by a genuine Pareto is an empirical question the model poses rather than settles. The felt orders-of-magnitude gaps are then $\exp(n\sigma \cdot \Delta z)$ for a difference $\Delta z$ of a few standard scores, which at the calibrated depths is the structure Bessis describes.

In Bessis's informal language, the ``capitalization process'' of stacking abstractions on top of abstractions is the consecutive-success chain: each abstraction must be correct and must build on its predecessor, and a single fundamental error resets the process. The ``neuroplasticity'' that Bessis invokes as the mechanism is, in the present framework, captured by the score~$s$: a measure of how completely the individual's cognitive processing closes the gap between baseline probability and certainty at each step. The felt-compaction concept of Section~\ref{ssec:modified-p} provides a concrete interpretation of what Bessis describes as Ramanujan's ability to ``see'' formulas without proofs: at sufficiently high~$s$, the steps of the inferential chain merge into single perceived moves, and the formal proof---the decomposition into individually verifiable steps---becomes a post hoc unpacking of something grasped as a unit.

This mechanism has an important ancestor that should be acknowledged and distinguished. Shockley \cite{shockley1957}, observing that publication rates in research laboratories are log-normally distributed with some individuals producing at fifty times the rate of others, proposed that scientific production requires the conjunction of several sequential factors---thinking of a problem, working it, recognizing a result, writing it up, and so on---so that individual output is a \emph{product} of factor-level competences, and modest normal variation in the factors compounds multiplicatively into log-normal variation in output. The present framework agrees with Shockley's structural intuition that sequential conjunction is the amplifier, but differs in both the placement of the log-normal and the resulting output family: for Shockley the log-normal is the \emph{explanandum} (output), whereas here it is the \emph{input} (the latent score distribution across the population), and the consecutive-success process with reset-on-failure then amplifies that input's log-dispersion by the task depth, in a suitable low-score range, through the exact nonlinear transformation of equation~\eqref{eq:expected-trials}, which departs from log-normality and saturates at a finite upper rate. The reset dynamics, the depth-dependent crossover, and that depth-proportional amplification with its saturation are the additions; whether the output's finite-range appearance is Pareto-like is a hypothesis to examine, not an established change of family. Relatedly, Simonton's equal-odds rule \cite{simonton1997}---that major works arrive in proportion to total output, quality tracking quantity stochastically---is complementary rather than competing: it governs the hit rate \emph{within} a realized body of work, while the present mechanism governs how large a completed body of work becomes at all; the two compose without tension.

One further relation must be stated, because it is the null hypothesis
any reader from the sociology of science will reach for. Cumulative
advantage---Merton's Matthew effect \cite{merton1968matthew}, formalized
by Price \cite{price1976cumulative}---generates Pareto-like output from
a population of \emph{equal} latent ability through recognition
feedback: visibility begets resources begets output begets visibility.
The present mechanism generates broad, heavy-looking output from latent heterogeneity with \emph{no} feedback anywhere in the process: the score
is fixed, the trials are independent, and nothing about early success
alters later probabilities. The two accounts are therefore rivals less over the shape of the output, which both make broad over a finite range, than over its \emph{source}, and they separate empirically. Cumulative advantage predicts
concentration that grows with career age and is mediated by visibility
channels---institutional placement, citation networks---while the
consecutive-success mechanism predicts concentration that is present
from the first completed work, tracks task depth rather than accumulated
recognition, and does not require the feedback channels to exist at all;
an isolated clerk in Madras is exactly the configuration cumulative
advantage cannot reach and this mechanism is indifferent to. In
realistic careers both surely operate, and nothing here denies the
feedback; the claim is that the latent-amplification channel can by itself produce output broad enough to look Pareto-like over a finite window, with the classical empirical regularities of
productivity concentration \cite{lotka1926, oboyle2012} neutral between
the sources until their age- and depth-dependence is examined. The
closest modern quantitative neighbor on the latent side is the Q-model
of Sinatra et al.\ \cite{sinatra2016quantifying}, in which a stable,
log-normally distributed individual parameter multiplies stochastic
project potential; the structural difference is the one that separated
the present account from Shockley's: the Q-model's composition is a product with no reset, its dispersion fixed by the product's factors, whereas here the reset dynamics make the amplification grow with task depth and saturate above a depth-dependent crossover.

\subsection{Remark on Continuity, Irreplaceability, and the Rejection of
            Categorical Exceptionalism}\label{ssec:exceptionalism}

The model developed in this section produces two conclusions that are
usually treated as contradictory but are in fact simultaneous and
inseparable.

\medskip
\noindent\textbf{Conclusion 1: Continuity.}\;
The individual is not a categorical alien. The consecutive-success process
is a universal mechanism applied uniformly to all individuals; it does not
posit that outliers are drawn from a different distribution, possess
qualitatively different cognitive architecture, or benefit from supernatural
endowment. The entire explanatory burden is carried by two features:
\begin{enumerate}
  \item[(i)]
    the \emph{continuous} log-normal distribution of ability, which
    permits extreme scores without categorical discontinuity; and
  \item[(ii)]
    the \emph{nonlinear amplification} inherent in the
    consecutive-success requirement, which converts continuous differences
    in ability into discrete-looking differences in output.
  \item[(iii)]
    Under the divergence interpretation of Section~\ref{ssec:divergence}, the score itself reflects not the presence of an extraordinary capacity but the \emph{absence} of a common impediment. The genius is not superhuman; the genius is unimpeded. This further demystifies the outlier by locating the explanatory burden not in the rare individual who succeeds but in the ubiquitous tendency to diverge that prevents everyone else from doing so.
\end{enumerate}

\noindent\textbf{Conclusion 2: Irreplaceability.}\;
Precisely \emph{because} the amplification exponent is the task depth~$n$,
the model simultaneously implies that the output of an extreme outlier is \emph{not practically reproducible by individuals at lower quantiles}: within any specified population and time budget of realistic size, additional lower-quantile workers provide negligible compensation. With unlimited workers or time, completion by someone becomes certain, since every individual with $\widetilde{p} > 0$ completes eventually. The
ratio~\eqref{eq:amp-ratio} is polynomial in~$s$ of degree~$n$: an individual at $s = 20$ cannot replicate the output of an individual at $s = 60$ by investing three times the effort: at the calibration of Section~\ref{sec:results}, $n = 7{,}700$, the exact ratio of their expected completion times is about $10^{145}$. (The low-score form $s^{n}$ of \eqref{eq:amp-ratio} does not apply at these scores, where $s\,p_0 \gg 1$; it would give $3^{n} \approx 10^{3674}$.) The difference is one of degree, made enormous by the depth, and \eqref{eq:horizon} says how negligible lower-score effort is within any realistic budget. The model therefore \emph{strengthens} the case that the loss of
such an individual represents a genuine, quantifiable deficit to the
collective record, not because they are mystically different, but because
the mathematics of consecutive success is unforgiving.

\medskip

These two conclusions together constitute a precise middle position
between two inadequate extremes. The first extreme, sometimes associated
with na\"{\i}ve readings of the ``great man theory \cite{hook2004great},'' treats outlier
individuals as categorically distinct from the population, a framing
whose historical entanglement with the systematic exclusion of women,
people of color, and other marginalized groups from education and
scientific institutions makes it both intellectually bankrupt and
actively harmful. The evidence base from which ``great individuals'' were
historically identified was corrupted by a selection bias so profound that
any causal claims derived from it should be treated with corresponding
skepticism; the conflation of denied access with diminished ability is
among the most persistent and damaging errors in the history of ideas.
The present model repudiates this framing categorically: the individual
is a draw from a universal human distribution, full stop.

The second extreme, the strong inevitability thesis (which holds that any
sufficiently developed collective will eventually produce any given
advance), underestimates the severity of the exponential barrier. The
model shows that for tasks of sufficient depth, the expected completion
time for individuals below the crossover score~$s_c(n)$ exceeds not just
human lifetimes but cosmological timescales. The advance is not
``inevitable''; it is contingent on the existence of a specific draw from
the tail.

The resolution is that both the individual and the collective are
structurally indispensable. The collective sustains the infrastructure,
accumulates the empirical base, and carries out the vast majority of
productive work. The extreme outlier, operating within exactly the same
mechanistic framework as everyone else, happens to sit at a quantile where
the consecutive-success process transitions from impossible to routine.
Neither can substitute for the other. To insist on choosing between them
is to ask whether the heart or the spine is more essential to survival:
removing either is catastrophic, and the question itself betrays a
misunderstanding of the system's architecture.



\section{Results, and an Illustrative Calibration}\label{sec:results}

We now apply the general framework of
Sections~\ref{sec:inference}--\ref{sec:completion} to a concrete empirical
setting. The purpose of this section is twofold: to show, as a numerical illustration from stipulated anchors, that the methodology produces internally consistent parameter values, and to illustrate the sharp crossover of Section~\ref{ssec:phase} in a domain where both the ability distribution and the output distribution could in principle be independently assessed.

\medskip
\noindent\textbf{Status of the calibration, declared once.}\;
The numbers in this section are offered at a stated grade, fixed here so
that it governs every table below. The two empirical anchors (E1 and E2)
are stipulations---defensible order-of-magnitude judgments, not
measurements---and every derived figure (the depth $n = 7{,}700$, the
$14$-second micro-step, the $10^{6440}$-year median entry) is an
internal-consistency consequence of the model plus those stipulations.
Two further cautions keep the calibration illustrative. The collective's $20\%$ replacement probability is one of the calibration's own equations, so the later finding that the collective has roughly that probability is a consistency check and not independent empirical support. And the identification of the theoretical-physics population's mode with the $1 - |B|/N$ quantile of latent ability treats the world population as a random sample of people with equal access to becoming theoretical physicists, which it is not; the quantile is a stipulated placement, not a measured one, as is the ``$3.5$ times the elite mode'' scalar, which corresponds to no observable in any conventional dataset. What would replace these with measurements is three separable empirical hypotheses, each testable on its own: that a latent score has multiplicative, Gumbel-domain extremes across populations of varying size; that individual success reliability follows a saturating link in that score; and that real high-level tasks exhibit the long-range restart of Section~\ref{ssec:sensitivity} rather than bounded rollback. The historical interpretation is evidence for the mechanism only once those three are; until then it is an illustration of what the mechanism would say.
Such figures carry the model's \emph{structure}: the existence and
sharpness of the crossover, the ordering of the quantiles, and
the depth-sensitivity of the gaps. They do not carry measurement-grade
content, and no digit beyond the leading order of magnitude should be
read as if they did. The robustness sweep reported below (stability of
the qualitative conclusions across $k \in [2.5, 5.0]$) is the check that
the structure, as opposed to the digits, survives variation in the
anchors.

\subsection{Theoretical Physics}\label{ssec:physics}

\subsubsection{Section~\ref{sec:inference} Parameter Estimates}

We take as our parent population~$A$ the world population, with
$N \approx 8 \times 10^9$. The elite subpopulation~$B$ is identified with
the set of actively publishing theoretical physicists worldwide. Drawing
on data from the American Institute of Physics \cite{aip2024trends} and
Physics Today \cite{physicstodaymillion}, we estimate $|B| = n \approx
10^5$ as an order-of-magnitude figure.

%%% Substitute your empirical values for m_B below.
%%% For the present analysis we normalize $m_B = 1$.
The quantile of the elite mode within~$A$ is
$q = \Phi(4.2148) \approx 0.99999$, giving $z_q = 4.2148$.
The single extreme observation is that of Einstein, whose peak
contribution we estimate at $v_{\max} = k\,m_B$ with $k = 3.5$.
This is a single-observation estimate, subject to the uncertainty
quantified in Section~\ref{ssec:uncertainty}: the coefficient of
variation of the population maximum is approximately $12\%$, so the true
$k$ plausibly lies in the range $3$ to $4$. Numerical experimentation
confirms that the qualitative results below are robust well beyond this
range: recalibrating at each $k \in [2.5,\, 5.0]$ so that the two empirical constraints (E1 and E2 below) hold, which fixes Einstein's completion time and the collective probability by construction, the $B$-mode completion time remains astronomically large throughout, from about $10^{270}$~years at $k = 2.5$ to about $10^{970}$~years at $k = 5$, against $10^{537}$~years at the central $k = 3.5$. The parameters that absorb the variation are $\sigma$ (the
log-space dispersion), $p_0$ (the baseline success probability), and
$\lambda$ (the step rate), whose calibrated values shift but whose
product remains constrained by the empirical anchors. Only the precise
numerical entries in the tables change; the crossover, the
order-of-magnitude separations between quantiles, and the comparison
across revolutions are unaffected.

Evaluating \eqref{eq:zmax} at $N = 8 \times 10^9$ gives
$z_{\max} = 6.3343$. The log-space solutions
\eqref{eq:sigma-sol}--\eqref{eq:mu-sol} then yield
\begin{equation}\label{eq:physics-sigma}
  \sigma
  \;=\;
  \frac{\ln 3.5}{6.3343 - 4.2148}
  \;=\;
  \frac{1.2528}{2.1195}
  \;\approx\; 0.5911\,,
\end{equation}
and $\mu = \ln m_B - 0.5911 \times 4.2148$.

Einstein's score relative to the population median is
\begin{equation}\label{eq:einstein-score}
  s_E
  \;=\;
  k \cdot e^{\sigma z_q}
  \;=\;
  3.5 \cdot e^{0.5911 \times 4.2148}
  \;\approx\; 42.3\,.
\end{equation}
The mode of the elite subpopulation~$B$ (at the selection threshold)
corresponds to a score of $s_{\mathrm{mode}} \approx 12.1$, and the
mean score within~$B$ is approximately $13.7$.

The relative uncertainty in $\sigma$ from \eqref{eq:rel-uncert-sigma} is $\delta\sigma/\sigma \approx 9.0\%$ once the factor $z_{\max} - z_q \approx 2.12$ is carried, and the first-order multiplicative uncertainty scale of the maximum from \eqref{eq:bolt-spread} is
$\mathrm{CV}(X_{(N)}) \approx 11.9\%$.

The score $s_E \approx 42.3$ is the \emph{typical} maximum, the log-maximum matched at its typical value~\eqref{eq:emax} and not at its expectation. As shown in
Section~\ref{ssec:gap}, the actual maximum is a random variable whose gap
to the second-ranked individual is exponentially distributed with scale
$\beta_N \approx 0.088$. For the calibration below, we work with a
representative Gumbel-process realization in which the record-holder's
score is $s_{(1)} \approx 47.7$ (corresponding to an upward fluctuation
of roughly one Gumbel standard deviation) and the second-ranked
individual is at $s_{(2)} \approx 39.7$, a $20\%$ gap. This gap lies
in the upper tail of the Gumbel-process gap distribution, not near its median
(Section~\ref{ssec:gap}, equation~\eqref{eq:median-gap}: median gap
$\approx 6.3\%$, and $P(g > 20\%) \approx 12.6\%$ by \eqref{eq:gap-tail}) and is chosen to illustrate the regime in which
the crossover produces maximal separation between the
record-holder and the field.

\subsubsection{Section~\ref{sec:completion} Parameter Calibration}

We model the production of a foundational theoretical advance, such as
the formulation of a new physical theory from ontological first
principles, as a completion type of depth~$n$ with homogeneous
achievement parameters $(p_0, \lambda)$. The Section~\ref{sec:inference}
parameters are held fixed; the Section~\ref{sec:completion} parameters
are calibrated from two empirical constraints:
\begin{enumerate}
  \item[\textbf{E1.}]
    \emph{Einstein's working time.}\;
    We hypothesize that Einstein, operating at his peak and working
    without interruption, could have produced a foundational advance
    (comparable in depth to special relativity or the pending \cite{spekkens2010einstein} realist quantum program) in approximately $500$
    person-hours of focused effort, roughly 21 working days, dispersed
    over a longer calendar period of say ten months to account for rest and other
    activities. This gives $\mathbb{E}[T_n \mid s_E] = 500$ hours, with time spent directly working occupying about 7\% of the total elapsed time over a given period, optimistically speaking.
  \item[\textbf{E2.}]
    \emph{Collective completion probability.}\;
    In the Gumbel-process model of Section~\ref{ssec:gap}, the top
    contemporaneous scores are random, with the record-holder separated
    from the field by an exponentially distributed gap. The second-ranked
    individual occupies a transitional position: close enough to the
    record-holder that their completion time is intermediate (typically
    centuries), yet far enough below that they cannot contribute on a
    human timescale. The operative collective is therefore the
    \emph{cluster} of competitors from rank~$\#3$ onward, whose scores
    are tightly bunched (Section~\ref{ssec:gap}, third consequence).
    We require that ranks $\#3$ through $\#50$, working at their respective rates over $1{,}000$ years of on-task time, the unit of the tables below, have a cumulative completion
    probability of approximately $20\%$:
    $1 - \exp(-R_{\mathrm{coll}} \cdot T_{1000}) = 0.20$, where
    $R_{\mathrm{coll}} = \sum_{j=3}^{50} \lambda / \mathbb{E}[M_n \mid s_{(j)}]$
    and $T_{1000} \approx 8.77 \times 10^6$ on-task hours, about twelve thousand calendar years at two on-task hours a day. Including the
    second-ranked individual raises the collective probability to
    approximately $80\%$, confirming that the second-ranked individual is
    the decisive swing factor.
\end{enumerate}

\noindent
The scores of the top $50$ contemporaneous individuals are drawn from
the Gumbel order-statistic process described in Section~\ref{ssec:gap},
conditioned---as an illustrative choice, not an inference, per Section~\ref{ssec:uncertainty}'s caution against using the maximum twice---on a representative realization in which the record-holder's
gap exceeds $15\%$ (occurring in approximately $20\%$ of draws from a population of $8 \times 10^9$ at $\beta_N \approx 0.088$; see~\eqref{eq:gap-tail}). In the
representative realization used below, the record-holder has score
$s_{(1)} = 47.7$, the second-ranked individual has $s_{(2)} = 39.7$
(a $20.0\%$ gap), and ranks $\#3$ through $\#10$ cluster between $39.0$
and $34.2$ (a $14\%$ spread).

The two constraints fix two of the three task parameters, not all three. Given the depth $n$, eliminating $\lambda$ via $\lambda = \mathbb{E}[M_n \mid s_{(1)}] / 500$ and matching the collective constraint determines $p_0$ and $\lambda$; the depth itself is not determined by them, and the value $n = 7{,}700$ used throughout is an illustrative choice, so that the triple below is one member of a family of calibrations indexed by $n$. For $n = 7{,}700$ the calibration gives:
\begin{equation}\label{eq:calibrated}
  p_0 \approx 0.146\,,
  \qquad
  \lambda \approx 252\;\text{steps/hour}\,,
  \qquad
  s_c(n) = \frac{\ln n}{-\ln(1-p_0)} \approx 56.7 \quad (n = 7{,}700)\,.
\end{equation}
The depth is $n = 7{,}700$ consecutive micro-steps, each lasting
approximately $14$~seconds on average. These are rapid inferential
micro-evaluations: assessments of whether a candidate logical move
coheres with the surrounding structure, analogous to the ``blitz checks''
a chess grandmaster performs when scanning candidate positions.
Einstein's $500$ hours of focused work consist of approximately
$126{,}000$ such evaluations, of which $7{,}700$ must succeed
consecutively. He restarts on average about $70$ times, $\widetilde{p}^{\,-n} - 1$, most failed runs ending early, so that his expected number of evaluations is $\mathbb{E}[M_n] \approx 16.4\,n$.

\subsubsection{Results}
\noindent\textit{How the tables are computed.} Quantile scores are $s_q = e^{\sigma z_q}$ with $\sigma = 0.5911$ and $z_q = \Phi^{-1}(q)$, and every completion time in the tables of this section is computed from the unrounded score, the displayed scores being rounded. At these depths the rounding is consequential: the $90$th-percentile score $2.133$ gives $10^{4192}$ years, while $2.1$ would give $10^{4236}$.


Table~\ref{tab:physics-completion} reports expected completion times by
population quantile.

\begin{table}[ht]
  \centering
  \caption{Expected completion time (on-task working hours) for a
           foundational theoretical advance ($n = 7{,}700$, $p_0 = 0.1455$,
           $\lambda = 252\,$hr$^{-1}$). Ranks $\#2$--$\#20$ are from
           a representative Gumbel-process realization with a $20\%$ gap
           between the record-holder and the second-ranked individual.
           Calendar time is longer by a factor of $10$--$20\times$ or more
           depending on the fraction of time devoted to the specific
           problem\protect\footnotemark. Note that we consider the sequential achievement process itself to run only during working hours, as if pausing during absent hours as opposed to our slowing its intrinsic rate down to fit that time uniformly.}
  \label{tab:physics-completion}
  \medskip
  \begin{tabular}{l r r r}
    \hline
    Quantile & Score $s$ & $\widetilde{p}(s)$ &
    $\mathbb{E}[T_n]$ \\
    \hline
    Median (50th) & $1.000$ & $0.146$ & $10^{6440}$ years \\
    90th percentile & $2.133$ & $0.285$ & $10^{4192}$ years \\
    99th percentile & $3.956$ & $0.463$ & $10^{2568}$ years \\
    99.9th percentile & $6.213$ & $0.624$ & $10^{1574}$ years \\
    $B$ mode (threshold)  & $12.1$  & $0.850$ & $10^{537}$ years \\
    \hline
    $\#20$ contemporaneous & $32.6$ & $0.994$ & $10^{16}$ years \\
    $\#10$ contemporaneous & $34.2$ & $0.995$ & $10^{11}$ years \\
    $\#5$ contemporaneous  & $35.1$ & $0.996$ & $10^{9}$ years \\
    $\#3$ contemporaneous  & $39.0$ & $0.998$ & $3{,}877$ years \\
    $\#2$ contemporaneous  & $39.7$ & $0.998$ & $755$ years \\
    Einstein (record-holder) & $47.7$ & $0.9994$ & 21 days \\
    \hline
  \end{tabular}
\end{table}
\footnotetext{
A good estimate might be that no more than two hours on average are spent on the project per day, on the hypothesis, conditional on the appendix's model and not established here, that the high-achiever is at origin more likely than not to be found outside academia. The supposition doing the work here is modest, and it is worth stating at exactly its strength: not that academic training is harmful, but that academia's selection operates on measured proxies that, per the appendix, decouple from the sequential-achievement criterion in the relevant tail---so that selection on them buys little enrichment precisely where enrichment would be needed. Per the population-weighting argument of Section~\ref{app:population}, this has a consequence stronger than mere non-enrichment. Academia is a small, sharply selected subpopulation, and---crucially---its selection criteria are measure-side: credentialing, examination performance, and conformity to standard training select on observed proxies of precisely the kind the appendix shows decouple from the sequential-achievement criterion in the relevant tail. Above the decorrelation onset, then, selection on those proxies buys essentially no lift in the per-person probability of satisfying the criterion, while the smallness of the selected population is paid in full: the expected number of criterion-satisfiers a subpopulation contains is the per-person probability times the population size, and with the first factor approximately flat across the selected range, the second factor dominates the count. The high-achiever person's probability of being in academia is therefore at most equal to the proportion of people who are in academia (or lower than this in the anticorrelation scenario), provided membership depends on the criterion only through those measured proxies. A plateau in the score's predictiveness does not establish that proviso: membership can depend on other variables relevant to the criterion, such as training, interest and opportunity, and where it does, the conclusion below does not follow. On the proviso, ironically enough, the heavy in-group dynamic of academia counts against it per its low population: the next Ramanujan-caliber individual is likelier, at origin, to be found outside academia simply because that is where almost everyone is. The same weighting bears on the measured profile of such individuals---the modal transformative observed score sits at the plateau onset under the proposition's conditions (Proposition~\ref{app:prop-mode})---so a population purpose-built to concentrate high measured scores is, doubly, not where the criterion concentrates. 
}

We highlight several aspects of these results.

\medskip
\noindent\textbf{The crossover is visible.}\;
The expected completion time drops by over \emph{six thousand
orders of magnitude} between the population median ($10^{6440}$~years)
and Einstein ($21$~days). The transition from the exponential regime
to the linear regime accounts for the entirety of this drop. The
typical theoretical physicist, at the $B$-mode score $s \approx 12$
with $\widetilde{p} \approx 0.85$, faces $10^{537}$~years. The
third-ranked contemporaneous individual, at $s \approx 39$ with
$\widetilde{p} \approx 0.998$, faces $3{,}877$~years. Einstein, at
$s \approx 48$ with $\widetilde{p} \approx 0.9994$, does it in three
weeks. The difference between $\widetilde{p} = 0.998$ and
$\widetilde{p} = 0.9994$ is tiny in absolute terms but is amplified by $n = 7{,}700$ consecutive steps into a factor of about $7 \times 10^4$ in completion time.

\medskip
\noindent\textbf{Einstein operates inside the crossover, not beyond it.}\; The crossover score at this depth is $s_c \approx 56.7$, and Einstein's score is $s_{(1)} \approx 47.7$, so $n(1-p_0)^{s_{(1)}} \approx 4.1$ and the scaling function \eqref{eq:scaling-function} gives an overhead $F(4.1) \approx 15$ over the ideal $n/\lambda$---the sixteen task-depths' worth of evaluations the figures below exhibit.
His effective per-step probability $\widetilde{p} \approx 0.9994$
means that in a chain of $7{,}700$ micro-steps, the expected number
of restarts is small: the expected total evaluations
$\mathbb{E}[M_n] \approx 126{,}000$, or roughly $16$ times the
depth. At $\lambda \approx 252$ steps per hour (one evaluation every
$14$~seconds), each micro-step is a rapid inferential check: does this
candidate move cohere with the surrounding logical structure? A
completion of depth $7{,}700$ at this rate represents approximately
$31$~hours of uninterrupted evaluation for a nearly error-free
individual, with the remaining hours of Einstein's $500$-hour total
accounted for by the roughly $70$ restarts, most of them ending early, and their re-traversals.

\medskip
\noindent\textbf{The Bolt gap separates Einstein from the field.}\;
The Gumbel-process gap (Section~\ref{ssec:gap}) places the
second-ranked individual at $s_{(2)} \approx 39.7$, a $20\%$ gap in
score but a factor of $13{,}000$ in completion time for the base
calibration ($755$~years vs.\ $21$~days). The third-ranked individual
clusters just below at $s_{(3)} \approx 39.0$ ($3{,}877$~years), and
ranks $\#4$ through $\#10$ drop rapidly into the billions of years. The
collective of ranks $\#3$--$\#50$ has a $24\%$ probability of completion within $1{,}000$ on-task years; including the second-ranked
individual raises this to $82\%$, confirming that the second-ranked
individual is the decisive swing factor between ``the revolution was
inevitable'' and ``only one person could have done it.'' (This
$20\%$-gap realization is representative of the high-depth regime;
at shallower depths or smaller gaps, the collective's probability
increases substantially, as discussed below.)

\medskip
\noindent\textbf{The collective contribution is dominated by the tail.}\;
The rate concentration is extreme: Einstein alone accounts for
$99.9\%$ of the total progress rate. The second-ranked individual
contributes $0.1\%$; ranks $\#3$--$\#50$ together contribute less than
$10^{-4}\%$. Yet the collective of ranks $\#3$--$\#50$ still has a
non-negligible ($24\%$) probability of completion within $1{,}000$ years
because the Poisson-process model accumulates rate over a long
time window. The task does not decompose into parallelizable subtasks:
each step must follow the last, and a single error resets the chain.
But many independent individuals can each attempt the chain, and their
rates sum. The collective's $24\%$ arises almost entirely from the
third-ranked individual ($3{,}877$~years); the remaining $47$ people
contribute negligibly.

\subsubsection{Comparison Across Historical Revolutions}

\noindent\textit{What these figures assume.} Every figure in this comparison is conditional on the full-restart assumption and on the illustrative calibration above. Under bounded rollback, Section~\ref{ssec:sensitivity}, the depth-dependent barrier is replaced by a fixed reliability threshold above which completion is linear in the depth, and the separations below would shrink accordingly; retained partial results, error correction, collaboration and reusable intermediate discoveries all move a real task in that direction. The comparison therefore illustrates what the model implies under its assumptions. It is not independent evidence about the historical cases, about any individual's latent score, or about how long a field would otherwise have taken.


The calibrated parameters $(p_0, \lambda)$ are properties of the
achievement domain (theoretical physics) and the population ($B$), not of
any particular task. We may therefore hold them fixed and vary only the
completion depth~$n$ to compare revolutions of different complexity.

A subtlety arises for Newton. Einstein's score $s_{(1)} \approx 47.7$
is the record-holder from a representative Gumbel-process realization of
$N = 8 \times 10^9$ contemporaneous individuals. Newton represents the
top performer of the cumulative pool since approximately 1000~CE,
roughly $N_{\mathrm{Newton}} \approx 5 \times 10^{10}$
\cite{haub2022howmany}. Taking Newton as the typical maximum of this larger population gives
\begin{equation}\label{eq:newton-score}
  s_N
  \;=\;
  e^{\sigma \cdot z_{\max}(N_{\mathrm{Newton}})}
  \;\approx\; 49.8\,,
\end{equation}
with $\widetilde{p}(s_N) \approx 0.9996$.

Table~\ref{tab:revolutions-summary} reports completion statistics for four
historical milestones. The depths are expressed in micro-steps, scaled
proportionally from the base calibration ($n = 7{,}700$ for a
depth-$60$ advance).

\begin{table}[ht]
  \centering
  \caption{Completion across revolutions of varying depth
           ($p_0 = 0.1455$, $\lambda = 252\,$hr$^{-1}$). Each revolution
           is evaluated at the score of its historical author and at the
           third-ranked contemporaneous individual.}
  \label{tab:revolutions-summary}
  \medskip
  \begin{tabular}{l r r r r}
    \hline
    Revolution & $n_{\mathrm{eff}}$ &
    Einstein & Newton & $\#3$ cluster \\
    \hline
    Special relativity   & $5{,}949$  &    8 days &    4 days &    86 years \\
    Foundational advance & $7{,}700$  &   21 days &    9 days & $3{,}877$ years \\
    General relativity   & $8{,}859$  &   40 days &   14 days & $48{,}000$ years \\
    Newtonian mechanics  & $11{,}453$ &  168 days &   40 days & $10^{7}$ years \\
    \hline
  \end{tabular}
\end{table}

Special relativity ($n_{\mathrm{eff}} = 5{,}949$) requires
approximately eight days of focused work,
consistent with the historical record of the theory's emergence
in a concentrated burst during 1905. General relativity
($n_{\mathrm{eff}} = 8{,}859$) requires approximately $40$~days,
consistent with the well-documented false starts (the Entwurf field
equations, the hole argument \cite{pais1982subtle}) that preceded the
final formulation in late 1915.

Newtonian mechanics ($n_{\mathrm{eff}} = 11{,}453$)
is the deepest completion in the table,
yet Newton's higher score ($s_N \approx 49.8$ versus Einstein's
$s_{(1)} \approx 47.7$) compresses it to $40$~days.
Einstein himself would require $168$~days on this task, nearly six months.
The depth is greater, but Newton is more extreme, and the
compound effect is that Newton experiences the \emph{Principia}'s
derivation with the same fluency that Einstein experiences
general relativity. This is a direct consequence of the model's
structure: the score and the depth interact multiplicatively, so a
sufficiently high score can ``absorb'' additional depth without crossing
back into the exponential regime.

Tables~\ref{tab:sr-full}--\ref{tab:newton-full} report selected
completion-time profiles for each revolution.

\begin{table}[ht]
  \centering
  \caption{Special relativity ($n_{\mathrm{eff}} = 5{,}949$).}
  \label{tab:sr-full}
  \medskip
  \begin{tabular}{l r r r}
    \hline
    Quantile & Score $s$ & $\widetilde{p}(s)$ &
    $\mathbb{E}[T_n]$ \\
    \hline
    Median (50th)              &  $1.0$  & $0.146$ & $10^{4974}$ years \\
    99.9th percentile          &  $6.2$  & $0.624$ & $10^{1218}$ years \\
    $B$ mode                   & $12.1$  & $0.850$ & $10^{414}$ years \\
    $\#3$ contemporaneous      & $39.0$  & $0.998$ & 86 years \\
    $\#2$ contemporaneous      & $39.7$  & $0.998$ & 25 years \\
    Einstein ($s_{(1)}$)       & $47.7$  & $0.999$ & 8 days \\
    Newton ($s_N$)             & $49.8$  & $1.000$ & 4 days \\
    \hline
  \end{tabular}
\end{table}

\begin{table}[ht]
  \centering
  \caption{General relativity ($n_{\mathrm{eff}} = 8{,}859$).}
  \label{tab:gr-full}
  \medskip
  \begin{tabular}{l r r r}
    \hline
    Quantile & Score $s$ & $\widetilde{p}(s)$ &
    $\mathbb{E}[T_n]$ \\
    \hline
    Median (50th)              &  $1.0$  & $0.146$ & $10^{7410}$ years \\
    99.9th percentile          &  $6.2$  & $0.624$ & $10^{1816}$ years \\
    $B$ mode                   & $12.1$  & $0.850$ & $10^{619}$ years \\
    $\#3$ contemporaneous      & $39.0$  & $0.998$ & $48{,}000$ years \\
    $\#2$ contemporaneous      & $39.7$  & $0.998$ & $7{,}200$ years \\
    Einstein ($s_{(1)}$)       & $47.7$  & $0.999$ & 40 days \\
    Newton ($s_N$)             & $49.8$  & $1.000$ & 14 days \\
    \hline
  \end{tabular}
\end{table}

\begin{table}[ht]
  \centering
  \caption{Newtonian mechanics ($n_{\mathrm{eff}} = 11{,}453$).}
  \label{tab:newton-full}
  \medskip
  \begin{tabular}{l r r r}
    \hline
    Quantile & Score $s$ & $\widetilde{p}(s)$ &
    $\mathbb{E}[T_n]$ \\
    \hline
    Median (50th)              &  $1.0$  & $0.146$ & $10^{9581}$ years \\
    99.9th percentile          &  $6.2$  & $0.624$ & $10^{2350}$ years \\
    $B$ mode                   & $12.1$  & $0.850$ & $10^{801}$ years \\
    $\#3$ contemporaneous      & $39.0$  & $0.998$ & $10^{7}$ years \\
    $\#2$ contemporaneous      & $39.7$  & $0.998$ & $10^{6}$ years \\
    Einstein ($s_{(1)}$)       & $47.7$  & $0.999$ & 168 days \\
    Newton ($s_N$)             & $49.8$  & $1.000$ & 40 days \\
    \hline
  \end{tabular}
\end{table}

The tables reveal several striking patterns.

For Einstein, increasing the task depth from special relativity to
general relativity increases the expected time from $8$~days to
$40$~days, a factor of $5$. For the third-ranked contemporaneous
individual, the same increase raises the expected time from $86$~years
to $48{,}000$~years, a factor of roughly $560$. The depth acts as a
\emph{sensitivity amplifier}: in the linear regime it scales time
moderately, while nearer the transitional regime it scales time
catastrophically.

Both Einstein and Newton appear in every table, permitting direct
comparison. On special relativity ($n_{\mathrm{eff}} = 5{,}949$), their
expected times are $8$ versus $4$~days: both are deep in the linear
regime, and Newton's higher score yields only marginal improvement on a
relatively shallow task. On Newtonian mechanics
($n_{\mathrm{eff}} = 11{,}453$), the gap widens: Einstein would require
$168$~days (nearly six months), while Newton completes it in $40$~days.
The deeper the task, the more Newton's additional score margin matters,
precisely because the exponential barrier's curvature is steeper at
higher~$n$, amplifying even modest score differences.

The third-ranked contemporaneous individual ($\#3$, $s = 39.0$) provides
a useful benchmark for collective accessibility. For special relativity,
$\#3$ faces $86$~years, beyond a single career but within a human
lifetime. For general relativity, $\#3$ faces $48{,}000$~years, and for
Newtonian mechanics, $10^7$~years. The model predicts that special
relativity was the shallowest of the four revolutions and the closest to
the boundary of collective accessibility, consistent with the assessment
that Poincar\'{e} and Lorentz possessed many of the mathematical
ingredients \cite{pais1982subtle} but that the full ontological chain
remained out of reach within a single generation.

\medskip
\noindent\textbf{How long would the 1905 physics community have taken?}\;
The special-relativity column of Table~\ref{tab:sr-full} provides a
direct answer, with one important caveat: the times in the table are
\emph{on-task working hours}, not calendar time. Einstein's $500$ hours
for the base calibration were dispersed over approximately $9$ calendar
months, a duty cycle of roughly $7.6\%$, or about $13$ on-task hours per
week. A Poincar\'{e}-level researcher devoting $1{,}000$--$2{,}000$
on-task hours per calendar year to a specific problem would experience a
calendar-time multiplier of approximately $4$--$9\times$ relative to
the tabled values.

Applying this correction, the second-ranked contemporaneous individual
faces approximately $110$--$220$ calendar years for special relativity;
the third-ranked faces $380$--$750$ calendar years. These are well
beyond a single human lifetime. On Pais's account, the Poincar\'{e}s and Lorentzes of the era possessed the mathematical apparatus but not the specific cognitive orientation that the ontological chain demanded \cite{pais1982subtle}; the model does not establish that reading, and the absence of a visible result is weak evidence about capacity, opportunity and the direction of effort mattering as well (Section~\ref{app:screen}).
Even so, special relativity remains the shallowest of the four
revolutions: the collective of ranks $\#2$--$\#50$ has a non-negligible
probability of producing it within a few centuries, consistent with the
widespread historical judgment that the mathematical ingredients were
``in the air'' even if the full ontological synthesis was not.

The deeper revolutions tell a different story. For general relativity,
the second-ranked individual faces $7{,}200$ on-task years, or
$30{,}000$--$60{,}000$ calendar years. On the model's illustrative calibration, this revolution was \emph{not} in the air; it required an individual able to sustain an
ontological chain of $n_{\mathrm{eff}} = 8{,}859$ micro-steps without
error.

This estimate underscores a subtle but important point about
\emph{paradigmatic} advances in the Kuhnian sense \cite{kuhn1970structure},
those whose completion type has substantial depth. For such tasks, the
collective completion probability is dominated not by the \emph{number} of
researchers but by the \emph{single best-suited individual}. Adding
thousands of researchers whose scores fall below the critical threshold
produces negligible improvement, because the exponential barrier renders
their individual contributions effectively zero. The collective is only
as fast as its fastest member, a direct consequence of the
consecutive-success mechanism's non-decomposability.

It bears emphasis, however, that the extreme concentration of
achievement documented above is specific to \emph{high-depth Kuhnian
revolutions}. The same consecutive-success mechanism governs everyday
productive work; the only difference is the depth~$n$. At $n = 1$ (a
single micro-step), the $B$-mode physicist ($\widetilde{p} \approx
0.85$) is only $1.2$ times slower than Einstein ($\widetilde{p} \approx
0.999$), and the top $50$ contemporaneous individuals are all within
$1.004$ times of his speed, entirely indistinguishable in practice. Even
the population median ($\widetilde{p} \approx 0.15$) is only $6.9$ times
slower. On single-step tasks, Einstein is unremarkable among
professionals and only modestly faster than anyone in the general
population. At moderate depth (e.g., $n_{\mathrm{eff}} \approx
3{,}000$--$5{,}000$), the second- and third-ranked individuals finish
within years to decades of on-task time, well within a human lifetime,
and the collective of ranks $\#3$--$\#50$ has a high probability of
independent completion. The crossover bites only at high depth,
where the exponential barrier separates those above the crossover score
from those below it.

Table~\ref{tab:depth-comparison} makes this transition explicit, and
Figure~\ref{fig:depth-comparison} displays the same data graphically.

\begin{table}[ht]
  \centering
  \caption{Expected completion time (on-task) and ratio to Einstein
           across task depths. Reading left to right, the sharp crossover becomes visible: at $n = 1$, the $B$~mode
           physicist is $1.2\times$ Einstein; by $n = 7{,}700$, the
           ratio is $10^{538}$. Reading top to bottom at fixed~$n$,
           the gap between the record-holder and the rest widens with
           score.}
  \label{tab:depth-comparison}
  \medskip
  \small
  \begin{tabular}{l r r r r r r r r r r}
    \hline
    & \multicolumn{2}{c}{$n = 1$}
    & \multicolumn{2}{c}{$n = 10$}
    & \multicolumn{2}{c}{$n = 100$}
    & \multicolumn{2}{c}{$n = 1{,}000$}
    & \multicolumn{2}{c}{$n = 7{,}700$} \\
    Quantile ($s$)
    & $\mathbb{E}[T]$ & ratio
    & $\mathbb{E}[T]$ & ratio
    & $\mathbb{E}[T]$ & ratio
    & $\mathbb{E}[T]$ & ratio
    & $\mathbb{E}[T]$ & ratio \\
    \hline
    10th pct (0.5)
    & 3 m & 14
    & $10^5$ y & $10^{11}$
    & $10^{109}$ y & $10^{113}$
    & $10^{1142}$ y & $10^{1145}$
    & $10^{8836}$ y & $10^{8837}$ \\
    25th pct (0.7)
    & 2 m & 10
    & $10^4$ y & $10^{9}$
    & $10^{94}$ y & $10^{98}$
    & $10^{993}$ y & $10^{996}$
    & $10^{7688}$ y & $10^{7690}$ \\
    Median (1.0)
    & 2 m & 6.9
    & 124 y & $10^{7}$
    & $10^{77}$ y & $10^{82}$
    & $10^{831}$ y & $10^{834}$
    & $10^{6440}$ y & $10^{6441}$ \\
    75th pct (1.5)
    & 1 m & 4.8
    & 4 y & $10^{6}$
    & $10^{62}$ y & $10^{66}$
    & $10^{674}$ y & $10^{677}$
    & $10^{5231}$ y & $10^{5232}$ \\
    90th pct (2.1)
    & 50 s & 3.5
    & 66 d & $10^{5}$
    & $10^{48}$ y & $10^{53}$
    & $10^{539}$ y & $10^{542}$
    & $10^{4192}$ y & $10^{4194}$ \\
    99.9th pct (6.2)
    & 23 s & 1.6
    & 1 h & 30
    & $10^{15}$ y & $10^{19}$
    & $10^{199}$ y & $10^{202}$
    & $10^{1574}$ y & $10^{1575}$ \\
    $B$ mode (12.1)
    & 17 s & 1.2
    & 6 m & 2.7
    & 33 y & $10^{6}$
    & $10^{65}$ y & $10^{68}$
    & $10^{537}$ y & $10^{538}$ \\
    \hline
    $\#10{,}000$ (17.9)
    & 15 s & 1.1
    & 3 m & 1.4
    & 1 d & 75
    & $10^{22}$ y & $10^{25}$
    & $10^{200}$ y & $10^{201}$ \\
    $\#1{,}000$ (22.0)
    & 15 s & 1.0
    & 3 m & 1.2
    & 3 h & 7.4
    & $10^{9}$ y & $10^{12}$
    & $10^{103}$ y & $10^{104}$ \\
    $\#100$ (26.9)
    & 14 s & 1.0
    & 3 m & 1.1
    & 55 m & 2.2
    & 78 y & $10^{5}$
    & $10^{45}$ y & $10^{46}$ \\
    $\#50$ (28.6)
    & 14 s & 1.0
    & 3 m & 1.1
    & 44 m & 1.8
    & 3 y & $10^{4}$
    & $10^{33}$ y & $10^{35}$ \\
    $\#10$ (34.2)
    & 14 s & 1.0
    & 2 m & 1.0
    & 30 m & 1.2
    & 4 d & 16
    & $10^{11}$ y & $10^{13}$ \\
    $\#3$ (39.0)
    & 14 s & 1.0
    & 2 m & 1.0
    & 27 m & 1.1
    & 14 h & 2.7
    & $3{,}877$ y & $10^{5}$ \\
    $\#2$ (39.7)
    & 14 s & 1.0
    & 2 m & 1.0
    & 26 m & 1.1
    & 12 h & 2.3
    & 755 y & $10^{4}$ \\
    \hline
    Einstein (47.7)
    & 14 s & 1.0
    & 2 m & 1.0
    & 24 m & 1.0
    & 5 h & 1.0
    & 21 d & 1.0 \\
    \hline
  \end{tabular}
  \\[4pt]
  \raggedright\footnotesize
  Abbreviations: s = seconds, m = minutes, h = hours, d = days, y = years.
  Ratio column gives $\mathbb{E}[T]/\mathbb{E}[T_{\mathrm{Einstein}}]$.
  Contemporaneous ranks are from the representative Gumbel-process
  realization (Section~\ref{ssec:gap}).
\end{table}

\begin{figure}[ht]
  \centering
  \includegraphics[width=\textwidth]{depth_comparison.pdf}
  \caption{The crossover across task depth.
    \textbf{(a)}~Expected completion time (log-scale) as a function of
    score, for five task depths. At $n = 1$ (teal), all curves are nearly
    flat; Einstein is barely faster than anyone. At
    $n = 7{,}700$ (olive), the curve plunges through the crossover score $s_c(7700) \approx 56.7$ (red dotted line; the single-step reliability point $1/a \approx 6.4$ is not the crossover), producing a separation of thousands
    of orders of magnitude between the $B$~mode and Einstein.
    \textbf{(b)}~Slowdown relative to Einstein as a function of task
    depth, for eight score levels. The median (red) and 75th percentile
    (orange) diverge almost immediately. The $B$~mode (blue) begins
    diverging around $n \approx 30$. The top contemporaneous ranks
    ($\#3$, $\#2$; pink and yellow) remain close to Einstein until
    $n \approx 1{,}000$, then separate rapidly. At the base
    calibration $n = 7{,}700$, only the record-holder operates in the
    linear regime.}
  \label{fig:depth-comparison}
\end{figure}

This depth-dependence produces a strongly divided structure across the Gumbel-process realizations of Section~\ref{ssec:gap}. At the base calibration ($n = 7{,}700$, $p_0 = 0.1455$, $\lambda = 252$ per on-task hour), over a horizon of $1{,}000$ on-task years and with near certainty taken as $95\%$, a simulation of $200{,}000$ realizations of the top fifty scores in a population of $8 \times 10^9$ (script in the supplementary material) gives three outcomes. In about half ($52\%$), the collective of ranks $\#2$ onward completes the task with near certainty: the revolution is ``in the air,'' and any of several individuals could have produced it. In about a quarter ($24\%$), the collective's probability is at most $5\%$, so that only the record-holder completes within that horizon and the revolution appears to require singular genius. The remaining quarter is intermediate. The split moves with depth: the collective is near-certain in $78\%$ of realizations at $n = 5{,}949$ and in $21\%$ at $n = 11{,}453$, where the record-holder-only outcome rises to $62\%$. All three outcomes arise from the same mechanism; the only variable is the Gumbel draw. The model's prediction is therefore not that every
revolution requires an irreplaceable individual, but rather that the
\emph{probability} of irreplaceability is governed by the gap
distribution~\eqref{eq:gap-tail} and the task depth~$n$: shallow tasks are mostly collectively accessible, deep tasks mostly dominated by the record-holder, and the transition between these regimes is sharp.

The model thus bridges in two directions simultaneously: it unifies the
mechanism across \emph{individuals} (everyone runs the same process,
differing only in score) and across \emph{task complexity} (routine work
and revolutionary breakthroughs are governed by the same dynamics,
differing only in depth). The concentration of achievement among extreme
outliers emerges not from a different mechanism but from the same
mechanism operating at higher~$n$.

The Newtonian row is the most extreme. At $n_{\mathrm{eff}} = 11{,}453$,
even the 99.9th percentile of the general population faces
$10^{2350}$~years. Newton's
achievement was not just difficult; it was structurally inaccessible
to every other human who has ever lived, save perhaps a handful of
individuals at comparable quantiles across the full span of history.
Within the present framework, Newton's singular status in the history
of physics is not a cultural judgment but a quantitative consequence of
having completed the deepest revolution on record from a score position
that rendered it tractable.

\medskip
\noindent\textbf{A caveat on cross-domain score comparison.}\;
The scores $s_{(1)} \approx 47.7$ and $s_N \approx 49.8$ are inferred from
their respective positions in the population-wide log-normal
distribution and are not directly comparable across achievement
categories. The achievement types comprising special or general
relativity need not be identical to those comprising Newtonian mechanics:
they may differ in category $\mathcal{K}$, baseline probability $p_0$,
or both. Newton's score of $49.8$ reflects his position as the top
draw from a larger cumulative population, but whether that score
translates at full value into Einstein's specific achievement categories
depends on the correlation structure between the relevant
categories, a quantity we do not estimate here. In the homogeneous
model used throughout this section, the scores are treated as
domain-general for simplicity, but a heterogeneous treatment
(Section~\ref{ssec:hetero}) would assign each individual a
category-specific score $s_j$ for each achievement type
$\mathfrak{a}_j$. It is entirely possible that Newton's score in
certain Einsteinian categories was lower than $49.8$, or, given that
he was Newton, that it was not. The tables above should be read with
this caveat in mind: the Newton rows represent an upper bound on his
performance in Einstein's domain, not a guaranteed translation.

\subsection{The Appearance of Categorical Difference in Recreational Games}\label{ssec:games}

Allow me to quickly unpack why Elo scores \cite{elo1978rating} in chess do not follow this same sequential achievement distribution, and indeed why this is directly the point of the model in the first place. To start, the very fact that Elo scores are lognormally-distributed \cite{downey2025chess} immediately means we can categorize\footnote{Notice how the scores in chess do not contradict or are otherwise incompatible with the sequential achievement idea---instead they are simply a special case of it.} the game of chess itself as an $n=1$ completion type, because it is specifically the single-step special case of the sequential achievement model that reduces directly to a lognormal distribution in scores. This is a quantitative categorization method, but it is also perfectly deducible from the categorical content of chess itself, which is perhaps the entire point of this model. That is, the true reason chess is single-step is simply because after each and every game, the process is reset exactly, and you begin each next game from precisely the same starting point. What this amounts to is artificially enforcing the ``reset to beginning" for everyone, thereby disallowing the cascade that would highlight the world's best as categorically like Ramanujan. 

To illustrate what \textit{would} enable this cascade, we simply need to imagine the variation to the game of chess which would be the counterpart to the multi-step completion types seen in pure mathematics, (such as those with which Ramanujan was engaging). For instance, we might construct a thought experiment for a game called ``superchess", to which a single chess game is as a set is to a match in tennis, or how a game is to a series in basketball or hockey. The idea of super-chess is that a ``series" consists of thirty separate chess games, such that players always play against random other players who have completed exactly the same number of games as themselves\footnote{Or, alternatively, (since playing against others at near to the same stage of the process would probably reduce winnings to 50-50 chances rather than rapidly approaching zero probability of winning), we might also imagine for the thought experiment that there exists a second tournament which each person is playing, with the superchess being the A-series and the other being the B-series. For the B-series, we might suppose that each person plays the part of whoever each A-series person is playing, such that the B-series person always gets to keep all their pieces, and gets to play a random A-series person during each game. Then perhaps each person, for every other round (or perhaps more chaotically than that if we want to have a chance of each person playing each other person at least once), would switch off between playing in the A-series versus in the B-series, that way they can compete for the leaderboard in both. Ironically enough, while I believe the A-series would inflate scores as per the aforementioned superchess mechanism, the B-series would very likely muffle Elo-advantages, as those people who happen to get a higher number of ``easy games" again people with low numbers of pieces get inflated scores regardless of their underlying talent. So, if you will, \textit{talent counts for more in the A-series, and for less in the B-series}.}. The catch, however, is that every time someone loses a game, they must choose a piece from the board to leave out when starting all remaining games in the series. Accordingly, the people who lose several games early on will almost surely lose their remaining games, because they will be down so many pieces when beginning each of them. The idea of this is that by the end of the series, if we rank players by total number of won games, then even the most minute perturbations in their intrinsic Elo score will present itself as perhaps doubly or triply as many won games. This is precisely illustrative of where attempting to put Ramanujan at an effectively vast z-score on the lognormal talent distribution of normal people becomes a category error: \textit{what we are actually measuring with Ramanujan's output is contingent on both his edged out\footnote{Although, in fairness, I suspect that Ramanujan's intrinsic z-score itself was a lucky enough dice roll that his quantile might be a fair bit higher than expected by extrapolating from his contemporaries, as per the residual analysis. So the more accurate framing is that he edges out on the timescale of perhaps several hundred years worth of people---thereby still maintaining logical realism in his quantile---but all the while he lives amongst a population of contemporaries who themselves hold a highest score on the timescale of decades at most. Hence, he actually has a fairly impressive---just not cosmically implausible---z-score already, which is why when we deal with the full multi-step completion type inflation he appears as though plainly superhuman.} lognormal intrinsic cognition (equivalently to the Elo score of chess players), and the nature of the discipline he occupies, which has as an inherent property to it that messing up even one step along the way is catastrophic to one's series score, such that lower probability of messing up intermediate steps amounts to orders of magnitude higher score for the entire series}. What we call ``Ramanujan's cognition ability", broadly speaking, is therefore a conjunction of both his Elo-analogous quality with the contingency in world class mathematics that you are always measuring cognition for the series, not the games\footnote{Although, for sub-world class individuals we often measure in terms of the games, or at least significantly lower-step-count series. This perhaps adds to the illusion that Ramanujan is categorically different, because we are inadvertently trying to stitch-together the distribution of score on lower-step series with his immense score on higher-step series, thereby appearing to inflate his z-score.}. Accordingly, both cases are indeed perfectly analogous, and Ramanujan still has a non-antirealist z-score\footnote{This is also the precise point at which the framework answers the unease Bessis \cite{bessis2024mathematica,bessis2025genes,bessis2026jaimungal} voices---that Ramanujan \emph{looks} categorically alien while one feels strongly that he ``should not be.'' The resolution vindicates the intuition rather than the appearance: the error, on this account, lies not in the log-normal model of intrinsic cognition, which the paper adopts as a hypothesis, but the extrapolation of game-outcome statistics to series outcomes---reading the inflated series output back as though it fixed an intrinsic $z$-score (in effect, ``extrapolating from typical individuals, at what $z$-score does this output sit?''). Undoing that step leaves a realist picture in which his intrinsic standing is reasonable and merely his \emph{output} is stretched by series depth, which is exactly the anti-alien intuition made precise. The measurement appendix (Section~\ref{IQ_individual}) supplies a second, independent wedge of the same kind on a different axis: even an observed psychometric score decouples from latent capacity in the extreme tail, so the gap between what is measured and what is intrinsically present opens there too.}---it just so happens that the analogous ``series probability" pertaining to his category of work stretches out the distribution.

\section{Conclusion}\label{sec:conclusion}

We opened with a question about Ramanujan, dying at thirty-two in a
nursing home in Madras, having produced a body of work that the
collective mathematical community has spent a century merely
unpacking. The framework developed in this paper offers a conditional, quantitative answer to the question ``how does one account for this?'', under its stated assumptions, in three interlocking parts.

\emph{First}, the proposed extreme-value diagnostic, the proportional scaling of rank-$\#1$ residuals above the elite subpopulation, motivates a log-normal working model of latent ability relative to the comparator families considered, through an illustrative inference scheme (Section~\ref{sec:inference}). The Gumbel process governing extreme
order statistics then produces, as a routine consequence of the
exponential spacing formula~\eqref{eq:score-gap}, a natural gap between
the record-holder and the field: a median gap of $6.3\%$, with a $34\%$
probability of exceeding $10\%$. This gap requires no special
mechanism; it is the generic statistics of the maximum of a large
sample from a log-normal population.

\emph{Second}, the consecutive-success mechanism
(Section~\ref{sec:completion}) converts this modest gap in latent
score into an exponential gap in realized output. The phase
crossover at the depth-dependent score $s_c(n)$ acts as a threshold:
individuals above it experience the task as routine (linear in
depth), while individuals below it face, under full restart, an exponential barrier that renders the task practically infeasible within civilizational time budgets, the finite-horizon probability of \eqref{eq:horizon} being negligible there. Table~\ref{tab:depth-comparison}
demonstrates this transition quantitatively: the $B$-mode
theoretical physicist, who is only $1.2$ times slower than Einstein
on a single micro-step, faces $10^{537}$~years on a foundational
advance of depth $7{,}700$. The mechanism is universal: everyone
runs the same process, differing only in score. Nobody is
categorically excluded; everyone below the threshold is
\emph{effectively} excluded by the mathematics of exponentiation.

\emph{Third}, the interaction between these two features produces
a bimodal prediction about scientific revolutions. In approximately
half of all Gumbel-process realizations, the gap between the
record-holder and the second-ranked individual is small enough that
the collective can produce the result independently: the revolution
is ``in the air.'' In the remaining realizations, only the
record-holder can complete the task within a civilizational
timescale, and the loss of that individual represents a genuine,
quantifiable deficit to the collective record. Which regime
obtains in a given historical instance depends on two quantities:
the depth of the required achievement chain, and the Gumbel draw.
Shallow tasks are almost always collectively accessible; deep tasks
are almost always dominated by the record-holder; and the
transition between these regimes is sharp
(Figure~\ref{fig:depth-comparison}).

Under the model's assumptions, this offers a resolution of the tension identified in the opening: under full restart, modest differences in individual reliability can produce very large differences in completion time, including parameter regimes in which one person's completion is feasible while the remaining population's is unlikely within a specified horizon. Ramanujan is
not a categorical alien: he is a draw from the same log-normal
distribution as everyone else, located at an extreme but finite
quantile. And the strong inevitability thesis is not wrong in
general: for shallow achievements, the collective \emph{would} have
produced the same results eventually, and does so routinely. What the inevitability thesis misses, on the model and under full restart, is that for deep achievements the consecutive-success barrier converts continuous differences in
ability into discrete-looking differences in output, so that, within civilizational time budgets, the task is feasible for the record-holder and practically infeasible for everyone else: a difference of degree, made enormous by exponential structure.

The framework makes several testable predictions. The distribution of completion rates across individuals within a field should be approximately log-normal, with log-dispersion of order $n\sigma$ in the low-score regime, and should saturate at the rate ceiling $\lambda/n$ at the top: it may resemble a power law over a limited window, but its curvature and saturation distinguish the two. The probability that a
revolution is ``in the air'' (i.e., that the collective could
have produced it independently within a generation) should decrease
with the depth of the required achievement chain, in a manner
quantified by~\eqref{eq:gap-tail}. And progress on deep open
problems, when it comes, would on the model tend to be sudden rather than incremental and to concentrate in individual work rather than diffuse collective convergence, not because of any mystical property of the individual but because the mathematics of consecutive success produces that tendency under its assumptions.

Ramanujan, in this reading, was not a miracle in the sense of requiring explanation beyond the human distribution---but he was extraordinary in the precise statistical sense that his position in the tail of that distribution is itself a rare event whose consequences are amplified beyond recognition by the consecutive-success mechanism. A population of billions, distributed log-normally in latent ability, produces individuals at extreme quantiles, and for tasks of suitable depth their consecutive-success probability can place them in the linear regime while the task remains exponentially infeasible for everyone else; for deeper tasks it remains infeasible for all. On the model, that such an individual appeared in early twentieth-century India, without formal training, reflects the universality of the distribution rather than the exceptionality of the individual. The mechanism that produced Ramanujan is the same mechanism that produces a median mathematician; it is operating at the same score, drawn from the same family, processed through the same consecutive-success dynamics. The only difference is where in the distribution the draw landed, and the staggering amplification that the task depth applies to that difference. The wonder is real; what is removed is only the need to place its source outside the framework that governs everyone else.

The divergence interpretation of Section~\ref{ssec:divergence} provides a final reframing. The standard question about Ramanujan---\emph{what special capacity did he possess?}---is replaced by its inverse: \emph{what common impediment did he lack?} The answer this paper proposes is divergence friction: the spontaneous tendency, at each step of an inferential chain, to extrapolate the structure of the problem incorrectly. In this reading, which is a hypothesis about cognition, Ramanujan's genius would consist not in possessing something extraordinary but in lacking something ordinary. The mechanism that produced his output is the same mechanism that produces a median mathematician's output; what differs is the divergence rate, and the consecutive-success process converts even modest differences in divergence rate into differences of many orders of magnitude in realized output. On the divergence hypothesis, the question ``what is Ramanujan?'' thus receives a precise answer: he is an individual whose divergence rate places him in the linear regime of the consecutive-success process for tasks whose depth renders them exponentially infeasible for everyone else. He is not a mutant, and not a messenger of the gods. He is a draw from the tail of a universal distribution, processed through a universal mechanism, whose output is extraordinary not because the input is categorically different but because the mathematics of consecutive success is unforgiving.

A closing remark on the nature of the explanation. Bessis \cite{bessis2024mathematica,bessis2025genes,bessis2026jaimungal} has argued that mathematical talent is best understood not through the lens of genetic endowment but through the lens of cognitive development: the brain's capacity to build, through practice and neuroplasticity, ever-richer intuitive representations. The present framework is fully compatible with this view. The score~$s$ is not identified with any specific biological substrate; it is a sufficient statistic encoding the net probabilistic consequence of whatever combination of genetic predisposition, developmental history, and cognitive practice produced the individual's current capacity. Whether a high score arises primarily from genetic fortune, from years of deliberate practice, from an unusual developmental trajectory, or---most likely---from some irreducible combination of all three, the consecutive-success mechanism operates identically. The model is therefore agnostic on the \emph{deepest} source of the score---the local neurological or developmental mechanisms that produce a given individual's divergence rate---and specific only about its \emph{consequences} under the model's assumptions: the crossover, the inaccessibility below it under full restart, and the depth-proportional amplification of dispersion. Between these two levels, however, the divergence interpretation of Section~\ref{ssec:divergence} occupies a realist middle ground. While we do not claim to know \emph{why} some individuals carry less divergence friction than others, we propose, as a hypothesis the model does not itself establish, that what the score measures, at the level of the inferential process itself, is divergence resistance: the degree to which the individual's extrapolation of structure at each step is free from the systematic deflections that reset the chain for typical individuals. This is not merely a statistical label but a characterization of the cognitive event that the per-step probability encodes.

The model thus operates at three levels of specificity: the statistical mechanism is fully specified, the inferential interpretation identifies what the mechanism is tracking, and the ultimate cognitive source remains open. It is at the middle level that the model's distinctive claim resides---that extreme outlier productivity is better quantitatively understood as the absence of a common impediment than as the presence of a rare gift---and it is this claim, rather than any commitment about neurology or genetics, that the framework offers as its contribution to the question of what Ramanujan was.

\printbibliography

\appendix

\section{Measurement, Decorrelation, and Population Structure at the Extreme Tail of Ability}
\label{IQ_individual}

This appendix develops the statistical structure governing measured ability at the extreme upper tail and the consequent misattribution of exceptional achievement to extreme scores. The argument proceeds from the behavior of measurement error in the tail, through an account of ability decorrelation as a measurement artifact, to a conditional model of outlier capacity and its implications for the distribution of ability \emph{among} transformative individuals.

\medskip
\noindent\emph{Scope and notation.}\;
Symbols in this appendix are local to it. In particular, $X$ denotes an \emph{observed} psychometric test score (IQ-type, approximately Gaussian by construction of the norming) and $\theta$ its latent true-score counterpart in the classical-test-theory sense $X = \theta + E$; these are distinct from the main text's $X$ (the latent log-normal ability trait) and $T_n$ (completion time). Conceptually, the observed score $X$ is an imperfect, ceiling- and validity-limited proxy for latent capacity, and the measurement pathologies developed below are precisely what cause $X$ to \emph{decouple} from latent capacity---and hence from the divergence-resistance score $s$ of Section~\ref{ssec:divergence}---in the extreme upper tail. The appendix is therefore complementary to the main development rather than in tension with it: where the main text shows how a modest latent advantage is amplified into vast realized output, this appendix shows why the \emph{observable proxy} for that advantage becomes uninformative exactly where the amplification is largest.

\subsection{Adjacent constructs and the conditional model}\label{app:conditional}
Within-person scatter across conceptually adjacent abilities---for example, pre-linguistic mathematical reasoning versus symbolic--procedural reasoning---is routine. Indeed, the very qualities that make an individual exceptional within a specialized subpopulation can, at ordinary score levels, be nearly interchangeable with a single coarse competence: pre-linguistic mathematical reasoning and generic ``math skill'' are scarcely distinguished across the typical population, and resolve into separate coordinates only among the rare individuals for whom one of them is load-bearing and exceptional. Within-person scatter is, on this view, what a coarse measurement scheme produces when it samples a distinction it does not itself resolve---the special case, at the level of a single person, of the category-transfer mechanism of Section~\ref{app:singularity}. The operative question is the conditional prospect of exceeding a high threshold $a$ on a measure $U$ given a score $b$ on a related measure $V$. Modeling $(U,V)$ as bivariate normal in standardized units with correlation $\rho$,
\[
P(U>a\mid V=b)=1-\Phi\!\Big(\tfrac{a-\rho b}{\sqrt{1-\rho^{2}}}\Big),
\]
the conditional mean $\rho b$ exhibits regression toward the mean and is the dominant reason extreme scores rarely co-occur. Under this model alone $P(U>a\mid V=b)$ is increasing in $b$ whenever $\rho > 0$, so no plateau or inversion follows from it; the plateau below requires the additional structure of a latent gate. The correlation $\rho$ conflates two quantities---true-construct overlap and reliability---separated by the correction for attenuation \cite{spearman1904}, $\rho_{\mathrm{obs}}=\rho_{\mathrm{true}}\sqrt{r_U r_V}$. Setting $\rho_{\mathrm{true}}=1$ isolates the divergence attributable to measurement alone: $U-V\sim\mathcal N\!\big(0,\,\mathrm{SEM}_U^{2}+\mathrm{SEM}_V^{2}\big)$, whose standard deviation is the standard error of the difference score.

\subsection{Tail amplification of measurement imprecision}\label{app:tail}
A fixed percentile-level error is negligible near the median but severe in the tail, through two independent channels. First, quantile-function nonlinearity: $\tfrac{d}{dp}\Phi^{-1}(p)=1/\varphi\!\big(\Phi^{-1}(p)\big)$ diverges as $\varphi\to 0$, so a fixed percentile displacement maps to an increasingly large score gap---on the order of a $70\times$ amplification from the median to the $99.9$th percentile. Second, under item response theory \cite{lord1980} the conditional standard error of measurement is the inverse square root of the test information function, which collapses at the extremes for want of sufficiently discriminating items; measurement error in ability units therefore grows in the tail, compounding the first effect. Jointly these render score-level imprecision practically decisive between roughly the $98$th and $99$th percentiles---that is, before the construct-validity failures of Section~\ref{app:three} even onset.

\subsection{Three sources of divergence at the tail}\label{app:three}
For an individual far out on a measured construct (``Ramanujan-tier'' in realized output, which the main text's series-versus-game distinction (Section~\ref{ssec:games}) shows need not correspond to a cosmically extreme latent quantile), observed scatter decomposes into three mechanisms with distinct distributions and distinct implications.

\emph{Profile jaggedness.} Abilities decorrelate at the top (Spearman's law of diminishing returns \cite{spearman1927}; the ability-differentiation hypothesis \cite{tuckerdrob2009}), so different instruments sample genuinely different coordinates of a jagged profile; the individual does not possess a single scalar ``IQ'' but a collection of loosely coupled peaks. This is the measurement-side counterpart of the class-specificity of the score in the main text (Sections~\ref{ssec:divergence} and~\ref{ssec:hetero}): a score is well-defined only within an achievement class, and cross-class comparison is not---though Section~\ref{app:singularity} will argue that even this jaggedness is largely an artifact of the standard metric's category scheme, so that the seeming divergence between a high scorer's different-category scores is a property of the instrument rather than a verdict on anything intrinsic to the person.

\emph{Ceiling artifact.} Standard instruments operate where their information function is near zero for such a person. Observed scores are biased downward toward the instrument ceiling, and rank ordering among extreme individuals is dominated by noise.

\emph{Validity failure and non-monotonicity.} At extreme ability the test increasingly measures something other than the target construct---breadth, speed, adherence to the expected solution path---so the map from true ability to observed score can flatten and locally invert. A depth specialist whose cognition operates in a non-standard mode may score \emph{below} a well-rounded but less exceptional individual. This is systematic misfit, not Gaussian noise: a $40$-point negative gap is $8$--$13$ standard errors of measurement and is effectively never realized as random error, so when such a gap appears its source is construct misfit rather than an unlucky draw. Only this mechanism can generate gaps of that magnitude; random measurement error and systematic validity failure behave oppositely in the tail, the former thin-tailed and bounded in its reach, the latter fat-tailed and able to displace a score arbitrarily far. Administering many instruments compounds the effect: the observed minimum across a battery is an order statistic, so even a uniformly capable extreme individual will, across enough tests, occasionally post a strikingly low score wherever some instrument samples a weak coordinate or rewards the wrong kind of breadth---which makes a handful of scores well below the rest expected rather than anomalous.

\subsection{Decorrelation as instrument-relative}\label{app:singularity}
The central claim is that tail decorrelation of abilities is instrument-relative rather than an intrinsic feature of high-ability individuals.

\emph{Measurement invariance.} Factor structure ($g$ and subtest loadings) is extracted from a reference population; factorial invariance \cite{meredith1993} across subpopulations is not guaranteed, and scoring an atypical group with the reference loadings manufactures spurious decorrelation that presents as diminishing returns.

\emph{Frame-dependence.} The decorrelation is a real feature of the measured structure, but plausibly a removable one: a change of frame---new instruments and a refitted category scheme---would dissolve it, while simultaneously rendering the reference population's structure indistinct under the new frame. What obscures this is the absence of tail-dense data with which to construct and validate the alternative categories.

\emph{Instrument- versus coordinate-relativity.} A mere rotation of the existing factor solution does not dissolve the decorrelation, since the eigenvalue spectrum and the first-factor variance share are rotation-invariant for a fixed battery. Dissolution requires a \emph{different} battery---instruments that discriminate in the tail and load on the operative faculties---that is, new measurement rather than reanalysis.

\emph{Diagnostic.} What would settle whether the decorrelation is frame-dependent or genuine is a frame-invariant characterization of ability structure. The practical handle is convergence across independently constructed batteries: decorrelation that replicates across batteries built on divergent theoretical commitments behaves as a genuine invariant; decorrelation that smooths out under an alternative battery is a frame artifact.

\emph{Normal subpopulations.} The motivation for positing such subpopulations is a prior question: why should the total population be privileged as the only aggregate constructible as satisfying the continuity criterion at issue? Consider, then, subpopulations defined ontologically by a cognitive feature rather than by a score, which downstream happen to score at or above (say) the $98$th percentile on \emph{some} instrument. Call such a subpopulation \emph{normal}---by loose analogy with a normal subgroup, coherence preserved under the ambient structure---if it re-coheres under its own refitted category scheme, with groupings substantially rearranged relative to the reference scheme---a rearrangement that may run all the way down to the reassignment of individual people. The full human population satisfies the underlying continuity criterion naturally: being a complete, exhaustive aggregate, between any two of its members one may regard arbitrary intermediate interpolations as available---interpolations with respect to the equivalence classes induced by the scoring criteria, or by any labeling scheme deductively interchangeable with them up to mutual redescription.\footnote{The relevant continuity is not tied to the IQ scale specifically but to any relabeling that preserves the same equivalence structure; the criterion is meant to be invariant under such redescriptions, which is precisely what would let it be stated without privileging a particular instrument. It is this invariance that the full population enjoys for free and that an artificial delimitation must earn.} An artificially delimited subpopulation enjoys no such guarantee; there the burden falls on us to ensure that the delimitation introduces no discontinuous jumps between segments of the wider population. If non-trivial normal subpopulations tile the fringe of the reference population, each locally coherent under its own scheme, then the reference population's tail decorrelation is frame-relative---a mosaic of locally coherent descriptions---and there is no principled basis for privileging the reference population as ontologically surviving decorrelation while these subpopulations do not. Put structurally, this is the abstract isomorphism the construction is reaching for: each such subpopulation can stand in for the whole, its own fringe playing the role that the full population's fringe plays in the standard account. This is offered as structural motivation, not a deductive result, and its limits are explicit: coherence is not guaranteed to be closed under taking fringes, so the recursion to ``fringes of fringes'' is not secured; the predicate of ontological cohesiveness requires further specification; and the argument supports the coordinate-dependence of the decorrelation rather than its vanishing. Its force is that the standard framing admits completion---it is not synonymous with the exhaustive set of options.

\emph{Category transfer and the source of decorrelation.} The refitted scheme is, in all likelihood, not derivable from the reference scheme---in that predictive sense largely \emph{irrelevant} to it---yet not independent of it either: its categories are constituted as a hard-to-predict, variable blend of the reference categories, a blend that is itself not a tractable function of them. Statistical independence is thus the wrong idealization; crossovers are inevitable, but the re-mapping is scrambled enough that the reference scheme affords no purchase on predicting the refitted one. Encountered without knowledge of the embedding, such a subpopulation would look like an unremarkable population in its own right, complete with its own recursively definable sub-structure; nothing internal to it would betray that it is a fringe carved from a larger whole. The two schemes remain mutually expressible---transferable up to redescription---but only through this scrambled re-mapping, and that is the crux: applying the reference scheme to a subpopulation whose natural coordinates are a scrambled blend of it is exactly what manufactures the observed tail decorrelation. On this reading the loss of correlation is not a fact about the people but an artifact of coordinate misalignment---measuring one population in another's categories---which is the mechanism the instrument-relativity diagnosis has been pointing at all along.

Made concrete, this is what becomes of a category like ``math skill.'' Tailored to the general population, the math items of a standard test are, for the relevant normal subpopulation, very nearly beside the point: ontologically there is no longer a single ``math skill'' but a different package altogether---whether one is a pre-linguistic thinker, a symbolic thinker, a verbal thinker, and so on---each a distinct flavor of the same underlying gift for high-sequential achievement. A person at, say, $170$ tops out the math criterion contingently, because the profile that happens to satisfy that test's question-topology happens to be theirs; but the topology of the subpopulation in which the high-sequential-achievement individuals actually live is organized around quite different ``spindles,'' dissimilar to the general population's ``good at math,'' and it is those spindles, not the test's category, that carry the correlation with the intrinsic capacity. And the conclusion follows from our ignorance rather than in spite of it: \emph{because} we already know we lack the subpopulation's categories, we know that the test we do have---extrapolated blind from the general population's scheme---amounts to a dice roll against the subpopulation's actual coordinates, and a blind extrapolation is overwhelmingly unlikely to land on those coordinates by chance. We can therefore be confident, without ever identifying the right categories, that the test almost surely fails to track whatever would predict the outlier---so correlation degrades into the tail no later than the vicinity of the $98$th percentile of the general population's roughly-closest score.

\subsection{An exemplar of dissolution: the chess case}\label{app:cam}
The dissolution just described is not a conjecture about tails; it has a canonical exemplar in the expertise literature, and the exemplar shows the whole mechanism at once. Chess masters recall real board positions vastly better than novices, and random positions only slightly better, the advantage shrinking sharply but persisting \cite{chase1973perception, gobet1996templates, gobet1996recall}. The capacity that makes the master is therefore largely invisible to the general-population instrument that would seem to measure it: a span test presents arbitrary items, and the master's advantage lives almost entirely in retrieval structures keyed on the game's structure, not in the volatile buffer the span test samples \cite{ericsson1995ltwm}. Read as a measurement fact, this is exactly the appendix's claim in miniature. The reference scheme's category---``memory,'' operationalized as span---and the operative category in the expert subpopulation---structure-keyed retrieval---are different coordinates, nearly interchangeable at ordinary levels and orthogonal at the top; a battery built on the first will rank the master unremarkably while a refitted battery built on the second would rank him at the ceiling; and the ranking inversion is not noise but the instrument sampling a distinction it does not resolve. The chess case is the diagnostic unification because it displays, in one setting with a clean manipulation (real versus random boards), the three things the appendix asserts about the tail in general: that categories which coincide in the general population come apart in the elite subpopulation, that the general-population instrument then loads on the wrong one, and that the resulting rank order can be wrong in a way no rescoring of the reference battery repairs---only new measurement, keyed to the operative structure, would. What the master shows for one domain, the normal-subpopulation thesis asserts for the tail as such.

\subsection{A conditional model: substitution under measurement convolution}\label{app:model}
Let outlier capacity aggregate two latent inputs---general ability $\theta$ (the axis the test loads on) and a non-$g$ structural ingredient $C$ (depth and originality; the faculty distinguishing a transformative mind\footnote{``Faculty'' here is not meant to reintroduce the possessed-organ picture that Section~\ref{ssec:divergence} argues against. It may be read deflationarily as simply \emph{a high score under the relevant achievement type}---equivalently, low divergence friction within that class---so that $C$ names whatever the observed test fails to load on, not a metaphysically special endowment. Nothing in the model turns on the stronger reading.})---through a constant-elasticity-of-substitution form \cite{arrow1961} taken over positive-valued versions of the inputs, $a = g(\theta)$ and $c = g(C)$ with $g$ positive and increasing, $g = \exp$ for instance, since fractional powers of a Gaussian $\theta$ are undefined on its negative support: \[ Z=\big[w\,a^{\beta}+(1-w)\,c^{\beta}\big]^{1/\beta},\qquad \sigma_{\mathrm{CES}}=\tfrac{1}{1-\beta}, \]
whose elasticity interpolates between a Leontief threshold ($\sigma_{\mathrm{CES}}\to 0$; a hard gate requiring both inputs) and Cobb--Douglas multiplicativity ($\sigma_{\mathrm{CES}}\to 1$; smooth substitution). ``Ramanujan-tier'' is the upper tail of $Z$; the bend in capacity across the $120$--$130$ band is the rounded CES corner, its sharpness governed by $\sigma_{\mathrm{CES}}$---a free shape parameter rather than one fixed by fiat.

Only $X=\theta+E$ is observed. Under joint normality the true-score posterior is Kelley's \cite{kelley1927},
\[
\theta\mid X=x\;\sim\;\mathcal N\!\big(\mu+\rho_{xx}(x-\mu),\,\sigma_{\ast}^{2}\big),
\qquad \sigma_{\ast}=\sigma_X\sqrt{\rho_{xx}(1-\rho_{xx})},
\]
so the probability of clearing an ability gate at $\gamma$ is the probit
\[
P(\theta>\gamma\mid X=x)=\Phi\!\Big(\tfrac{\mu+\rho_{xx}(x-\mu)-\gamma}{\sigma_{\ast}}\Big).
\]
This is the operative interpolation: a normal CDF in the observed score, obtained as the exact image of the latent gate---however sharp---under Gaussian measurement error, with location and width fixed by the gate position and the test's reliability rather than chosen. The gate location $\gamma$ has one empirical anchor, weaker than an achievement threshold would be: the breakpoint analysis of \cite{jauk2013} finds thresholds for creative potential, the highest near $120$, and none for creative achievement, so a gate near $120$ is motivated for potential only and its use for achievement is an assumption. That location is distinct from the psychometric information knee in the $140$s, so the two decay phenomena---substantive gating and information collapse---occur at different locations and should not be conflated.

In the separable hard-gate case, the Leontief limit with $\theta$ and $C$ independent, the full conditional factorizes as this probit times the $C$-contribution. For general $\sigma_{\mathrm{CES}}$, or with $\rho_{\theta C} \neq 0$, it does not factorize, and the conditional is the CES upper-tail probability integrated against the Kelley posterior and the conditional law of $C$ given $\theta$; the saturation statement that follows belongs to the hard-gate case. In the hard-gate case with $C$ independent of measured ability the curve saturates sharply above the gate. Independence alone does not give that in the general CES model, since a higher $\theta$ lowers the $C$ needed to clear a fixed threshold: in the independent Cobb--Douglas case $Z = a^{w} c^{1-w}$, the condition $Z > z_{\star}$ requires $c > z_{\star}^{1/(1-w)}\,a^{-w/(1-w)}$, a requirement that falls as $a$ grows, so a plateau there must be assumed or obtained from a suitable parameter regime. Within the hard-gate case, a mildly positive $\rho_{\theta C}$ instead yields a shallow continued slope above threshold, the binding constraint sliding from $\theta$ toward $C$---consistent with longitudinal evidence (e.g.\ the Study of Mathematically Precocious Youth \cite{lubinski2006}) that returns persist within the top $1\%$ while diminishing. The residual slope above the gate is governed by $\rho_{\theta C}$; the sharpness of the bend by $\sigma_{\mathrm{CES}}$ and $\sigma_{\ast}$ jointly. A minute prior on transformative status does not by itself make the conditional probability small at every $x$, since that depends on how sharply the score separates the classes; where it is small, the informative quantity is the evidence the score carries. By Bayes's theorem, with $\pi(x) = P(\mathcal{R} \mid X = x)$, the likelihood ratio is $f(x\mid \mathcal{R})/f(x\mid\neg \mathcal{R}) = \big[\pi(x)/(1-\pi(x))\big]\,\big[(1-P(\mathcal{R}))/P(\mathcal{R})\big]$: if $\pi$ approaches a plateau $c > P(\mathcal{R})$, it approaches a constant greater than one, and while $\pi$ increases it increases. What may approach zero past the gate is the incremental change in log evidence, $\tfrac{d}{dx}\ln\big(\pi(x)/(1-\pi(x))\big)$: further increments of score become uninformative while the score already attained remains positive evidence.

Two consequences constrain interpretation. First, the transition cannot be only a few points wide: the $10$--$90$ width of the probit alone is $\approx 2.56\,\sigma_X\sqrt{(1-\rho_{xx})/\rho_{xx}}\approx 9$ points at high reliability ($\sigma_X=15,\ \rho_{xx}=0.95$), and realistically $12$--$18$ once index-level reliability and CES corner-rounding are included; a stated ``$120$ to $125$'' is therefore the \emph{center} of the transition, not its width. Second, on the $120$--$130$ interval the early portion carries little cumulative mass---an extreme individual observed at $125$ sits below roughly the $25$th percentile of that interval's mass---with the conditional approaching, but not reaching, its uniform asymptote by $130$.

The model is offered as structure rather than as a fitted object: most of its parameters could in principle be pinned from ordinary-range psychometric and creativity--ability data, but the gate location $\gamma$ and the elasticity $\sigma_{\mathrm{CES}}$ would require tail-dense cohorts that do not exist, so the curve is best read qualitatively rather than estimated.

\subsection{Population structure among extreme individuals}\label{app:population}
The conditional probability that a person at a given observed score is transformative, $\pi(x):=P(\mathcal{R}\mid X=x)$, must be distinguished from the distribution of the score \emph{among} transformative individuals, the two related by Bayes's theorem,
\[
P(X=x\mid \mathcal{R})\;\propto\;\pi(x)\,f(x),
\]
with $f$ the population density of $X$.

\begin{proposition}\label{app:prop-mode}
Let $\pi$ be bounded ($\pi \le 1$), nondecreasing, and of bounded logarithmic derivative on $[x_0, \infty)$, and let $f$ have a super-exponentially decaying upper tail with $-(\ln f)' \to \infty$ (e.g.\ Gaussian, including a log-normal variable expressed in its log-scale score $(\ln x - \mu)/\sigma$; a log-normal density on the original or an affinely standardized scale does not qualify, since there $-(\ln f)' \to 0$). Then $\pi(x) f(x)$ is eventually decreasing and integrable, so the score distribution among transformative individuals has a finite mode. If moreover $\pi f$ is nondecreasing on $(-\infty, x^{\star}]$, $\pi$ is constant beyond $x^{\star}$, and $f$ is nonincreasing beyond $x^{\star}$, the mode lies at $x^{\star}$. Boundedness of $\pi$ alone does not give eventual monotonicity, since a bounded $\pi$ may oscillate or spike; and constancy of $\pi$ beyond $x^{\star}$ makes the product a constant multiple of $f$ there, which decreases after $x^{\star}$ only where $f$ does (a Gaussian $f$ with constant $\pi$ and $x^{\star} = -1$ meets every other condition, yet its mode is at $0$), so the global mode is at $x^{\star}$ only under both monotonicity conditions.
\end{proposition}

\noindent Under the proposition's derivative conditions, the density among transformative individuals eventually decreases, so its mode is finite. A plateau in the conditional probability does not by itself locate the mode; that location also depends on the conditional's rise before the plateau and on the population density, the mode lying at the plateau onset under the proposition's further monotonicity conditions, and the plateau itself being an assumption. Equivalently, the \emph{absolute count} in a score band is $\pi(x)\,N(x)$; since $N$ collapses in the tail, under a plateau, a band near its onset contains far more transformative minds than a band deep in the tail, not because high ability is detrimental but because the population there is vastly larger and each member is nearly equally likely to qualify.

Two cautions fix the reading, and together they invert an expectation that is easy to carry in by default. First, the score here is strictly the measure-side quantity---the thing actually measured, and the thing to which assessments of standing almost exclusively refer---whereas the criterion of transformative status established earlier in the paper is abstract-side; the two are distinct objects, and the whole force of the preceding sections is that they come apart in the upper tail. Second, once they have come apart---and on the measurement analysis's model this decorrelation sets in abruptly, plausibly no earlier than about $125$---each further increment of measured score, precisely by thinning the population it labels, \emph{lowers} the expected number of criterion-satisfiers that population contains rather than raising it: above the decorrelation onset the per-person probability is essentially flat, so the count is governed by $N(x)$, which is collapsing. This is the population-weighting principle the main text invokes for reference-class size (membership in any small, sharply selected subpopulation counts \emph{against} the prior odds of harboring the next outlier, simply by scarcity of population), now in its least intuitive form: the rarer the score, the fewer transformative individuals the population at that score can hold, so that the next such individual is far likelier to be found in the populous middle than in the sparse extreme---the reverse of the reading on which a more extreme score ought to concentrate transformative standing.

A documented modest observed score for a transformative figure functions, under the equal-population idealization, as an anchor for the conditional curve---and is consistent with the main text's reading that such a figure's \emph{intrinsic} quantile is impressive but not cosmically extreme---whose parameters are then carried to the population-weighted regime.

\subsection{The criterion is abstract; visible work is an asymmetric proxy}\label{app:screen}
Transformative status is best located not in the observed score and not in the visible record as such, but in the abstract quantity the main text models directly---the score and sequential-completion capacity of Sections~\ref{ssec:divergence} and~\ref{ssec:process}, the abstract-side object of which the observed IQ score is only a measure-side proxy. A body of work $W$ is then evidence about that criterion, but evidence of a sharply \emph{asymmetric} kind, and it is worth being precise about the asymmetry rather than treating $W$ as a sufficient statistic for the status.

A visible record of revolutionary work is strong positive evidence: such work is extraordinarily improbable absent the criterion, so its likelihood ratio is large. The \emph{absence} of such a record, by contrast, is close to uninformative. Ramanujan-tier individuals spend long stretches---sometimes their entire recorded lives---without visible output, and the absence is heavily opportunity-gated: denial of education, audience, or institutional access depresses the probability of visible work for reasons orthogonal to the criterion, as the paper stresses elsewhere in connection with the systematic exclusion of women and members of marginalized groups. Conditioning on ``no revolutionary work yet'' therefore barely updates against the criterion. The evidence turns sharply negative in only one case: the \emph{demonstrated} production of low-quality work for reasons internal to the person---failures of rigor, discipline, or intellectual honesty---is strong evidence against, because the abstract criterion, low divergence friction sustained across long inferential chains, is close to incompatible with reliably generating such work. In short, brilliant work counts heavily for, characteristically shoddy work counts heavily against, and the mere absence of work counts for very little either way.

The observed IQ score sits on the other axis entirely. It is the measure-side quantity---the thing actually measured, and the thing to which claims of standing almost always refer---and the preceding sections show that it comes apart from the abstract criterion in the extreme tail, so that neither a high nor a low score is decisive there: a claimed extreme score is not evidence of meeting the criterion, and, symmetrically, an ordinary score does not refute it. What distinguishes the outlier is the abstract capacity, of which realized work is the better (if asymmetric and opportunity-gated) proxy, and of which the measured score is a proxy that fails exactly where it would be needed.

\subsection{Against monotonic scaling}\label{app:monotone}
Claims that intellectual eminence is a monotonically increasing function of \emph{measured} ability rest on three premises: (i)~eminence scales monotonically with the observed test score; (ii)~a higher observed score therefore strictly dominates in expected contribution; (iii)~figures at a moderate observed score are categorically inferior to those at an extreme one. The form of the argument is valid, so the contest is over soundness, and the load-bearing premise is~(i). The preceding analysis denies it for the \emph{observed} score: the conditional $\pi(x)$ plateaus, the marker of transformative status is the abstract criterion rather than the measured score (Section~\ref{app:screen}), and population weighting (Section~\ref{app:population}, Proposition~\ref{app:prop-mode}) places, under the proposition's conditions, the modal transformative observed score near the plateau onset, so that the observed score ceases to track eminence exactly where the claim needs it to. It is essential that this is a statement about the observed \emph{proxy}, not about latent capacity: the main text's amplification result---in which realized output rises sharply, indeed super-linearly, in the latent score $s$---is untouched, and is in fact the reason the effect matters. What this appendix adds is the complementary half: the ceiling, validity, and decorrelation pathologies of the upper tail sever the observed score from that latent capacity precisely where the amplification is largest, so that a large measured score is neither necessary nor sufficient for transformative output. Under the corrected reading the plateau onset is not a floor but the mode of the observed score among transformative figures; premise~(i) fails, and with it the inference from measured score to rank---while the main text's inference from \emph{latent} score to realized output stands.

A final observation sharpens the same point from the other side. The very individual who tops a standard test---well-rounded at precisely the competences the general population implicitly prizes, contingent as that valuation is---is no more legible, as a candidate for transformative work, than anyone else; the high score certifies fit to the test's topology rather than possession of the unpredictable capacity. And that capacity is unpredictable in a strong sense: it is recognized only once the work is in hand, never before it, and never from any standardized metric---however much a discerning human may sometimes sense it in another through particular acquaintance. The persistence of that pattern---that genius is named only in retrospect and never read off a scale---is itself the best evidence that no scale captures it.

\begin{center}
\emph{If IQ paralleled genius, then it would parallel genius.}
\end{center}

\section{Additional Arguments and Speculations}

\subsection{The Newton Condition, an Application of the Nexic Accordance Principle}\label{app:newton}

One application of a different standing is recorded here by pointer only. The inference from Newton's documented birth to a variance-inflating perinatal class, and from that class to a two-dimensional refinement of the root score, rests on the Accordance Principle of the Nexic companion \cite{NexicFoundation}---a typicality premise this paper does not establish---rather than on any result of the main text. It is therefore set out in a footnote,\footnotemark{} where its evidential standing can be weighed apart from the main argument; nothing in the main text depends on it.
\footnotetext{Symbols in this footnote follow the main text rather than Appendix~\ref{IQ_individual}: $Y = \ln X$ is the latent log-trait, $s = e^{Y-\mu}$ the score of Definition~\ref{def:score}, and $\sigma \approx 0.5911$ the log-space dispersion calibrated in Section~\ref{ssec:physics}. Standardized log-scores are written $z = \ln s/\sigma$, so that Newton's calibrated position~\eqref{eq:newton-score} is $z_N = z_{\max}(N_{\mathrm{Newton}}) \approx 6.61$, with $\bar\Phi = 1 - \Phi$ the standard-normal survival function.


\par\medskip\noindent\textbf{The observation, and what the record supports.}\enspace Newton was born on Christmas Day 1642 (Julian), some three months after his father's death, and the account of his birth is a family legend recorded late. Newton told John Conduitt, eight decades on, that he had heard from his mother that he was ``very much below the usual size of children,'' and so unlikely to live that two women sent to a neighbour for a remedy for him sat down on a stile on the way, saying ``they need not make haste, for the child would be dead before they could get back'' \cite{conduitt1727memoranda}; the version Westfall quotes adds that he was so little ``they could put him into a quart pot,'' and so weak that a bolster was needed to keep his head on his shoulders \cite{westfall1980never}. Westfall's own gloss is careful---``apparently premature''---and the physician Milo Keynes argued that the low birth weight is ``more likely to have been due to poor, or improper, nourishment while in utero than due to premature birth by more than a few days,'' since a markedly premature infant would scarcely have survived in the seventeenth century \cite{keynes1995personality}. The quart pot is conventional hyperbole rather than a measurement: read literally, it would bound the infant near one kilogram. The defensible reading is therefore a very small neonate, whether very preterm or growth-restricted---which is to say, in the terms this footnote needs, an unusual fetal environment.

Taken alone, this is an anecdote. Taken together with a single application of the Accordance Principle, it becomes a quantitative inference, developed here in six steps: the principle indicates that Newton's standing relied, in the sense of carrying most of the statistical weight, on his having been born as unusually as he was, with a strength that depends on which reading of the record is taken; standard statistics, read three ways, agrees on the size of the effect; the explanation is a variance inflation of exactly the kind biological messiness should be expected to produce; the inflation can be bounded from both sides; the bound implies that the log-normal root score of Part~I, while exact for every purpose the bulk of the population serves, is slightly an idealization of a distribution with a second axis; and a second record-holder, born at the opposite tail of fetal growth, shows the class to be a parent state defined by the mechanism rather than by any shared descriptor.


\par\medskip\noindent\textbf{The Accordance Principle, imported.}\enspace The principle is stated, derived, and priced in the companion \cite{NexicFoundation}; we import it as stated. Let $A$ be an attribute of a partial information state---a partial specification in which, before the attribute is interrogated, it stands blank (The companion frameworks name such a state an \emph{Aethus}; nothing below depends on the name)---with complement $A'$, and let $H$ be the production of the outcome under study. Two mediocrity commitments carry the argument: \emph{self-mediocrity}, that the realized value of the attribute is typical among the $H$-producing cases, $P(A'\mid H) \lesssim P(A \mid H)$; and \emph{cosmic-mediocrity}, that the weighting in which these probabilities are taken is the untailored occurrence weighting, fixed before the observation rather than fitted to it. By Bayes' theorem,
\begin{equation}\label{eq:app-bayes}
  \frac{P(A'\mid H)}{P(A\mid H)} \;=\; \frac{P(H\mid A')}{P(H\mid A)} \cdot \frac{P(A')}{P(A)}\,,
\end{equation}
% R836 (Opus 5.5, Oct 21 2026), at the author's direction: P(H|A')/P(H|A) named "the canonical ratio".
so self-mediocrity is exactly the following bound, on the middle factor $P(H \mid A')/P(H \mid A)$, the \emph{canonical ratio} of the companion \cite{NexicFoundation}.

\begin{principle}[Accordance \cite{NexicFoundation}]\label{app:prin-accordance}
\[
  \frac{P(H \mid A')}{P(H \mid A)} \;\lesssim\; \frac{P(A)}{P(A')}\,:
\]
the prior improbability of a realized attribute must be accorded---compensated---by that attribute's relevance to the production of the outcome.
\end{principle}

\noindent The relation $\lesssim$ is made exact in the companion by the \emph{Copernican draw}, under which the bound acquires a complete tail profile with no free tolerance: for every $\kappa \ge 1$, at $\kappa = 1$ in strict form,
\begin{equation}\label{eq:app-exceed}
  \Pr\!\left[\,\frac{P(H \mid A')}{P(H \mid A)} \;\geq\; \kappa\,\frac{P(A)}{P(A')}\,\right] \;\leq\; \frac{1}{1+\kappa}\,,
\end{equation}
the probability taken over the realization \cite{NexicFoundation}. Equivalently, as a quantile ladder, the posterior odds $R = P(A' \mid H)/P(A \mid H)$ against the realized attribute satisfy $\Pr[R \geq (1-p)/p] \leq p$ for $p \in (0, \tfrac12)$, so that the ninetieth percentile of $R$ is at most $9$ and the ninety-ninth at most $99$. The principle's display is the ladder's middle rung, available in the strict form $\Pr[R > 1] \leq \tfrac12$: the realized attribute sits on the majority side of its pair \emph{more often than not}, and it is the \emph{median} of the canonical ratio, not its mean, that the display constrains. Two further features of the full principle bear on what follows. At independence it is silent: if the attribute is statistically unrelated to the outcome, the canonical ratio is identically one and the whole family of bounds collapses onto the draw itself, so the method returns nothing where no relevance is present---the companion's Theorem of Aimless Optimation, recovered as the principle's own value there \cite{NexicFoundation}. And the threshold $\kappa$ must be fixed in advance: choosing it from the realized prior, $\kappa = (1-P(A))/P(A)$, would yield a bound on the prior alone that independence immediately refutes, so no statement that the null has probability at most $P(A)$ is available from the principle.

The companion states the principle for $H$ the production of observers. The present application runs it at the finer grade the principle's origin note anticipates, with the interrogated state specified down to an individual, where ``a smaller population becomes interpretable simply as a lower probability of generating such a person in particular'' \cite{NexicFoundation}. Here the background state is a person born within the main text's cumulative pool $N_{\mathrm{Newton}}$, with perinatal class blank; $A$ is the joint trait \emph{born at least as small as Newton---by prematurity or by growth restriction---and surviving under the care of the era}; and $H$ is the attainment of a score of at least $s_N$. At the individual grade $H$ is binary per person, so the companion's existence-versus-count convention does not arise.


\par\medskip\noindent\textbf{The computation, read three ways.}\enspace The prevalence of the joint trait is the product of a birth-size base rate and a survival probability, and the record fixes neither precisely, so we carry three scenarios. Modern registries place about $0.4$--$0.55\%$ of live births below $28$ weeks and about $1.5\%$ below $32$ weeks \cite{ohuma2023national}; pre-modern rates plausibly lay within a factor of two of these either way, since iatrogenic and multiple-gestation preterm births inflate modern counts while poorer maternal nutrition and infection raised spontaneous ones. Survival without incubators, oxygen, or tube feeding was low but nowhere near $10^{-4}$: the Paris Maternit\'e series of Tarnier and Auvard recorded roughly a third of infants under $2{,}000$~g surviving before the incubator, Budin's hypothermic infants under $2{,}000$~g survived at about $2\%$, Hess's Chicago infants of $605$--$1{,}260$~g at about $27\%$ even with incubators and oxygen, and contemporary special-care units without intensive care report about $14\%$ survival below $28$ weeks and near $60\%$ at $28$--$31$ weeks. Home care in a Lincolnshire farmhouse in midwinter plausibly fell below every such series. On that evidence, the scenarios are: \emph{(A)} an extremely preterm or extremely-low-birth-weight infant, the quart pot taken literally ($\pi \approx 2 \times 10^{-5}$); \emph{(B)} a very preterm or very-low-birth-weight infant of roughly $1{,}000$--$1{,}800$~g, surviving at a few percent ($\pi \approx 5 \times 10^{-4}$, the central case); and \emph{(C)} a moderately preterm or growth-restricted infant, Milo Keynes's reading ($\pi \approx 5 \times 10^{-3}$). The survival filter cuts both ways here: surviving is itself evidence against the most extreme bracket, which is one reason~(B) rather than~(A) is taken as central. In the central case the joint-trait subpopulation of the cumulative pool numbers about $\pi N_{\mathrm{Newton}} \approx 2.5 \times 10^{7}$, and Newton was in it. Write
\begin{equation}\label{eq:app-r}
  r \;=\; \frac{P(H \mid A)}{P(H \mid A')}
\end{equation}
for the per-capita enrichment of the outcome in the joint-trait subpopulation, the reciprocal of the canonical ratio. Three standard readings of the single observation then agree on one scale (Table~\ref{tab:app-scen}).

\emph{The accordance reading.} At the median rung---more often than not---Principle~\ref{app:prin-accordance} gives $r \gtrsim (1-\pi)/\pi$: in the central case the complementary region of the population has at most about a two-thousandth of the joint-trait region's per-capita propensity to house a Newton-class mind. The exceedance form~\eqref{eq:app-exceed} grades the claim rather than merely asserting it: in the central case, with probability at least $0.9$ over the realization, $r \geq 2.2 \times 10^{2}$; at $0.99$, $r \geq 20$.

\emph{The frequentist reading.} Under the null hypothesis of no enrichment ($r = 1$), the record-holder carries the joint trait with probability exactly $\pi$, so the observation rejects the null at $p \approx \pi$: about $5 \times 10^{-4}$ in the central case, $2 \times 10^{-5}$ under~(A), and only $5 \times 10^{-3}$ under~(C). The $p$-value is the chance of the observation under the null; it is not a bound on the chance of the null, and reading it as one would be exactly the substitution the full principle withholds.

\emph{The likelihood reading.} Under enrichment $r$ the record-holder carries the trait with probability $\pi r / (\pi r + 1 - \pi)$, so the likelihood ratio of $r$ against the null is $r / (\pi r + 1 - \pi)$: increasing in $r$ and bounded above by $1/\pi$, the most evidence any single observation of this kind can carry. At the accordance bound itself the likelihood ratio is exactly half its ceiling, $1/(2\pi)$---the typicality premise and the evidential ceiling are two faces of one quantity.

\par\smallskip\begin{center}\footnotesize
  \captionof{table}{The three scenarios for the joint trait. For each: the prevalence $\pi$; the accordance floor on the enrichment at even odds and at $0.9$ confidence; the frequentist $p$; the minimal inflation $\gamma_{\min}$ of the mechanism paragraph by the per-capita and record-holder routes ($\delta = 0.8$); the two-sided $90\%$ band of the paragraph bounding the second axis; and the crossover $z^{\ast}$.}
  \label{tab:app-scen}
  \medskip
  \resizebox{\linewidth}{!}{%
  \begin{tabular}{l r r r r c c r}
    \hline
    Scenario & $\pi$ & $r \gtrsim$ & $r_{0.9} \geq$ & $p$ & $\gamma_{\min}$ (two routes) & $90\%$ band & $z^{\ast}$ \\
    \hline
    (A) extremely preterm/ELBW & $2\times10^{-5}$ & $5.0\times10^{4}$ & $5.6\times10^{3}$ & $2\times10^{-5}$ & $1.556$ / $1.555$ & $[1.39,\,1.64]$ & $1.44$ \\
    (B) very preterm/VLBW      & $5\times10^{-4}$ & $2.0\times10^{3}$ & $2.2\times10^{2}$ & $5\times10^{-4}$ & $1.379$ / $1.378$ & $[1.26,\,1.43]$ & $2.11$ \\
    (C) growth-restricted      & $5\times10^{-3}$ & $2.0\times10^{2}$ & $22$              & $5\times10^{-3}$ & $1.283$ / $1.283$ & $[1.19,\,1.33]$ & $2.83$ \\
    \hline
  \end{tabular}}
\end{center}

\emph{Robustness.} The bound scales as $1/\pi$, so the conclusion's strength is set by the reading of the record rather than by fine details of the inputs (Table~\ref{tab:app-sens}): strong under~(A) and~(B), where even the $0.9$-confidence floor stays in the hundreds or above, and markedly weaker under~(C), where the $0.99$-confidence floor falls to $r \geq 2$. The one statistical hazard specific to biographical inference is the look-elsewhere effect: had the trait been singled out from $M$ candidate traits of comparable rarity, the effective null probability would be about $M\pi$---in the central case $5 \times 10^{-3}$ for $M = 10$ and $0.05$ for $M = 100$. Two facts keep $M$ small here. Birth size is a standard perinatal variable that would have been interrogated whatever its value, which is what cosmic-mediocrity demands; and Newton was selected as the record-holder by the main text's calibration before any perinatal fact entered, so the individual was not chosen for the trait. That establishes the individual's independence from the trait, not the category's: the birth-history class and its thresholds were formulated with Newton's recorded birth in view, and the $p$-values below assume an independence of their choice that is not established here.

\par\smallskip\begin{center}\footnotesize
  \captionof{table}{Accordance floor $r_{\min} = (1-\pi)/\pi$ across the input grid: birth-size prevalence against survival under the care of the era. The central case sits near the $1.5\%$ row between the $1\%$ and $4\%$ survival columns.}
  \label{tab:app-sens}
  \medskip
  \begin{tabular}{l r r r r r}
    \hline
    & \multicolumn{5}{c}{survival under the care of the era} \\
    prevalence & $0.1\%$ & $1\%$ & $4\%$ & $10\%$ & $50\%$ \\
    \hline
    $0.5\%$ & $2.0\times10^{5}$ & $2.0\times10^{4}$ & $5.0\times10^{3}$ & $2.0\times10^{3}$ & $4.0\times10^{2}$ \\
    $1.5\%$ & $6.7\times10^{4}$ & $6.7\times10^{3}$ & $1.7\times10^{3}$ & $6.7\times10^{2}$ & $1.3\times10^{2}$ \\
    $5\%$   & $2.0\times10^{4}$ & $2.0\times10^{3}$ & $5.0\times10^{2}$ & $2.0\times10^{2}$ & $39$ \\
    \hline
  \end{tabular}
\end{center}


\par\medskip\noindent\textbf{The mechanism: a risky dice roll.}\enspace With the Accordance Principle itself enforcing the categorical conclusion, the accompanying explanation is straightforward. An awkward fetal environment of the class that Newton's birth instantiates creates, qualitatively, a risky dice roll: heuristically, a standard-deviation expansion in cognitive potential. Like mutations under the theory of evolution, the roll is in most cases probably detrimental; but conditioning on landing a score of the kind the record-holders exhibit effectively requires it, up to the vast majority of statistical weight, because the expansion alone brings the dispersion to where landing such a score becomes reasonably probable.

Enrichment alone does not identify its mechanism: variance inflation is one mechanism that would produce it, taken here as a hypothesis, and others could produce the same enrichment. Formally, let the complementary population keep the main text's law, $\ln s \mid A' \sim \mathcal{N}(0, \sigma^{2})$, and let the joint-trait class carry an inflated and possibly lowered law,
\begin{equation}\label{eq:app-mix}
  \ln s \mid A \;\sim\; \mathcal{N}\bigl(-\delta\sigma,\; \gamma^{2}\sigma^{2}\bigr), \qquad \gamma > 1,\;\; \delta \geq 0,
\end{equation}
with $\gamma$ the inflation and $\delta$ the mean detriment in standardized units. The per-capita enrichment above a standardized threshold $z$ is then
\begin{equation}\label{eq:app-R}
  R(z) \;=\; \frac{\bar\Phi\bigl((z+\delta)/\gamma\bigr)}{\bar\Phi(z)}\,.
\end{equation}

\begin{proposition}[Tail dominance]\label{app:prop-tail}
For $\gamma > 1$ and any $\delta$---a detrimental mean shift included---$R(z) \to \infty$ as $z \to \infty$, so that for any class prevalence $\pi > 0$ the posterior
\[
  P\bigl(A \mid \ln s \geq \sigma z\bigr) \;=\; \frac{\pi R(z)}{\pi R(z) + 1 - \pi} \;\longrightarrow\; 1.
\]
Moreover, $R(z) > 1$ exactly when $z > z^{\ast} := \delta/(\gamma - 1)$.
\end{proposition}

\noindent\textit{Proof.} By the Mills-ratio asymptotic $\bar\Phi(u) = \varphi(u)\,u^{-1}\bigl(1 + O(u^{-2})\bigr)$,
\[
  \ln R(z) \;=\; \frac{z^{2}}{2} - \frac{(z+\delta)^{2}}{2\gamma^{2}} + \ln\frac{z\gamma}{z+\delta} + o(1),
\]
whose leading term $\tfrac12(1-\gamma^{-2})z^{2}$ diverges because $\gamma > 1$. For the crossover, $\bar\Phi$ is strictly decreasing, so $R(z) > 1$ if and only if $(z+\delta)/\gamma < z$, that is, $z > \delta/(\gamma-1)$. $\square$

\medskip
\noindent This is the mutation analogy made exact: a class whose members do worse on average ($\delta > 0$) nonetheless owns the extreme tail, because its tail decays as $e^{-z^{2}/2\gamma^{2}}$ against the bulk's $e^{-z^{2}/2}$. Combined with the accordance floor, the proposition fixes how large the roll must be: the class must reach Newton's threshold at a per-capita rate of at least $r_{\min}\,\bar\Phi(z_N)$, so that
\begin{equation}\label{eq:app-gmin}
  \gamma \;\geq\; \gamma_{\min} \;=\; \frac{z_N + \delta}{\bar\Phi^{-1}\bigl(r_{\min}\,\bar\Phi(z_N)\bigr)}\,,
\end{equation}
whose denominator is about $5.38$ in the central case. The same number arrives by a second, independent route: the typical maximum of the joint-trait subpopulation, $-\delta + \gamma\, z_{\max}(\pi N_{\mathrm{Newton}})$ with $z_{\max}$ from~\eqref{eq:zmax}, exceeds the typical maximum $z_N$ of the rest exactly when $\gamma \geq (z_N + \delta)/z_{\max}(\pi N_{\mathrm{Newton}})$, and the per-capita route and the record-holder route agree to three figures in every scenario (Table~\ref{tab:app-scen}). For the mean detriment, meta-analyses of very preterm and preterm cohorts report deficits of about $0.73$ to $0.89$ standard deviations on measured scores \cite{bhutta2002cognitive,twilhaar2018cognitive,eves2021association}, while extremely preterm cohorts show larger ones; taking $\delta \approx 0.8$ as central, the roll must inflate the log-dispersion by a factor of about $1.38$, ranging from $1.23$ at $\delta = 0$ to $1.60$ at an extremely-preterm-sized $\delta = 2$.

Two further features deserve statement. First, the posterior is sharp: in the central case at $(\gamma, \delta) = (1.38, 0.8)$ it rises from $0.17$ at $z = 6$ to $0.51$ at Newton's $z_N$, $0.76$ at $z = 7$, and $0.986$ at $z = 8$---a switch from minor to near-certain over about two standardized units, echoing without duplicating the sharp crossover of Section~\ref{ssec:phase}. Second, the requirement \emph{relaxes} as the achievement grows more exceptional: the minimal inflation falls from $1.47$ at $z = 6$ to $1.38$ at $z_N$, $1.26$ at $z = 8$, and $1.17$ at $z = 10$, tending to $1$ in the limit; below $z \approx 3.5$, no finite pure inflation meets the central-case floor at all. Variance-carried typicality is a deep-tail phenomenon---which is precisely where the main text places Newton.


\par\medskip\noindent\textbf{Bounding the second axis.}\enspace The inflation is bounded from both sides, and the two bounds meet in a narrow band. From below, the accordance floor gives $\gamma_{\min}$. From above, the calibrated height of the record: were the class's inflation much larger than $\gamma_{\min}$, the class would be expected to have produced a record higher than the calibration admits. Requiring the calibrated record $z_N$ to lie within the central $90\%$ of the distribution of the class maximum---$\Pr[-\delta + \gamma M_n \leq z_N] = \Phi\bigl((z_N+\delta)/\gamma\bigr)^{n}$, with $M_n$ the maximum of $n = \pi N_{\mathrm{Newton}}$ standard normals---pins the inflation, in the central case, to
\begin{equation}\label{eq:app-band}
  1.26 \;\lesssim\; \gamma \;\lesssim\; 1.43 \qquad (\delta = 0.8,\ \pi = 5 \times 10^{-4}),
\end{equation}
with the corresponding bands $[1.39, 1.64]$ and $[1.19, 1.33]$ under~(A) and~(C). The bound is relative to the main text's calibration of the record, and inherits its uncertainty. It also treats the record as the maximum of the class alone, with $z_N$ from the main text's single-log-normal calibration; under the mixture the record is the maximum over both components, whose law is $\Phi(z)^{(1-\pi)N_{\mathrm{Newton}}}\,\Phi\bigl((z+\delta)/\gamma\bigr)^{\pi N_{\mathrm{Newton}}}$, and $z_N$ would itself need recalibrating from that law, so the band is a first approximation and not a consequence of the mixture.

The data on modern preterm cohorts then bear on the band directly, and they cut in an instructive direction. Dispersion is inflated: the EPICure cohort of infants born before $26$ weeks scored $82$ (SD $19$) at age six against $106$ (SD $12$) for classmates \cite{marlow2005neurologic}, a ratio of $1.58$, and a Swedish extremely preterm cohort scored $93$ (SD $15.4$) at eighteen against $106$ (SD $12.5$) \cite{hallin2010followup}, a ratio of $1.23$. But the inflation is carried by the lower tail---in EPICure, some two-fifths of the children scored more than two standard deviations below the classmate mean---and under a normal approximation to these figures only about $0.6$--$0.7\%$ of the preterm children exceed $+2$ standard deviations of the control distribution, against $2.3\%$ of controls, with the true upper tail plausibly thinner still given the skew. The relevant quantity for the mechanism is therefore not the total dispersion ratio but the \emph{upper-tail} dispersion, and for modern extremely preterm survivors the data suggest an upper-tail ratio of about $1.0$ to $1.2$---at or below the lower edge of the band~\eqref{eq:app-band}.

The tension is modest, and it has four candidate resolutions, not mutually exclusive. The survival filter: a seventeenth-century survivor of a very small birth is a strongly selected member of the class, and modern neonatal care has removed most of that selection, so modern cohorts may not share the historical class's dispersion. Latent against measured: the prediction concerns the latent score, and observed scores decouple from latent capacity in exactly the tail where the mechanism operates (Appendix~\ref{IQ_individual}). The axis itself: under Milo Keynes's reading the variance-relevant deviation is growth restriction, a different fetal environment from the one modern extremely preterm cohorts sample. And the typicality premise may simply fail for this realization, which the exceedance form~\eqref{eq:app-exceed} prices rather than forbids. Which resolution holds is an empirical matter; what the band establishes is that the required inflation is modest---some forty percent in the log-dispersion, not a multiple---and that the data in hand constrain its form more than its existence.


\par\medskip\noindent\textbf{Invisible in the cohort, decisive at the record.}\enspace The prevalence of the class and the size of the roll trade against each other: the rarer the class, the larger the inflation the accordance floor demands, and the nearer the crossover $z^{\ast}$ sits to the mean (Table~\ref{tab:app-scen}). At the central parameters, per thousand members the class has $21.2$ above $z = 2$ against the bulk's $22.8$---slightly fewer, consistent with reported mean deficits---yet $2.9$ above $z = 3$ against $1.35$, $0.92$ above $z = 3.5$ against $0.23$, and $0.25$ above $z = 4$ against $0.03$: an enrichment no cohort of thousands can resolve, and one that nonetheless decides who holds the record. For the modern extremely preterm class, with its larger mean detriment ($\delta \approx 1.6$--$2.0$), the crossover at the same inflation lies near $z \approx 4.2$--$5.3$, so the depletion observed at $+2$ standard deviations is exactly what the model predicts rather than evidence against it. The mechanism hides by construction. It is invisible in the cohort and decisive at the record.


\par\medskip\noindent\textbf{A two-dimensional root score.}\enspace We might therefore picture that even a truly log-normal root score is itself slightly an idealization: due to the messiness of biological factors, the real distribution is closer to a two-dimensional one, for which the usual log-normal is one axis and a second axis inflates the first axis' distributional standard deviation. Formally, with $D$ a standardized fetal-environment deviation,
\begin{equation}\label{eq:app-2d}
  \ln s \mid D \;\sim\; \mathcal{N}\bigl(\sigma\, m(D),\; \sigma^{2} v(D)^{2}\bigr), \qquad v(0) = 1,\;\; v \text{ nondecreasing in } |D|,\;\; m \leq 0,
\end{equation}
so that the standard deviation of the normal cross-section at $D$ is $\sigma\,v(D)$, and the observed root score is the marginal, a location--scale mixture of normals in log-space. The joint-trait class of the preceding sections is a tail of this axis, $A = \{D \geq d_N\}$, and the analysis above constrains the multiplier $v$ at that tail. Four consequences follow.

\emph{(i) The multiplier is bounded, and saturates.} At $D = 0$ it is $1$ by construction; at the Newton class it lies in the band~\eqref{eq:app-band}; and it cannot grow without limit. If $v$ grew without bound in $|D|$---say $v = 1 + \beta|D|$ with $D$ Gaussian---the log-space tail would become exponential, as for a product of Gaussians, and the score would acquire a power-law tail, $P(s > x) \sim x^{-1/(\sigma\beta)}$: an input-side route to Pareto-like behavior. The residual scaling by which Part~I selects the log-normal family (Section~\ref{ssec:motivation}) rules that out at the scales probed, so the natural form is a saturating one, $v(D) = 1 + (v_{\infty} - 1)\bigl(1 - e^{-D/\ell}\bigr)$ for $D \geq 0$, with $v_{\infty}$ near the band and the scale $\ell$ unconstrained by present data. The endpoints are pinned; the shape between them is not.

\emph{(ii) The cross-section is not symmetric.} The cohort data show that the inflation of~\eqref{eq:app-2d} is realized chiefly in the lower tail, so the conditional law at large $D$ is left-skewed rather than normal. The mechanism runs only on the upper-tail dispersion $v_{+}(D) \leq v(D)$, and the band~\eqref{eq:app-band} is a bound on $v_{+}$. The two-dimensional score is thus an idealization twice over: the log-normal is one cross-section of it, and each cross-section is itself only approximately normal.

\emph{(iii) The slightness is quantifiable.} For the two-point case of~\eqref{eq:app-mix} at $\delta = 0$, the log-space kurtosis is
\[
  \frac{3\,\bigl[(1-\pi) + \pi\gamma^{4}\bigr]}{\bigl[(1-\pi) + \pi\gamma^{2}\bigr]^{2}} \;\approx\; 3 + 3\pi\bigl(\gamma^{2}-1\bigr)^{2},
\]
an excess of about $1.2 \times 10^{-3}$ in the central case. No fit to the bulk can see it---which is why Part~I's log-normal is exact for every purpose the bulk serves. The overall log-variance is likewise inflated only by the factor $1 + \pi(\gamma^{2} - 1 + \delta^{2})$, far inside the calibration's own uncertainty in $\sigma$.

\emph{(iv) The record-holder is drawn from the inflated component.} The dispersion that Part~I reads from the record-holder residual is therefore effective---tail-weighted rather than bulk. The main text's completion-time results are calibrated to observed record scores and are unaffected; what changes is the reading of the record-holder, now the maximum of the high-variance component rather than of a single log-normal. This supplies a mechanism for the main text's suspicion that Ramanujan's intrinsic $z$-score, ``a lucky enough dice roll,'' might sit ``a fair bit higher than expected by extrapolating from his contemporaries'' (Section~\ref{ssec:games})---the same idiom, reached for independently there.


\par\medskip\noindent\textbf{Neighbors in the eminence literature.}\enspace The adversity-as-variance reading has psychosocial neighbors. Eisenstadt's study of $699$ eminent figures found that a quarter of those with data had lost a parent by age ten and over half by twenty-five, while noting the methodological difficulty of fixing base rates \cite{eisenstadt1978parental}; Damian and Simonton list orphanhood among the ``disruptive events'' that diversify development, with the explicit caution that cumulative adversity can prove excessive and inhibit rather than enhance creative growth \cite{damian2014diversifying}---the mutation analogy in empirical form. The present treatment adds a prenatal locus, a quantification from a single record-holder, and a bound on the inflation. The neighborhood also sharpens a caution. Newton belongs to more than one rare early-adversity class: he was born posthumously, and was left with his grandmother at about three when his mother remarried. The joint trait analyzed here is not isolated from these; they are competing members of the same class, and the accordance argument identifies the class, not its member.


\par\medskip\noindent\textbf{A second instance, and the class as a parent state.}\enspace Einstein is the second instance. No reliable source reports him premature; the documented newborn was large and heavy, with a head so large and misshapen that his family feared a deformity until it normalized within weeks \cite{isaacson2007einstein}. Newton and Einstein thus sit at opposite tails of fetal growth, one very small and one very large, and no surface descriptor is common to the two except the extremity itself. That is the right result, because the class the mechanism requires was never a surface descriptor. Call a fetal state \emph{awkward} when it lies at a surprisingly high or low quantile on at least a handful of perinatal attributes---birth weight, length, head size and shape, and gestational timing among them---and write $\mathrm{Awk}$ for the class of such states. The class $C$ that realizes the risky dice roll of the mechanism paragraph is a refinement of it, $C \sqsubseteq \mathrm{Awk}$: every member of $C$ is awkward, but awkwardness is necessary for membership rather than sufficient, a feature of the class's description rather than the whole of it. (In the companion frameworks' terms, $\mathrm{Awk}$ is a parent state of $C$, and $C$ a state-decomposable attribute whose observable constituents are the awkward quantiles, the variance-realizing constituent remaining blank at the grade of any birth record \cite{NexicFoundation}) Three consequences follow.

\emph{(i) Observed descriptors approximate the class from inside its scope.} What two members happen to exhibit---extreme birth size, in both directions---is a two-member sample of the class's forms, narrower than the class itself, which extends to awkwardness in any of the handful of attributes. The common descriptor of the two is therefore not the class but an approximation of it, and likely a less general one.

\emph{(ii) Bounds proven for the parent transfer to the child, concentrated.} The accordance principle applies to what a record can interrogate, which is awkwardness or its observed instance, not the unobservable roll. But if the enrichment within $\mathrm{Awk}$ is carried by $C$---so that $\mathrm{Awk} \setminus C$ is no more productive of the outcome than the complement---then with $r_{\mathrm{Awk}}$ and $r_{C}$ the per-capita enrichments of the two classes and $\pi_{\mathrm{Awk}}, \pi_{C}$ their prevalences,
\begin{equation}\label{eq:app-concentrate}
  r_{C} \;\geq\; 1 + \frac{\pi_{\mathrm{Awk}}}{\pi_{C}}\,\bigl(r_{\mathrm{Awk}} - 1\bigr),
\end{equation}
so the child is enriched beyond the parent by the concentration factor $\pi_{\mathrm{Awk}}/\pi_{C}$, and the inflation it must carry is correspondingly larger than the parent-level bound, the rarer class demanding the larger $\gamma_{\min}$ (Table~\ref{tab:app-scen}).

\emph{(iii) A second instance compounds the evidence.} If the class is fixed by the mechanism independently of either record-holder's history, which the present formulation, made with both histories in view, does not establish, and the two perinatal draws are independent, the chance under the null that both record-holders fall in $\mathrm{Awk}$ is $\pi_{\mathrm{Awk}}^{2}$, and the likelihood ceiling rises to $\pi_{\mathrm{Awk}}^{-2}$. The class is broad: at an illustrative prevalence of $10^{-2}$, counting both tails of a handful of correlated attributes, each instance alone yields only $p \approx 10^{-2}$, yet the pair yields $p \approx 10^{-4}$---stronger than the single-instance central case, from a class some twenty times broader---and the exceedance bound~\eqref{eq:app-exceed}, applied to each independent draw, confines a $\kappa$-fold violation in both to probability at most $(1+\kappa)^{-2}$. The documentary grades differ and should be kept apart: Newton's record is a family legend written down eight decades on, Einstein's a family account in a modern biography, and neither records a weight.


\par\medskip\noindent\textbf{Predictions.}\enspace \emph{The next Newton.} Because the conclusion is a general property of Newton's score under a single accordance application, it can be imported to the present: the next Newton-achieving intelligence would, on the model and if the enrichment is carried by this class, be expected to have had an unusual fetal environment in some sense---several standard deviations from the population median under one or several measurable perinatal traits, premature birth, growth restriction, and unusually large birth being examples. The accordance bound is a repeated-sampling statement and supplies no posterior over mechanisms; a posterior would need priors over the candidate mechanisms, which are not supplied here. The prediction concerns membership in the awkward-state class of the second-instance paragraph, of which any listed trait is an observed instance rather than the definition. The prediction gains robustness from its disjunctive form: with several independent variance-inflating axes, the union class has prevalence approximately $\sum_{i}\pi_{i}$, and ``at least one axis'' is the claim. It carries force only with the class fixed in advance.

\emph{The survival filter.} Newton's joint trait bundles a fetal-environment deviation with survival under seventeenth-century conditions, a severe selection in its own right. Modern neonatal care has largely removed that filter, so the prediction ranges over the class of variance-inflating perinatal deviations rather than over any single modern diagnostic category, and whether such deviations inflate the upper tail absent the historical filter is a question the model poses rather than answers.

\emph{The cohort test.} Proposition~\ref{app:prop-tail} discriminates the variance mechanism from a pure mean shift: a mean-shift-only model predicts per-capita depletion of the class at every threshold, while inflation predicts enrichment above the crossover $z^{\ast} = \delta/(\gamma-1)$ despite the lower mean. At the central parameters and a very-preterm-sized mean detriment, the class would be roughly twice as frequent as the bulk above $+3$ standard deviations and four times as frequent above $+3.5$. Published outcome summaries emphasize means; the test requires the upper quantiles of individual-participant data pooled at a scale of order $10^{4}$ very preterm adults, beyond the roughly thousand-per-arm of the largest current synthesis \cite{eves2021association}, and it inherits the measurement caveats of Appendix~\ref{IQ_individual}.


\par\medskip\noindent\textbf{Scope.}\enspace Five limits bound the argument. First, the inference for any one trait rests on a single observation, and its whole weight rides on the typicality premise; the exceedance form~\eqref{eq:app-exceed} prices that weight, and the second instance, with its class declared by the mechanism, multiplies the evidence rather than merely adding to it. Second, its strength depends on the reading of the birth record: decisive under the extremely preterm reading, strong under the central very-preterm reading, and modest under the growth-restriction reading, where the high-confidence floor on the enrichment falls to single digits. Third, the documentary basis is a family legend recorded some eighty years after the event, and the scenario analysis is the discipline that legend requires. Fourth, the class is not decomposed: prematurity, growth restriction, the survival filter, and the posthumous and early-separation adversities are not separated here. Fifth, the analysis is descriptive, not normative. The same model that assigns the class the record also says most of its rolls are detrimental ($\delta > 0$); it identifies which histories the extreme tail is drawn from and says nothing about the worth of any person inside the tail or out of it, in keeping with the commitment stated in Section~\ref{app:lonegenius}. Einstein enters as a second instance of the awkward-state class, not of prematurity (see the second-instance paragraph).}

\subsection{A Speculative Application to Quantum Foundations}
A speculative application of the framework concerns the quantum
measurement problem, widely regarded as the central open problem in the
foundations of quantum mechanics \cite{spekkens2010einstein}. We include
this application as an illustration of how the model's general
predictions bear on a contemporary question; the validity of the
framework developed in Parts I and II does not depend on the specifics
of what follows, which should be read independently of any contemporary
work in quantum foundations.

Despite sustained collective effort spanning nearly a century\footnote{I want to emphasize that I feel the situation with quantum foundations is literally identical to what would have happened with special relativity in the counterfactual where Einstein was born later or not at all. What we are seeing with special relativity is not that it was ``up in the air", but rather than Einstein was coincidentally and luckily born early enough to cause the illusion of this by means of actually solving it. I suspect that once the measurement problem is solved, (specifically to the standards put forward by Spekkens himself \cite{spekkens2010einstein}), it will have the exact same character of feeling ``up in the air" over this past century---\textit{and yet}---where was the solution that whole time? This is not meant to dismiss the academic community itself so much as to highlight the hard truth that this sort of ontological multi-step sequence of an idea requires people on the very tails of the distribution to have even a reasonable time complexity of completion in the first place. The community is perfectly adequate, and indeed vastly \textit{better} in terms of time complexity at solving problems with low numbers of steps in sequence, because everyone there is in the linear regime, so it only comes down to that many hands are quicker than one alone. Hence, the heart and spine analogy itself---the open questions of the universe are themselves bimodal in what it takes to achieve them, so dropping either the many or the few would stall human progress indefinitely.\par
So in short, yes the mathematical pieces for special relativity very much directly accessible in the literature, but like quantum foundations the field easily would have circled special relativity for a century without landing it until the outlier person came along, (because, combinatorially speaking, it physically could not have resolved any other way). What I want to emphasize as well is that the lack of the ontology of special relativity leading up to 1905 is not a feature of the game but \textit{the game} concerning what was absolutely necessary for forward progress. Pais and Holton treat Poincaré's and Lorentz's pre-1905 work on the dynamics of the electron and the Lorentz transformations as constituting most of the mathematical structure of special relativity, with the ontological completion being the missing piece \cite{pais1982subtle, holton1988thematic}. Their framing, however, seems to presuppose that the ontological component is on equal footing with the other components, as if just another brick to be placed on the metaphorical wall. However, the reality as I see it is that the ontological component is like the Christmas tree upon which the mathematical components are strung about like lights, so removing that leaves you with shattered light bulbs on the floor and no theory. More explicitly speaking, the ontological substrate of special relativity \textit{is the high-step number} sequence in its case. There seems to be a bias in viewing scientific revolutions retrospectively that there is a natural fit between the ontology and the mathematical structure as if obvious, but I wish to point out that this is systematically not how such revolutions are viewed \textit{before}, when the mathematical pieces are scattered about and the ontological substrate---like the Christmas tree---is a puzzle of weaving those pieces about the topology of whatever philosophy of the universe will receive them. There is no need to view this in the abstract---just look at modern quantum foundations for goodness sake! The felt sense of living in it is an insurmountable Gordian knot of a possibility space, even though our descendants are sure (by my prediction) to look back at us with a puzzling expression at how we managed not to associate the obvious ontology with the obvious mathematical structures in the first place. From their perspective, so too will it seems as though anyone could have assembled, in an afternoon, all the loose mathematical bits that we already possess the exhaustive supply of. The problem, however, is in the framing---the having of all the loose mathematical pieces is both the reason it should be solvable and the reason it is not solved.}, a
satisfactory realist resolution to the measurement problem\footnote{To give an example of the relationship with special relativity, my prediction is that at some point it will become rather obvious that all these no-go theorems in quantum foundations were entirely orthogonal to solving the measurement problem. Did they build insight in how to engage with quantum mechanics operably? Yes, of course. Did they move us even a hair closer to the measurement problem resolution? Not by a longshot, because not a single one of them is the kind of move that could disturb the category error Spekkens \cite{spekkens2010einstein} points out as existing: a no-go theorem takes the standing categories as given and proves what they exclude, which is linear-regime work, while relocating a category is the high-step sequence, and no accumulation of the first kind of move sums to the second. Indeed, the no-go theorems in this view are just a crawl in person-years against the very combinatorially-nightmarish time-spans of this paper. So, with whatever winds up solving the measurement problem as my witness, what I'm trying to express is that exactly what happened with these no-go theorems is what would have happened had Einstein not explicitly pointed out the relativity of simultaneity in 1905---decades upon decades of desperately trying to perturb the paradigm without even the awareness to realize that everyone is one category error away from the resolution.\par
Also note that by my earlier use of ``orthogonal" I mean that yes, the probability space is indeed lowered, but that happens by the first couple of no-go theorems and not after, and furthermore I expect that it will if anything be easier to identify the category error without them, since even defining them smuggles it in anyway. Then once the category error is found, it should be easy enough to restate its resolution as a refutation of that definition.} remains
elusive\footnote{My feeling is that, in a specific sense made precise below, it is just as easy, and as hard, to do revolutionary work in any era of history as in any other. The perception that the modern era is not particularly primed for natural philosophy is simply a consequence of more jobs being available which do not require it.\par
By this I mean that the people in history will have felt their now-resolved open questions equally intrinsically difficult to how we find ours, such that if it would have been straightforward then, it would be straightforward now, which it is not. And by the modern natural philosopher thing I simply mean that the intrinsic environment of unsolved problems is primed for it just as it was in any similar historical vacancy, just if you ask a random scientist on the street they will be aware of a thousand non-natural philosopher-requiring jobs to every one that requires it, and so like a statistical overshadowing effect people falsely assume that the smaller category does not have a presence. It is similar to if people said ``doctors do not exist because there are far more farmers than doctors'', such that the only reason doctors are known to exist anyway is because we have a selected category for them. Meanwhile, ``natural philosophy-requiring tasks'' are potentially even vastly more proportionally rare, and so that together with not having a category for those tasks leads to the perception that they do not exist, despite the actual shape of the open questions---so far as anyone can read it from the questions themselves---requiring and responding to those types of methods as much now as before. So, if you will, the intrinsic natural philosophy-requiring tasks remained unchanged in numerosity, but the independent inflation of tasks not requiring them created the illusion of wiping them out as a category.}. If its resolution constitutes a completion type of substantial
depth---requiring a sequence of consecutive ontological insights each of
which must be correct and must build on its predecessor without
error---then the relevant reference class is the small population of
researchers actively working on quantum foundations, on the order of
$10^3$. The crossover score at this population's depth lies well above the mode of
the broader theoretical physics population~$B$
($s_{\mathrm{mode}} \approx 12.1$), meaning that most quantum
foundations specialists sit above the critical threshold in terms of
per-step probability. Nevertheless, the depth of the required
ontological chain is large enough that even scores of $s \approx 30$--$35$
(the range of the contemporaneous top 10--20) would produce completion
times on the order of billions of years under the model's assumptions.

Several caveats are important\footnote{I think something of an elephant in the room worth mentioning here is the implications of people trying to use the concepts of this paper to impose hierarchies on aspects of the world. This is a feeling I get while staring at something like the Landau scale \cite{mermin1984homage}, which while perhaps mildly academically interesting due to the historical subject matter, also seems outwardly to impose something onto that matter which was never there to begin with. That is, \emph{why do humans have to make unsolicited hierarchies out of everything?} There is obviously the moral unease itself that such thinking imposes when it has devolved far enough, but I think its fallacious nature is more direct still.
Do you think a tree in a forest cares about who gets to be the ``alpha" among the lichen on its bark or the leaves on its branches? Of course not. It is not only reductionist but attempting to ascribe reducible meaning to the image of the tree based on further reduction of the reduced aspects to an artificially imposed model that is both arbitrary to the modeler and irrelevant to the tree. (That is, even in cases where evolutionary completion is at play like a vine trying to block a smaller tree for sunlight, and still the ascription of a darn ``hierarchy" to this imparts to the thinker about as much real insight as a spreadsheet of the guessed number of words on each page of a book imparts about the allegorical devices throughout said book).
In my mind a hierarchy model is a corruption of the organic nature of life, (think lichen and bark), which is a central reason behind why their often being problematic is intertwined with their explicit incompatibility with the real world to begin with.}. The depth parameter for the measurement
problem is not directly measurable; the sequential-success structure
with full restart is one of several possible models of how foundational
work proceeds, and alternatives (partial credit across independent
insights, collaborative aggregation, conceptual shortcuts that
effectively compress depth) could substantially alter the quantitative
picture. The framework should therefore be read as one structural
perspective on why deep foundational problems can resist solution for
extended periods, not as a definitive account of this or any specific
case.

Within these caveats, the reading the model suggests is that the
continued absence of a solution is not, by itself, evidence that the
problem is insoluble or that existing approaches are misguided---the
task structure may simply place it in the exponential regime for all
but the most extreme outliers in a very small population. The
framework's mathematical structure also suggests, as a generic
consequence rather than a prediction about any specific instance, that
when progress on problems of this structural type does arrive, it will
tend to appear suddenly rather than incrementally, and will tend to be
attributable to concentrated individual work rather than diffuse
collective convergence. This is not a claim about the properties of any
individual; it is a claim about how the sharp-crossover mathematics concentrates realized output when task depth is large and population
size is small.

\subsection{The ``Lone Genius" Narrative is a Problematic Cause}\label{app:lonegenius}

\noindent\emph{(Status: this closing section is authorial moral reflection on the topic's ethics; nothing in the paper's formal or empirical content depends on it.)}

I want to take a moment to state my thoughts on the lone genius narrative. Specifically, I notice that any rational agent who believes there is truth to that narrative, (be it through the mathematical illustration of this paper or otherwise), will be met with an uncomfortable moral dilemma upon analyzing their support for it: \textit{who even does this narrative serve?}
That is, I might ask, \textit{what is the goal here, really}, and upon thorough reflection the only reasonable answer is that for society to adopt something closer to this narrative would probably be a net negative. For those who are already labeled geniuses, it would be a neutral effect, because that label is unaffected by the peripheral sociological contextualization of the narrative, and for those who are labeled mediocre the narrative seems only to be negative, because now they have to fight an uphill battle against it when advertising their own research. So, in that sense, it serves no one, and for that reason seems like a suspiciously immoral ``cause", if you can even call it that in the first place. Think especially about the egalitarian\footnote{To be abundantly clear about this, I would say that the major reason the egalitarian view ought to be defended tooth and nail is that it is the only conceivable way at present to preserve human dignity and, yes, even meritocracy. It is naive and foolish to think that you must leave egalitarianism to increase meritocracy, because once you abandon egalitarianism, suddenly the measure for merit is the imperfect and corrupt human opinion of those with authority, and if history is any example individual humans have about as much taste in judging merit as \textit{Yersinia pestis} has in practicing pacifism, (this does not apply to wisdom of the crowd, however, which I maintain is our species' saving grace throughout history). So the only end state of an anti-egalitarian society is a fascistic nightmare dystopia, and you could expect this to compound drastically the more the petty authorities pushing the transformation realize what they can get away with.} view which this undermines---a view that everyone deserves a voice and a contribution---which is hopefully quite obviously a wonderful view of humanity, and should be aspired to. So is the backtracking against this not itself immoral?

The specific philosophical dilemma I want to express here is the moral difficulty of contesting elements of the egalitarian view while still maintaining its beauty as a principle. For example, I have heard before people saying things like ``considering people like Ramanujan or Einstein, our society just wants to worship old white men\footnote{I want to support this resistance to the total authority of old white men with particular ferocity. Looking back to the great man theory \cite{hook2004great}, once again, my personal stake in the game is that \textit{I do not want the patriarchal entity to have the satisfaction of feeling validated by my views on the irremovability of high performers}. As I have emphasized earlier in this paper, the actions of white-centric patriarchy throughout history has blood on its hands, and altogether it has done what it has done out of ignorance masking as unquestionable authority. So it turns out to be embarrassing and undignified on both fronts of vision and morality. Hence, why even humoring this is a bridge too far---it ought only to be denounced for its actions.}, so they happen to be the people who were selected. But there are plenty of brilliant people." This is a perfect illustration of the moral unease I get, because I feel the need to first assert that \textit{no, the reason Einstein is respected the way he is is because he is very objectively at a different caliber of contribution compared to anyone else of the last several centuries}. The goal of this paper, in part, is the quantify what exactly I mean by that---Einstein is not just chosen arbitrarily as a genius, he is quite measurably operating in a different regime from everyone else, and for that reason is irreplaceable (a claim the model makes quantitative only conditionally, on its full-restart assumption and illustrative calibration). This is visible forwards and backwards once you analyze his real-world contributions, not merely in terms of their numerosity but in the unique sub-structural subtleties composing each as if a web\footnote{That is, people say that there are plenty of people equally intelligent to Einstein around today, and yes, I concede that that is true if we are talking about sheer cognition (as after all I do not think Einstein could compete with your typical modern abstract algebraist in their domain), but that is not actually what I am talking about when I call Einstein unique. What I am talking about is his ability to convert raw, impenetrable stretches of empirical mystery into the structural content of an ontologically-motivated theory. Given the contrast between what he did in 1905 versus what anyone else on the planet was doing---while I cannot promise he would be deciphering a theory of everything outright, if he were around today he would almost certainly making a substantial dent in the empirical overgrowth somewhere in the landscape of open questions. That is, he knew how to isolate from the crows rich spots in that landscape, and upon finding them he knew how to handle them, so the fact that nowhere do we see the metaphorical lawn being mowed---ore converted to new ingots which are passed around by the day---is decisive evidence to me that no one in the public sphere has this same capacity. If they had the capacity, we would see it, but we see nothing of the kind so the capacity must not be available. (If you want something more concrete, I might point out how the foundational open questions in physics have not changed very much in the past several decades. Say what you will about whether or not this indicates a lack of progress, but I am in the least making the more general argument that it indicates the lack of Einstein cognition---as a kind---among working researchers). \par
This does not rule out people on the fringes who are merely being ignored, however, and one of the great tragedies of history---however speculative---is that perhaps most of the world's Einstein's live and die in obscurity due to the crushing limitations of their settings to fostering their education and creativity. The asymmetry is surely severe---systemic exclusion has operated far more totally on women than on men---though I decline to decorate the point with invented percentages; the mechanism, not a number, is what matters. In all cases I imagine it is something simple: a young girl was told at the age of five by a trusted adult that her eccentric interests subtract from the family, and like a bitter gust of wind more than a conscious-level recognition she is swayed away.}. 
But then, if I were to state this objection, we again run into the problem of \textit{who does this narrative I am proclaiming serve?} I feel the need to correct a miscalculation for the sake of honesty, but at once I find myself reflecting on that honesty and feeling as though I have made the world a slightly worse place by stating it, which then runs the risk of compounding over time. 

I do not at the moment have a solution for this moral ambiguity, but I must stress that it ought to be taken with central seriousness within any future literature on the topic, as cites this paper or not. 

\end{document}